\section{Number-theory consequences with the hypotheses separated}
\label{sec:number-theory}

Throughout this section $A$ denotes the user-designated Alp\"oge theorem
for the Riemann zeta function and Dirichlet L-functions.  It is an adopted
cited input, not a claim that this run rebuilt the unseen Lean dependencies.
$M$ denotes the conditional all-number-field unitary Hecke-GRH assertion
of Masullo, \S1.2.  Every use of $M$ below remains conditional.  The
accompanying catalogue is \texttt{work/number\_theory/rows.json}; its
literature records include retrieval evidence and distinguish a checked
statement from a reread of an entire classical proof.

There are three different geometric translations.  Rational primes are
the closed points of $\operatorname{Spec}\mathbb Z$; prime ideals are the
closed points of $\operatorname{Spec}\mathcal O_K$; and a conjugacy condition
on Frobenius is a condition on a finite \'{e}tale cover of an open arithmetic
curve.  The last dictionary is the Galois-category equivalence in SGA~1,
Expos\'e~V, \S7.  The class group is
$\operatorname{Pic}(\operatorname{Spec}\mathcal O_K)$ and its class fields
are finite \'{e}tale covers.  These translations identify the objects to
which the estimates apply.  They do not turn number-field GRH into a
consequence of finite-field purity.  The inherited cote-1 trace programme
is a comparison of explicit formulas; no new archive scan is certified by
this section.  Exact archive work belongs to the geometry section.

\subsection{The classical arithmetic bounds being used}

\paragraph{NT-P01--NT-P04: rational primes and gaps.}
Schoenfeld, \emph{Sharper bounds for the Chebyshev functions
$\theta(x)$ and $\psi(x)$. II} (1976), Theorem~10, printed p.~337,
equations (6.2) and (6.3), gives, under RH,
\[
 |\psi(x)-x|<\frac{\sqrt{x}\log^2x}{8\pi}\quad(x\ge73.2),\qquad
 |\theta(x)-x|<\frac{\sqrt{x}\log^2x}{8\pi}\quad(x\ge599).
\]
Equation (6.1) replaces $\log^2x$ by $(\log x-2)\log x$ for
$x\ge23\cdot10^8$.  Corollary~1 gives
$|\pi(x)-\operatorname{li}(x)|<\sqrt{x}\log x/(8\pi)$ for
$x\ge2657$.  Here $\operatorname{li}$ has its usual normalization.
The constants in Theorem~10 were checked in the primary indexed text.
An independently retrieved primary proof locator for the same
prime-counting bound is Johnston (2022), Theorem~3.2, equation (3.6),
\emph{The Ramanujan Journal} \textbf{59}, 1307--1321.
It requires RH only to a height $T$ satisfying
$9.06\sqrt{x/\log x}/\log\log x\le T$; under $A$ one may choose
such a $T$ for each $x$.  Its \S2 explicitly says its main theorem applies
also to the ordinary, non-half-weight counting functions.  Equation (3.3)
there also improves the threshold for the displayed $\psi$ inequality
to $x\ge59$.  Thus the inaccessible original Corollary~1 page is no
longer a proof-source gate for NT-P03.

Dudek--Greni\'e--Molteni, \emph{Primes in explicit short intervals on RH}
(2016), Theorem~1.1, author PDF p.~1, gives a prime in
\[
 (x-c\sqrt{x}\log x,\ x+c\sqrt{x}\log x),
 \qquad c=\tfrac12+2/\log x,\qquad x\ge2.
\]
The same theorem gives at least $\sqrt{x}$ primes after replacing $c$ by
$c+1$.  These results imply the classical
$p_{n+1}-p_n=O(\sqrt{p_n}\log p_n)$ consequence of RH.  Neither this
bound nor the assertion of RH supplies a proof here of Cram\'er's
$O(\log^2p_n)$ conjecture, twin primes, or the binary Goldbach conjecture.
Those problems require additional information about prime correlations.

\paragraph{NT-P05--NT-P08: progressions and nonresidues.}
Greni\'e--Molteni--Perelli (2016), Theorem~1, author PDF p.~2, proves that
absolute $x_0,c_1,c_2>0$ work in
\[
 \pi(x+h;q,a)-\pi(x;q,a)\ge
 c_2\frac{h}{\varphi(q)\log x},\qquad
 c_1\varphi(q)\sqrt{x}\log x\le h\le x,\quad x\ge x_0,
\]
assuming the Dirichlet GRH modulo $q$ and $(a,q)=1$.  The following
sentence in that source gives the asymptotic if the ratio to the lower
scale tends to infinity.  We have not converted the unspecified absolute
constants into numerical constants.

Lamzouri--Li--Soundararajan (LLS), corrected
\texttt{arXiv:1309.3595v3}, is used for the sharp explicit conductor
statements.  Corollary~1.2, author p.~3, gives the least prime in a reduced
class modulo $q>3$ at most $(\varphi(q)\log q)^2$.
Corollary~1.1, p.~2, gives a quadratic nonresidue below $(\log q)^2$ for
prime $q\ge5$.  Theorem~1.1(1), p.~2, gives the proper-subgroup bound
recorded with its explicit functions $A(q),B(q)$ in NT-P08.
These are the already carried IA02 consequences, with the corrected
version and hypotheses rechecked; they are not new findings of this run.
Geometrically the progression is a Frobenius class in a cyclotomic cover,
and the nonresidue is detected by a finite rank-one character.

\paragraph{NT-P09: the least primitive root is a different quantifier.}
Shoup, \emph{Searching for primitive roots in finite fields} (1992),
Theorem~3, author PDF p.~4, gives under Dirichlet GRH
\[
 g(p)\ll r^4(\log r+1)^4(\log p)^2,
 \qquad r=\omega(p-1).
\]
The printed theorem was checked on the page image.  Here is the short
simplification to $g(p)\ll\log^6p$.  If $q_1<\cdots<q_r$ are the distinct
prime divisors of $p-1$, then $q_j\ge j+1$ and therefore
$p-1\ge\prod_{j=1}^r(j+1)$.  For $r\ge4$ the last at least $r/2$ factors
are at least $r/2$, so
$\log(p-1)\ge(r/2)\log(r/2)$.  Consequently
$r(\log r+1)\ll\log p$; enlarging the absolute constant handles $r<4$.
Substitution proves the assertion.  This finds a small generator for a
fixed $\mathbb F_p^\times$.  The density theorem for one fixed integer
as $p$ varies uses different, generally nonabelian, fields (NT-A01).

\subsection{A quantitative Frobenius trace transfer}

\paragraph{NT-C01 and NT-C02: the precise cited Chebotarev input.}
Let $L/K$ be a finite Galois extension, $G=\operatorname{Gal}(L/K)$,
$n_L=[L:\mathbb Q]$, and $\Delta_L=|\operatorname{disc}L|$.
For a conjugacy class $C$ put $\delta_C=|C|/|G|$.  For unramified
primes only, define
\[
 \psi_C(x)=\sum_{\substack{\mathfrak p\text{ unramified},\ m\ge1\\
                    (N\mathfrak p)^m\le x}}
 \mathbf1_C(\operatorname{Frob}_{\mathfrak p}^{m})\log N\mathfrak p,
 \qquad
 \pi_C(x)=\sum_{\substack{\mathfrak p\text{ unramified}\\N\mathfrak p\le x}}
             \mathbf1_C(\operatorname{Frob}_{\mathfrak p}).
\]
Set
\begin{align*}
 E_L(x)&=\sqrt{x}\left[
       \left(\frac{\log x}{2\pi}+2\right)\log\Delta_L+
       \left(\frac{\log^2x}{8\pi}+2\right)n_L\right],\qquad x\ge1,\\
 B_L(x)&=\sqrt{x}\left[
       \left(\frac1{2\pi}+\frac3{\log x}\right)\log\Delta_L+
       \left(\frac{\log x}{8\pi}+\frac1{4\pi}+\frac6{\log x}\right)n_L
                         \right],\qquad x\ge2,\\
 \operatorname{Li}_2(x)&=\int_2^x\frac{dt}{\log t}.
\end{align*}
Greni\'e--Molteni, \emph{An explicit Chebotarev density theorem under GRH}
(2019), Theorem~1.1 and Corollary~1.2, author PDF pp.~2--3, give
\[
 |\psi_C(x)-\delta_Cx|\le\delta_C E_L(x),\qquad
 |\pi_C(x)-\delta_C\operatorname{Li}_2(x)|\le\delta_C B_L(x),
\tag{NT-C}
\]
under RH for $\zeta_L$.  The theorem explicitly includes the unramified
version of $\psi_C$.  Its constants and ranges were checked on the page
images.  In particular Artin holomorphy is not an additional hypothesis.
$M$ supplies this input for every $L$; $A$ supplies it if $L/\mathbb Q$
is abelian by the factorization into primitive Dirichlet functions.

\paragraph{NT-F01. Statement and full proof.}
For a complex class function $f$ on $G$ put
\[
 \mu_f=\frac1{|G|}\sum_{g\in G}f(g),\qquad
 \|f\|_{1,G}=\frac1{|G|}\sum_{g\in G}|f(g)|,
\]
and define $\Psi_f$ and $\Pi_f$ by replacing $\mathbf1_C$ in the two
sums above by $f$.  Then
\[
 |\Psi_f(x)-\mu_fx|\le\|f\|_{1,G}E_L(x),\qquad
 |\Pi_f(x)-\mu_f\operatorname{Li}_2(x)|\le\|f\|_{1,G}B_L(x).
\tag{NT-F01}
\]
If $f=\chi_V$ is the character of a finite-dimensional complex
representation, then $\mu_f=\dim V^G$ and
$\|f\|_{1,G}\le\sqrt{\langle\chi_V,\chi_V\rangle_G}$.

To prove the estimates, choose one representative $g_C$ from each
conjugacy class.  The finite identities
\[
 f=\sum_C f(g_C)\mathbf1_C,
 \qquad \mu_f=\sum_C\delta_Cf(g_C)
\]
give $\Psi_f-\mu_fx=\sum_Cf(g_C)(\psi_C-\delta_Cx)$.
The triangle inequality and (NT-C) bound its modulus by
$E_L(x)\sum_C\delta_C|f(g_C)|=E_L(x)\|f\|_{1,G}$.
The same finite argument proves the assertion for $\Pi_f$.
The averaging operator $|G|^{-1}\sum_g\rho(g)$ is the projection to $V^G$,
so its trace is $\mu_{\chi_V}=\dim V^G$.  Cauchy--Schwarz for the
probability measure assigning mass $1/|G|$ to each group element yields
\[
 \left(\frac1{|G|}\sum_g|\chi_V(g)|\right)^2
 \le\frac1{|G|}\sum_g|\chi_V(g)|^2
 =\langle\chi_V,\chi_V\rangle_G.
\]
This proves every assertion.  This is a linear functional estimate on the
complexified representation ring; virtual differences need not be
nonnegative, so their error is controlled by the displayed norm.

\paragraph{NT-F02. Finite \texorpdfstring{\'{e}}{e}tale point counts, with the boundary.}
Let $U$ be a nonempty open subscheme of
$\operatorname{Spec}\mathcal O_K$, and let $X\to U$ be finite \'{e}tale
of degree $d$.  Let $S_X$ be the set of geometric points of its generic
fibre, let $L/K$ be its finite Galois splitting field, and let $r$ be the
number of $G$-orbits on $S_X$, equivalently the number of generic connected
components.  Let $S$ be the finite set of primes not in $U$ which are
unramified in $L$.  Under RH for $\zeta_L$,
\[
 \left|\sum_{\substack{\mathfrak p\in U,\ m\ge1\\(N\mathfrak p)^m\le x}}
       \#X(k_{\mathfrak p^m})\log N\mathfrak p-rx\right|
 \le rE_L(x)+d|S|\log x,\qquad x\ge1.
\tag{NT-F02}
\]
Here $k_{\mathfrak p^m}$ denotes the degree-$m$ extension of
$k_{\mathfrak p}$.  In particular, when no additional unramified prime is
removed, the error is $rE_L(x)$, independent of the degree $d$.

For the proof, finite \'{e}taleness implies that inertia at a prime of $U$
acts trivially on the generic geometric fibre.  The action of
$\operatorname{Gal}(L/K)$ on this fibre is faithful by the definition of
the splitting field, so $L/K$ is unramified at every prime of $U$.
Let $\chi(g)=\#S_X^g$ be the permutation character.  The elementary
Frobenius fixed-point formula for a finite \'{e}tale scheme over a finite
field is
$\#X(k_{\mathfrak p^m})=\chi(\operatorname{Frob}_{\mathfrak p}^{m})$:
its geometric points rational over that field are precisely the points
fixed by the $m$th Frobenius power.  Also $0\le\chi(g)\le d$.
Double counting pairs $(g,s)$ with $gs=s$ gives
\[
 \sum_{g\in G}\chi(g)=\sum_{s\in S_X}|G_s|=r|G|.
\]
Thus both $\mu_\chi$ and $\|\chi\|_{1,G}$ equal $r$.  NT-F01 bounds
the sum over every unramified prime by $rx$ with error $rE_L(x)$.
For each omitted prime, the nonnegative lost contribution is at most
\[
 d\left\lfloor\frac{\log x}{\log N\mathfrak p}\right\rfloor
          \log N\mathfrak p\le d\log x.
\]
If $N\mathfrak p>x$ the contribution is zero.  Summing this bound over
$S$ and applying the triangle inequality proves (NT-F02), including
$x=1$.  Notice that replacing the nonnegative character norm by the square
root of its second moment would give a weaker bound; the orbit count is
the relevant parameter here.

\paragraph{NT-F03. Splitting types.}
For a partition $\lambda$ of $d=|S_X|$, let $D_\lambda\subseteq G$ be
the elements with cycle partition $\lambda$ on $S_X$, and put
$\delta_\lambda=|D_\lambda|/|G|$.  This is a union of conjugacy classes.
Sum the second inequality (NT-C) over these classes and use the triangle
inequality.  Their densities sum to $\delta_\lambda$, proving
\[
 \left|\#\{\mathfrak p:\ N\mathfrak p\le x,\ \mathfrak p
 \text{ unramified},\ \operatorname{Frob}_{\mathfrak p}\in D_\lambda\}
 -\delta_\lambda\operatorname{Li}_2(x)\right|
 \le\delta_\lambda B_L(x).
\]
If $D_\lambda$ is empty both sides before the absolute value are zero.
For the canonical finite \'{e}tale model of the generic algebra over the
unramified locus, each Frobenius orbit is one closed point of its fibre,
and its orbit length is that point's residue degree.  Hence this cycle
partition is exactly its splitting type.  For the originally specified
model $X$ over $U$, restrict the counted primes to $U$; at most $|S|$
terms have been removed, so the same estimate holds with the additional
error $|S|$.  This proves both interpretations without assigning a fibre
of $X$ outside its specified base.

\paragraph{NT-F04. Avoiding a prescribed finite set.}
Let $D\subseteq G$ be a nonempty conjugacy-stable set of density
$\delta_D$, and let $S$ be any finite set of primes of $K$.  The same
finite-union argument gives
$\pi_D(x)\ge\delta_D(\operatorname{Li}_2(x)-B_L(x))$.
At most $|S|$ of the counted primes belong to $S$.  If
\[
 \delta_D\bigl(\operatorname{Li}_2(x)-B_L(x)\bigr)>|S|,
\]
the count outside $S$ is strictly positive.  It is an integer, so an
eligible prime exists.  This proves the avoidance assertion without adding
an unsupported arbitrary-exclusion clause to Bach--Sorenson's least-prime
theorem.  That theorem is NT-C03, Theorem~5.1, printed p.~1729, with
RH hypotheses specified in its \S2, p.~1720.

\subsection{Class numbers, Kummer towers, and Euclidean algorithms}

\paragraph{NT-H01--NT-H04: the regulator normalization.}
Write
\begin{align*}
 U(q)&=2e^\gamma\left(\log\log q-\log2+\tfrac12+
                                      \frac1{\log\log q}\right),\\
 V(q)&=\frac{12e^\gamma}{\pi^2}\left(\log\log q-\log2+\tfrac12+
                   \frac1{\log\log q}+\frac{14\log\log q}{\log q}\right).
\end{align*}
LLS Theorem~1.5, author pp.~3--4, gives
$V(q)^{-1}\le|L(1,\chi)|\le U(q)$ for primitive conductor
$q\ge10^{10}$.  Its Corollary~1.3 is the already carried imaginary
quadratic bound in NT-H02.  The inequality is not a two-sided estimate of
order $\log\log q$: the reciprocal lower bound must be retained.

Here is the full deduction of NT-H03, including its normalization.
For $D>0$ fundamental put $K=\mathbb Q(\sqrt D)$ and let $R_D$ be the
regulator of the full unit group modulo its torsion, so
$R_D=\log\varepsilon$ for a fundamental unit $\varepsilon>1$, regardless
of its norm's sign.  The analytic class-number formula is
\[
 \mathop{\rm Res}_{s=1}\zeta_K(s)=\frac{2h_DR_D}{\sqrt D}.
\]
An exact independent locator is Milne, \emph{Class Field Theory}, v4.03,
Chapter~VI, Corollary~2.12 and Example~2.13, printed p.~191.
Since $\zeta_K(s)=\zeta(s)L(s,\chi_D)$ and $\zeta$ has residue one,
$L(1,\chi_D)=2h_DR_D/\sqrt D>0$.
Multiplying the two inequalities in NT-H01 by $\sqrt D/2$ proves
\[
 \frac{\sqrt D}{2V(D)}\le h_DR_D\le\frac{\sqrt D}{2}U(D),
 \qquad D\ge10^{10}.
\]
This completes the short proof; no bound on $R_D$ was suppressed.

For completeness, the already inherited abelian-field formulation
NT-H04 is the following precise product calculation.  If $K/\mathbb Q$
is abelian with Dirichlet-character group $X_K$, then
\[
 h_KR_K=\frac{w_K\sqrt{\Delta_K}}{2^{r_1}(2\pi)^{r_2}}
              \prod_{\substack{\chi\in X_K\\\chi\ne1}}L(1,\chi).
\tag{NT-H04}
\]
Indeed Artin factorization gives the product for $\zeta_K/\zeta$,
and taking the residue at one gives this formula by the same class-number
formula.  Its left side and the prefactor are positive, so taking moduli
turns the product into $\prod|L(1,\chi)|$.  For any factors whose
conductors are at least $10^{10}$ insert $U$ and $1/V$; retain the finite
remaining product exactly or bound it separately.  This proves every
claimed product bound.  Dividing by an unestimated regulator does not
prove the corresponding bound on $h_K$.  This is the correction already
made in RETURN\_ZETA\_TWO IA02 and IA09, not a new result of this run.

\paragraph{NT-A01 and NT-A02: the nonabelian input in Artin's primitive-root problem.}
For a fixed integer $g\ne-1$ that is not a square, put
$K_n=\mathbb Q(\zeta_n,g^{1/n})$.  Hooley's classical theorem gives a
positive density
\[
 \delta(g)=\sum_{n\ge1}\frac{\mu(n)}{[K_n:\mathbb Q]}
\]
of primes where $g$ is a primitive root, assuming RH for these Kummer
fields.  We cite Hooley, \emph{On Artin's conjecture} (1967), the
unnumbered theorem on pp.~219--220.  The original internal page has not
been freshly obtained here; a fully inspected precise proof locator for
the density implication is Lenstra, \emph{On Artin's conjecture and
Euclid's algorithm in global fields} (1977), Theorem~(3.1), printed
p.~206, specialized to $K=F=\mathbb Q$, $W=\langle g\rangle$, $k=1$,
with density (2.15), p.~206, and positivity criterion (4.1), p.~207.
The series uses actual field degrees; replacing it by the universal
Artin constant for every $g$ loses the entanglement correction.

Lenstra's version also imposes a conjugacy-stable Frobenius condition
$C\subseteq\operatorname{Gal}(F/K)$ and asks that the residue image of
a finitely generated positive-rank $W\subseteq K^\times$ have index
dividing $k$.  Its auxiliary fields are formed from
$K(\zeta_{q(\ell)},W^{1/q(\ell)})$ with
$q(\ell)=\ell^{v_\ell(k)+1}$, taking the compositum for primes dividing
a squarefree $n$.  Theorem~(3.1) assumes RH for these fields' zeta
functions.  Theorem~(4.1) gives the finite obstruction criterion for
positive density.  $M$ supplies the hypotheses; $A$ alone does not
supply them because the Kummer fields need not be abelian over
$\mathbb Q$.  In geometric terms this is a tower of Kummer torsors of
$\mathbb G_m$, not merely the cyclotomic torsion cover.

\paragraph{NT-A03: a PID becomes Euclidean under the stated Kummer GRH.}
Lenstra's Theorem~(9.1), printed p.~218, states the result recorded in
NT-A03 for the ring of $S$-integers with $|S|\ge2$ and trivial ideal
class group.  Equation (9.2) gives its least Euclidean function.
The original page image confirms the weak inequality $|S|\ge2$;
the OCR's apparent strict inequality was not adopted.  The result is a
classical cited theorem with an exact proof locator, not a new proof of
the Euclidean algorithm.  The independent PID hypothesis and the
distinction from a norm-Euclidean algorithm remain in the catalogue.

\subsection{Growth, cancellation, and zero statistics}

\paragraph{NT-L02 and NT-L03. The fixed-character Lindel\"of deduction.}
The classical input is Simoni\v c, \emph{Estimates for L-functions in the
critical strip under GRH with effective applications} (2023), Theorem~1,
equations (4)--(6), proved in its \S3.  For a function $F$ in the Selberg
class with polynomial Euler product, zero-free to the right of the
critical line, it implies, for fixed $\delta>0$ and large $|t|$,
\[
 \log|F(\sigma+it)|,\qquad -\log|F(\sigma+it)|
 \ll_{F,\delta}(\log|t|)^{1-2\delta}\log\log|t|,
 \quad \tfrac12+\delta\le\sigma\le1.
\tag{NT-L-log}
\]
For $\delta\ge1/2$ one uses instead a smaller positive $\delta$ and
the absolutely convergent Euler product on $\sigma>1$.
The modulus bound for $\log F$ in the cited theorem gives both signs of
its real part.  Dividing the right side of (NT-L-log) by $\log|t|$
shows that it tends to zero for fixed $0<\delta<1/2$.  Hence, for every
$\epsilon>0$,
\[
 F(\sigma+it),\quad F(\sigma+it)^{-1}
 \ll_{F,\delta,\epsilon}(1+|t|)^\epsilon
 \qquad(\sigma\ge\tfrac12+\delta),
\tag{NT-L-right}
\]
with pole neighbourhoods omitted if $F$ has a pole.  The finite range of
$t$ is absorbed into the constant on any closed set avoiding poles and
zeros.  On $\sigma\ge1+\delta$ the Euler product gives the required
uniform bound directly; compact intermediate strips are already covered
by (NT-L-log) or the cited theorem's adjacent range.

Here is the passage to the critical line, often suppressed in invoking
Lindel\"of.  Fix $0<\delta<1/2$.  If $F$ has degree $d$, its functional
equation and Stirling's formula imply
\[
 |F(\tfrac12-\delta+it)|
 \ll_{F,\delta}(1+|t|)^{d\delta}
       |F^\vee(\tfrac12+\delta-it)|.
\]
The dual has the same zero-free hypothesis, so NT-L-right bounds this by
$O((1+|t|)^{d\delta+\eta})$ for every $\eta>0$.  The right boundary
$\sigma=1/2+\delta$ has exponent $\eta$.
The Phragm\'en--Lindel\"of principle on this strip gives at its midpoint
\[
 |F(\tfrac12+it)|\ll_{F,\delta,\eta}
                       (1+|t|)^{d\delta/2+\eta}.
\]
Its hypotheses hold because the standard functional equation and
finite-order analytic continuation give polynomial growth in each fixed
vertical strip; if there is a pole, apply the principle above and below
a fixed height containing it.  Choose $\delta$ and $\eta$ so that
$d\delta/2+\eta<\epsilon$.  The remaining bounded interval is compact,
and the critical line contains no pole for these L-functions.  This
proves $F(1/2+it)\ll_{F,\epsilon}(1+|t|)^\epsilon$.

For primitive Dirichlet functions the hypotheses are classical, and $A$
is the zero-free input; this proves NT-L02.  For a fixed non-pure-norm
unitary Hecke character, the Hecke L-function is entire of finite order,
its gamma shifts have nonnegative real parts, and its degree is
$[K:\mathbb Q]$.  On grouping its ideal Euler factors by rational primes,
one obtains a polynomial Euler product of that degree with local roots
of modulus at most one.  The Dirichlet coefficients are bounded by
$d_{[K:\mathbb Q]}(n)=O_\epsilon(n^\epsilon)$: each local coefficient
is bounded by the corresponding coefficient of $(1-p^{-s})^{-[K:\mathbb Q]}$.
Thus it meets the explicit axioms in Simoni\v c \S2.1, and $M$ supplies
zero-freeness.  If the character is a pure unitary norm power, its
L-function is a vertical translate of $\zeta_K$, to which the same
argument applies before translation.  Imprimitive Euler factors, when
included in the convention, form a fixed finite product; they and their
inverses are bounded on each $\sigma\ge1/2+\delta$, since their roots
have modulus at most one.  This proves NT-L03 in all the conventions just
specified.  Constants may depend on the field, character, infinity type,
and $\epsilon$; this proof makes no uniform conductor-aspect claim.

\paragraph{NT-L05 and NT-L06. Mertens, Liouville, and reciprocal tails.}
Titchmarsh, second edition, Theorem~14.25(C), printed p.~370, is the exact
classical equivalence
$\mathrm{RH}\Longleftrightarrow M(x)=O_\epsilon(x^{1/2+\epsilon})$
for every $\epsilon>0$.  We use that theorem rather than reproduce its
contour proof.  The following short deductions are included in full.
For a prime power $p^a$,
\[
 \sum_{d^2\mid p^a}\mu(p^a/d^2)=
 \begin{cases}1,&a\text{ even},\\-1,&a\text{ odd},\end{cases}
\]
because only the term with remaining exponent zero or one is nonzero.
Both sides are multiplicative, hence
$\lambda(n)=\sum_{d^2\mid n}\mu(n/d^2)$ for every $n$.
Consequently
\[
 \left|\sum_{n\le x}\lambda(n)\right|
 \le\sum_{d\le\sqrt x}|M(x/d^2)|
 \ll_\epsilon x^{1/2+\epsilon}
                \sum_{d\ge1}d^{-1-2\epsilon}
 \ll_\epsilon x^{1/2+\epsilon}.
\]
This proves Liouville cancellation, with the dependence on $\epsilon$
retained in the convergent sum.

Let $s=\sigma+it$ and choose $\eta>0$ with $1/2+\eta<\sigma$.
Partial summation for $y>x$ gives
\[
 \sum_{x<n\le y}\frac{\mu(n)}{n^s}
 =M(y)y^{-s}-M(x)x^{-s}
     +s\int_x^y M(u)u^{-s-1}\,du.
\]
The first term tends to zero as $y\to\infty$, and the integral converges
absolutely after using $|M(u)|\ll_\eta u^{1/2+\eta}$.  Its modulus is at
most a constant times
$|s|x^{1/2+\eta-\sigma}/(\sigma-1/2-\eta)$.
Together with the boundary term this proves the stated tail estimate.
It also proves local uniform convergence of the reciprocal Dirichlet
series on $\operatorname{Re}s>1/2$ by choosing $\eta$ below the distance
of a compact set from the boundary.  It agrees with $1/\zeta(s)$ on
$\operatorname{Re}s>1$ by the absolutely convergent Euler product and
therefore throughout that half-plane by analytic continuation.  At
$s=1$, $1/\zeta$ has value zero.  Hence $\sum\mu(n)/n=0$, as asserted.

\paragraph{NT-L04, NT-L07, NT-L08: what moments and zero density say.}
Harper (2013), Theorem~1, author p.~1, gives
$\int_T^{2T}|\zeta(1/2+it)|^{2k}dt\ll_kT(\log T)^{k^2}$ for fixed
$k\ge0$ under RH.  Its proof does not supply the conjectural leading
constant for every $k$.  Simoni\v c (2023), Theorem~2, equations (7),(8),
provides the explicit effective Mertens constants in NT-L07.
For NT-L08, if $\sigma>1/2$, every nontrivial zero covered by $A$ or
conditionally $M$ has real part $1/2<\sigma$, so the indicated counting
set is empty.  At $\sigma=1/2$ this reasoning says nothing about
multiplicity.  Trivial zeros from gamma factors and zeros of removed
finite Euler factors are not silently included in this definition.

\paragraph{NT-Z01 and NT-Z02: the exact pair-correlation range.}
For zeta zero ordinates $0<\gamma,\gamma'\le T$, counted with
multiplicity, put
\[
 F(\alpha,T)=\left(\frac{T}{2\pi}\log T\right)^{-1}
 \sum_{\gamma,\gamma'}T^{i\alpha(\gamma-\gamma')}
                    \frac4{4+(\gamma-\gamma')^2}.
\]
Montgomery (1973), \S1, unnumbered theorem, printed p.~181, proves under
RH the formula in NT-Z01 uniformly on each $0\le\alpha\le1-\epsilon$.
Corollary~2, p.~182, gives the proportion $2/3+o(1)$ of simple zeros.
The extension of the pair-correlation formula beyond that range is
explicitly a conjecture in that paper.  Adopting RH removes the theorem's
hypothesis; it does not replace its range by all $\alpha$.

\paragraph{NT-Z03--NT-Z05: low zeros of the prime-modulus Dirichlet family.}
For prime $q$ average over the $q-2$ nonprincipal characters, with scaled
ordinates $x_{\chi,j}=\gamma_{\chi,j}\log q/(2\pi)$.
Hughes--Rudnick, \emph{Linear statistics of low-lying zeros of
L-functions} (2003), Theorem~3.1, author p.~5, gives the unitary mean for
admissible even test functions with Fourier support in $[-2,2]$.
Theorem~4.1, p.~7, gives variance
$\int_{-1}^1|u|\widehat f(u)^2du$ with support in $[-1,1]$.
Theorem~5.1, p.~9, gives the Gaussian moments only for
$m<2/\alpha$ when the support is in $[-\alpha,\alpha]$.
These identities are unconditional with complex $\gamma$ allowed;
$A$ supplies the real-ordinate interpretation.  They are not an
unrestricted all-moment Gaussian theorem for one fixed test function.

In that source's ordering on p.~1, $\gamma_{\chi,1}$ is the first
\emph{nonnegative} ordinate, so a central zero is included with value zero.
With exactly this convention, Corollary~8.2, p.~20, gives
$\liminf_q\min_{\chi\ne\chi_0}x_{\chi,1}\le1/4$ under $A$.
Theorem~8.3, pp.~20--21, gives a positive proportion for
$x_{\chi,1}<\beta$ if $\beta\ge0.633$.
No strictly-positive-zero variant was inferred by discarding central
zeros.  The geometric comparison is with finite-field family monodromy
and its symmetry group; it does not identify the unrestricted
number-field limiting distribution.

\paragraph{NT-Z06 and NT-Z07: the additional gates.}
No argument in the cited sources proves full GUE, simplicity of every
zero, or central nonvanishing for every quadratic Dirichlet function
from zero location alone.  We record the specific missing correlation
or nonvanishing estimates, without claiming a logical independence
theorem.  For Iwaniec--Luo--Sarnak, \emph{Low lying zeros of families of
L-functions} (2000), the standing assumption on printed p.~56 explicitly
includes RH for the GL$_2$ functions in the family as well as Dirichlet
GRH.  Theorem~1.1, pp.~58--59, gives support $(-2,2)$ for fixed even
weight and squarefree level tending to infinity.  $A$ supplies its
Dirichlet part.  A universal Hecke statement does not identify a general
non-CM GL$_2$ newform with a rank-one Hecke character.  Thus the GL$_2$
hypothesis remains an additional input in this invocation.

\paragraph{NT-Z08 and NT-Z09: two exact classical equivalents.}
Rodgers--Tao, \texttt{arXiv:1801.05914v5}, Theorem~1, author p.~3,
proves $\Lambda\ge0$ for the classical de Bruijn--Newman constant.
The definition and RH equivalence in its \S1 give
$A\Rightarrow\Lambda\le0$, hence $\Lambda=0$.
For the deformation as normalized there, real-rootedness holds exactly
for $t\ge\Lambda$, yielding the two assertions in NT-Z08.
Lagarias, \texttt{arXiv:math/0008177v2}, Theorem~1.1, author p.~1,
is the exact locator for NT-Z09's divisor-sum inequality.
These are cited classical consequences, not additional new proofs of RH.

\subsection{Fixed-character cancellation and the local zero count}

\paragraph{NT-L09 and NT-L10. A full Perron transfer.}
Here is the elementary analytic form used for both rational and ideal
M\"obius sums.  Suppose
$G(s)=\sum_{n\ge1}a_nn^{-s}$ is absolutely convergent for
$\operatorname{Re}s>1$, $|a_n|\le d_D(n)$ for a fixed positive integer
$D$, and $G$ extends holomorphically to $\operatorname{Re}s>1/2$.
Assume that for every $\delta,\eta>0$,
\[
 G(\sigma+it)\ll_{\delta,\eta}(1+|t|)^\eta
 \quad\text{uniformly for }\tfrac12+\delta\le\sigma\le2.
\tag{NT-Perron-input}
\]
Then $\sum_{n\le x}a_n\ll_\epsilon x^{1/2+\epsilon}$ for every
$\epsilon>0$.

We prove the assertion, recording the truncation error as well as the
contour displacement.  It suffices to take $0<\epsilon<1/4$ and
half-integral $x\ge4$, because replacing any real $y$ by
$\lfloor y\rfloor+1/2$ leaves the partial sum unchanged and changes its
size by a bounded factor; the bounded remaining range is absorbed into
the constant.  Put
$c=1+1/\log x$, $T=x^2$, $\delta=\epsilon/4$ and
$\eta=\epsilon/8$.  The classical truncated Perron kernel gives
\[
 \sum_{n\le x}a_n=
 \frac1{2\pi i}\int_{c-iT}^{c+iT}G(s)\frac{x^s}{s}\,ds+R,
\]
with
\[
 |R|\ll\sum_{n\ge1}|a_n|(x/n)^c
           \min\left(1,\frac1{T|\log(x/n)|}\right).
\tag{NT-Perron-error}
\]
This is the usual Perron formula; an exact classical locator is
Titchmarsh, Lemma~3.12 and its proof, printed pp.~61--62.  The kernel
estimate follows by moving the integral of $y^s/s$ to the left or right
according to $y>1$ or $y<1$, taking the residue at zero in the first
case, and integrating the horizontal tails; the alternative trivial
bound handles $|\log y|\le T^{-1}$.

For $n\le x/2$ or $n\ge2x$ the denominator $|\log(x/n)|$ is bounded
below by $\log2$, so that part of (NT-Perron-error) is
\[
 \ll x^cT^{-1}\sum_{n\ge1}d_D(n)n^{-c}
 =x^cT^{-1}\zeta(c)^D\ll xT^{-1}(\log x)^D.
\]
For $x/2<n<2x$, use $d_D(n)\ll_{D,\eta}x^\eta$ and
$|\log(x/n)|\gg |x-n|/x$.  Since $x$ is half-integral,
$\sum_{x/2<n<2x}|x-n|^{-1}\ll\log x$.  Also $(x/n)^c\ll1$,
uniformly for these $x,n$.  This part of the error is therefore
$O(x^{1+\eta}T^{-1}\log x)$.  Both parts are smaller than
$x^{1/2+\epsilon}$.

Move the vertical line to $\sigma_0=1/2+\delta$.  The integrand has no
pole in the intervening rectangle: $G$ is holomorphic there and
$\sigma_0>0$.  On either horizontal segment,
(NT-Perron-input) bounds the integral by
$O(x^cT^{\eta-1})=O(x^{-1+2\eta})$.
On the new vertical segment it gives
\[
 \ll x^{\sigma_0}\left(1+\int_1^Tt^{\eta-1}\,dt\right)
 \ll_\eta x^{1/2+\delta}T^\eta
 =O(x^{1/2+\epsilon/2}).
\]
The bounded part $|t|\le1$ is harmless because $\sigma_0>0$ and $G$
is bounded on its compact segment.  This proves the asserted partial-sum
bound with room in the exponent.

Apply this lemma first to $G(s)=1/L(s,\chi)$.
Its coefficients are $a_n=\mu(n)\chi(n)$, including the usual zero
values at integers not coprime to the conductor.  Under $A$, the
reciprocal is holomorphic on $\operatorname{Re}s>1/2$; the principal
L-function's pole produces a zero of the reciprocal, not a pole.
NT-L02 supplies (NT-Perron-input), so the twisted bound follows.
Character orthogonality, for $(a,q)=1$, is the finite identity
\[
 \sum_{\substack{n\le x\\n\equiv a\ ({\rm mod}\ q)}}\mu(n)
 =\frac1{\varphi(q)}\sum_{\chi\ ({\rm mod}\ q)}
       \overline{\chi(a)}\sum_{n\le x}\mu(n)\chi(n).
\]
Imprimitive functions differ by finitely many Euler factors whose zeros
lie on $\operatorname{Re}s=0$, and their reciprocals are bounded in any
closed half-strip to the right of $1/2$.  Thus each term satisfies the
same fixed-character argument.  A finite sum proves the fixed-$q$
progression bound, with no uniformity in $q$ claimed.

For a unitary Hecke character, expand
\[
 \frac1{L(s,\eta)}
 =\sum_{(\mathfrak a,\mathfrak f)=1}
         \mu_K(\mathfrak a)\eta(\mathfrak a)(N\mathfrak a)^{-s}
 =\sum_{n\ge1}a_nn^{-s}.
\]
For each $n$, $|a_n|$ is at most the number of ideals of norm $n$.
If $[K:\mathbb Q]=D$, comparison of the rational-prime Euler factors
of $\zeta_K$ with $(1-p^{-s})^{-D}$ gives this number at most $d_D(n)$.
Indeed there are at most $D$ prime ideals above $p$, each of residue
degree at least one, and replacing every local factor by a degree-one
factor can only increase its nonnegative coefficients.  Holomorphy and
(NT-Perron-input) follow conditionally on $M$ from NT-L03, including
pure norm twists.  The lemma proves NT-L10.  Averaging over the finite
character group of any fixed ray class group then projects onto one ray
class exactly as above; all sums are over ideals prime to its modulus.
This completes every claimed cancellation statement.

\paragraph{NT-Z10. The fixed infinite-order Hecke zero count.}
Let $\eta$ be any fixed unitary Hecke character of a number field $K$.
Its number of nontrivial zeros, counted with multiplicity, in
\[
 0\le\operatorname{Re}s\le1,\qquad |\operatorname{Im}s-T|\le1
\]
is $O_{K,\eta}(\log(|T|+3))$, unconditionally.  This assertion includes
infinite-order characters; its constant is not asserted uniform in
their infinity types.

Here is a full local argument.  On any fixed vertical strip the finite
Hecke L-function has a polynomial upper bound in $|t|$, away from its
possible pole.  In the present argument this follows directly from the
unconditional first part of Simoni\v c's Theorem~1, equation (2), on
$\sigma\ge1/2$, and from the functional equation and Stirling's formula
on the reflected strip.  The Selberg-class hypotheses for a non-pure-norm
character were verified in NT-L03 without using GRH; a pure norm twist
is a translate of the Dedekind function.  Consequently constants $C,A>0$
exist such that, for large $|T|$,
\[
 \max_{|s-(2+iT)|\le4}|L(s,\eta)|\le C(|T|+3)^A.
\tag{NT-J-upper}
\]
For a pure norm twist the only pole has a fixed imaginary coordinate, so
choose $|T|$ sufficiently large that this disk does not contain it.
In all other cases the function is entire.

The Euler product at the centre gives a lower bound that a general
Schwartz Mellin transform need not have:
\[
 |L(2+iT,\eta)|
 \ge\prod_{\mathfrak p\nmid\mathfrak f}(1+(N\mathfrak p)^{-2})^{-1}
 \ge\prod_{\mathfrak p}(1-(N\mathfrak p)^{-2})
 =\zeta_K(2)^{-1}>0.
\tag{NT-J-lower}
\]
These comparisons follow term by term from $|\eta(\mathfrak p)|=1$
and $1/(1+u)\ge1-u$ for $0\le u<1$, and persist on taking limits of
the absolutely convergent products.

Let $n_T(3)$ count zeros in the disk of radius $3$ about $2+iT$.
Jensen's formula on the radius-$4$ disk, or its limiting-radius version
if a zero lies on the outer circle, implies
\[
 n_T(3)\log(4/3)
 \le\log\max_{|s-(2+iT)|\le4}|L(s,\eta)|
                -\log|L(2+iT,\eta)|.
\]
The right side is $O_{K,\eta}(\log(|T|+3))$ by (NT-J-upper) and
(NT-J-lower).  Every point with $0\le\operatorname{Re}s\le1$ and
$|\operatorname{Im}s-T|\le1$ is at distance at most $\sqrt5<3$
from the centre, so its zeros are included in $n_T(3)$.  The remaining
bounded interval of $T$ meets only finitely many zeros, by meromorphic
analyticity and the absence of an identically zero function; increasing
the constant covers it.  This proves the claim.

The finite-order conductor-explicit counterpart is
Greni\'e--Molteni2019, Appendix~A, Theorem~A.1, author p.~26.
NT-Z10 instead proves the fixed-character form needed for the local
arithmetic divisor, and supplies the infinite-order version of root
MT-M04.  The argument uses the lower bound at its Jensen centre and
polynomial vertical growth.  It does not repair a statement quantified
over every function in an arbitrary Schwartz Mellin-transform space.

\subsection{An explicit uniform cyclotomic specialization}

\paragraph{NT-P11 and NT-P12.}
For $q\ge3$ and $(a,q)=1$, define
\[
 \psi(x;q,a)=\sum_{\substack{p^m\le x\\p^m\equiv a\ ({\rm mod}\ q)}}\log p,
 \qquad
 \pi(x;q,a)=\#\{p\le x:p\equiv a\ ({\rm mod}\ q)\}.
\]
Put $e_q=1$ if $q\equiv2\pmod4$ and $e_q=0$ otherwise.
Then $A$ implies
\begin{align*}
 \left|\psi(x;q,a)-\frac{x}{\varphi(q)}\right|
 &\le\sqrt{x}\left[
 \left(\frac{\log x}{2\pi}+2\right)\log q+
 \frac{\log^2x}{8\pi}+2\right]+e_q\log x,\qquad x\ge1,\\
 \left|\pi(x;q,a)-\frac{\operatorname{Li}_2(x)}{\varphi(q)}\right|
 &\le\sqrt{x}\left[
 \left(\frac1{2\pi}+\frac3{\log x}\right)\log q+
 \frac{\log x}{8\pi}+\frac1{4\pi}+\frac6{\log x}\right]+e_q,
 \qquad x\ge2.
\end{align*}

For the full proof, take $L=\mathbb Q(\zeta_q)$ and $K=\mathbb Q$.
The Galois group is $(\mathbb Z/q\mathbb Z)^\times$, of order
$n_L=\varphi(q)$, and its conjugacy classes are singletons.
The cyclotomic discriminant formula gives
\[
 \Delta_L=q^{\varphi(q)}
    \prod_{p\mid q}p^{-\varphi(q)/(p-1)},\qquad
 \log\Delta_L\le\varphi(q)\log q.
\]
At every prime $p\nmid q$, arithmetic Frobenius acts on $\zeta_q$ by
$\zeta_q\mapsto\zeta_q^p$.  Thus its $m$th power is the element
represented by $p^m$ modulo $q$.  The use of arithmetic or geometric
Frobenius merely replaces the labelled class by its inverse; the estimate
is identical after that relabelling.  The Dirichlet factorization of
$\zeta_L$ makes RH for it a consequence of $A$.
Apply NT-C01 and NT-C02 with $\delta_C=1/\varphi(q)$ and insert the
discriminant upper bound; all displayed coefficients multiplying
$\log\Delta_L$ are positive on the stated ranges.

If $q$ is odd or $4\mid q$, the primes unramified in $L$ are exactly
those not dividing $q$, so the resulting sums are the progression sums
just defined.  If $q\equiv2\pmod4$, write $q=2m$ with $m$ odd.
Then $L=\mathbb Q(\zeta_m)$ and the prime $2$ is unramified, whereas it
is excluded from a reduced class modulo $q$.  For odd primes the class
condition modulo $m$ is equivalent to that modulo $q$, since reduction
$(\mathbb Z/q\mathbb Z)^\times\to(\mathbb Z/m\mathbb Z)^\times$ is
an isomorphism.  The unramified Chebotarev sum can therefore differ only
by the powers of $2$ whose Frobenius powers lie in the designated class.
Their total weight is at most
$\lfloor\log x/\log2\rfloor\log2\le\log x$.
For the prime-counting sum the discrepancy is at most one.  Applying the
triangle inequality proves both formulas, including this degeneration.

The small moduli are transparent separately:
\begin{align*}
 \psi(x;1,0)&=\psi(x),& \pi(x;1,0)&=\pi(x),\\
 \psi(x;2,1)&=\psi(x)-\lfloor\log x/\log2\rfloor\log2,&
 \pi(x;2,1)&=\pi(x)-\mathbf1_{x\ge2}.
\end{align*}
For $1\le x<2$ the floor is zero.  Hence NT-P01 and NT-P03 give the
corresponding bounds in their stated ranges, with the explicit powers-of-two
correction retained for modulus two.  These are estimates on the open
cyclotomic arithmetic curve, and the additional term is exactly its
removed-prime contribution.

\subsection{A specified quadratic family and central nonvanishing}

\paragraph{NT-Z11. The exact one-level-density input.}
Let $\mathcal D$ consist of the nonzero odd squarefree integers, of both
signs.  For $d\in\mathcal D$, $8d$ is a fundamental discriminant and
$\chi_{8d}$ is its primitive quadratic character.  Put
$\ell_X=\log(X/(2\pi e))$.  Let $w$ be an even, nonnegative, nonzero
Schwartz function and write
\[
 W_w(X)=\sum_{d\in\mathcal D}w(d/X),\qquad
 \mathcal Z_{w,\phi}(X)=\frac1{W_w(X)}
  \sum_{d\in\mathcal D}w(d/X)
       \sum_{\rho=1/2+i\gamma}\phi(\gamma\ell_X/(2\pi)),
\]
where the inner sum counts every nontrivial zero with multiplicity.
Use the Fourier convention
$\widehat\phi(u)=\int_{\mathbb R}\phi(t)e^{-2\pi iut}\,dt$.
For a real even Schwartz function $\phi$ with
$\operatorname{supp}\widehat\phi\subset(-2,2)$, $A$ implies
\[
 \lim_{X\to\infty}\mathcal Z_{w,\phi}(X)
 =\widehat\phi(0)-\frac12\int_{-1}^1\widehat\phi(u)\,du.
\tag{NT-Z11}
\]
This is the leading term of Fiorilli--Parks--S\"odergren,
\texttt{arXiv:1601.06833}, Theorem~1.1, author p.~2, with definitions
(1.1)--(1.2) on the same page.  Its proof uses Theorem~3.5, author p.~12, and is concluded on p.~14.
The GRH input in the prime-sum step is Dirichlet GRH
(Lemma~2.5, author p.~7), which $A$ supplies.  No unrestricted Fourier
support, individual central nonvanishing, or nonabelian GRH is being
asserted here.

\paragraph{NT-Z12. A full nonvanishing deduction, including deweighting.}
Under $A$, at least $15/16$ of this specified family have a nonzero
central value in the following precise sense:
\[
 \liminf_{X\to\infty}
 \frac{\#\{d\in\mathcal D:|d|\le X,\ L(1/2,\chi_{8d})\ne0\}}
      {\#\{d\in\mathcal D:|d|\le X\}}\ge\frac{15}{16}.
\tag{NT-Z12}
\]
The same lower bound holds for the weighted proportions defined by any
of the fixed weights $w$ in NT-Z11.  This is a deduction from the cited
classical density theorem, not a claim of a new nonvanishing theorem.
Here are all approximation and counting details.

Let $m_d$ be the order of vanishing at $1/2$.  The completed primitive
quadratic function has root number $+1$.  One way to see the sign without
choosing a Gauss-sum square root is to divide the symmetric completed
Dedekind functional equation of $\mathbb Q(\sqrt{8d})$ by that of
$\mathbb Q$: the factorization
$\zeta_{\mathbb Q(\sqrt{8d})}=\zeta L(s,\chi_{8d})$ and the real or
complex gamma factors identify the quotient, up to a nonzero constant,
with the completed quadratic function.  The gamma duplication formula
handles the imaginary quadratic case.  Thus this entire function is
even about $s=1/2$.  Its gamma factor is finite and nonzero at the
centre, so $m_d$ is an even nonnegative integer.  In particular
\[
 2\mathbf1_{L(1/2,\chi_{8d})=0}\le m_d.
\tag{NT-central-parity}
\]
These functional equations are the standard Hecke equations; a precise
source for the completed convention is Ramakrishnan--Valenza,
Theorem~7-19, printed p.~279.

Choose real even smooth functions $g_j$ compactly supported in $(-1,1)$
such that $\int g_j=1$ and
$g_j\to g=\tfrac12\mathbf1_{[-1,1]}$ in $L^2(\mathbb R)$.
For example, smooth even cutoffs of $[-1+1/j,1-1/j]$, normalized by
their integrals, have these properties.  Define
\[
 f_j(t)=\int g_j(u)e^{2\pi iut}\,du,\qquad
 \phi_j(t)=|f_j(t)|^2.
\]
Then $\phi_j$ is real, even, nonnegative and Schwartz,
$\phi_j(0)=1$, and
$\operatorname{supp}\widehat\phi_j\subset(-2,2)$.
The support assertion follows because
$\widehat\phi_j=g_j*\widetilde g_j$, where
$\widetilde g_j(u)=\overline{g_j(-u)}$.
Plancherel gives $f_j\to f$ in $L^2$, with
\[
 f(t)=\frac{\sin(2\pi t)}{2\pi t}.
\]
Moreover, Cauchy--Schwarz yields
\[
 \|\phi_j-|f|^2\|_1
 \le\|f_j-f\|_2(\|f_j\|_2+\|f\|_2)\longrightarrow0.
\]
Consequently their Fourier transforms converge uniformly.  Since
\[
 \widehat{|f|^2}(u)=
 \begin{cases}\tfrac12-\tfrac14|u|,&|u|\le2,\\0,&|u|>2,
 \end{cases}
\]
we obtain
\[
 \lim_{j\to\infty}
 \left(\widehat\phi_j(0)-\frac12
                   \int_{-1}^1\widehat\phi_j(u)\,du\right)
 =\frac12-\frac12\cdot\frac34=\frac18.
\tag{NT-test-limit}
\]
Under $A$ every ordinate in NT-Z11 is real.  The nonnegativity of
$\phi_j$ and $\phi_j(0)=1$ therefore imply
\[
 \frac1{W_w(X)}\sum_{d\in\mathcal D}w(d/X)m_d
 \le\mathcal Z_{w,\phi_j}(X).
\]
First take $\limsup_{X\to\infty}$ for a fixed $j$, use NT-Z11, and
then let $j\to\infty$.  The result and (NT-central-parity) give
\[
 \limsup_{X\to\infty}
 \frac{\sum_{d\in\mathcal D}w(d/X)
               \mathbf1_{L(1/2,\chi_{8d})=0}}{W_w(X)}\le\frac1{16}.
\tag{NT-weighted-vanishing}
\]
This proves the weighted assertion without applying the density theorem
directly to a nonsmooth Fourier transform or to support touching $2$.

For completeness, the ordinary count follows by a separate elementary
argument.  If $A_+(Y)$ counts positive odd squarefree integers up to $Y$,
the squarefree indicator gives
\begin{align*}
 A_+(Y)&=\sum_{\substack{k\le\sqrt Y\\k\ {\rm odd}}}\mu(k)
       \#\{n\le Y/k^2:n\ {\rm odd}\}\\
 &=\frac Y2\sum_{k\ {\rm odd}}\frac{\mu(k)}{k^2}+O(\sqrt Y)
 =\frac4{\pi^2}Y+O(\sqrt Y).
\end{align*}
Here the omitted tail is $O(Y\sum_{k>\sqrt Y}k^{-2})=O(\sqrt Y)$,
and the Euler product is
$\sum_{k\ {\rm odd}}\mu(k)/k^2=8/\pi^2$.
Thus $N(X):=\#\{d\in\mathcal D:|d|\le X\}
 =8X/\pi^2+O(\sqrt X)$.
Partial summation of this counting formula gives, for every fixed
smooth compactly supported $w$,
\[
 W_w(X)=\frac4{\pi^2}X\int_{\mathbb R}w(t)\,dt+O_w(\sqrt X).
\]
Indeed apply Stieltjes integration to the positive count $A_+$ and
its reflected negative count; the integrated error is bounded by a
constant depending on $\int_0^\infty\sqrt t\,|w'(t)|dt$.

Fix $\delta>0$.  Take such an even smooth $w$ equal to $1$ on
$[-1,1]$, between $0$ and $1$ everywhere, and supported in
$[-1-\delta,1+\delta]$.  The number of vanishing members with
$|d|\le X$ is bounded above by its weighted counterpart.  Dividing by
$N(X)$ and using (NT-weighted-vanishing) gives a $\limsup$ of at most
$\frac1{16}\cdot\frac12\int w\le(1+\delta)/16$.
Letting $\delta\downarrow0$ proves (NT-Z12).

The characters in this statement are quadratic rank-one Artin objects.
Their varying conductors describe a family of double covers of open
arithmetic curves.  The symplectic expression in NT-Z11 matches the
finite-field hyperelliptic-family comparison discussed immediately after
Theorem~1.1 of the cited paper.  This comparison does not enlarge the
analytic support range or turn the $15/16$ bound into nonvanishing for
each member.  The ordinary proportion proved here concerns the
$\chi_{8d}$ family exactly; it has not been silently relabelled as all
fundamental discriminants.

\subsection{Growing fields, class numbers, and restricted unramified towers}

\paragraph{NT-H05 and NT-H06. Exact growing-field inputs.}
For pairwise nonisomorphic number fields $K_i\ne\mathbb Q$, write
$g_i=\tfrac12\log\Delta_{K_i}$ and $n_i=[K_i:\mathbb Q]$.
Let $N_q(K_i)$ count prime ideals of norm $q$, where $q$ ranges over
rational prime powers.  The family is \emph{asymptotically exact} when
the following limits exist for every indicated index:
\[
 \phi_q=\lim_i\frac{N_q(K_i)}{g_i},\qquad
 \phi_{\mathbb R}=\lim_i\frac{r_1(K_i)}{g_i},\qquad
 \phi_{\mathbb C}=\lim_i\frac{r_2(K_i)}{g_i}.
\]
Let $h_i,R_i$ be the ordinary class number and full unit regulator, and
let $\kappa_i=\operatorname{Res}_{s=1}\zeta_{K_i}(s)$.
Tsfasman--Vl\u adu\c t, \texttt{arXiv:math/0205129},
GRH Theorem~7.2, author p.~44 (proof pp.~44--49), gives
\begin{align*}
 \lim_i\frac{\log(h_iR_i)}{g_i}
 &=1+\sum_q\phi_q\log\frac q{q-1}
      -\phi_{\mathbb R}\log2-\phi_{\mathbb C}\log(2\pi),\\
 \lim_i\frac{\log\kappa_i}{g_i}
 &=\sum_q\phi_q\log\frac q{q-1}.
\end{align*}
Only RH for the family of Dedekind functions is needed: $M$ supplies it;
$A$ suffices when every $K_i/\mathbb Q$ is abelian.  Neither normality
nor a tower is required in this theorem.  If $n_i/g_i\to0$, then
$N_q(K_i)\le n_i$ makes all the displayed invariants zero, so the
first limit is $1$.  This proves the usual Brauer--Siegel asymptotic
in this range without adding a normality hypothesis.

The same paper's GRH Theorem~3.1, author p.~19 (proof pp.~19--21),
supplies the basic inequality
\[
 \sum_q\frac{\phi_q\log q}{\sqrt q-1}
 +\alpha_{\mathbb R}\phi_{\mathbb R}
 +\alpha_{\mathbb C}\phi_{\mathbb C}\le1,
\tag{NT-basic}
\]
where
\[
 \alpha_{\mathbb R}=\tfrac12\log(8\pi)+\tfrac\pi4+\tfrac\gamma2,
 \qquad
 \alpha_{\mathbb C}=\log(8\pi)+\gamma,
\]
and $\gamma$ is Euler's constant.  Both original theorem pages were
inspected as images.  These are cited classical theorems.  The next
two consequences and their exact tower interpretations are proved below.

\paragraph{NT-H07. A root-discriminant bound with prescribed splitting.}
Suppose $n_i\to\infty$,
\[
 \Delta_{K_i}^{1/n_i}\longrightarrow\varrho\in(1,\infty),
 \qquad r_1(K_i)/n_i\longrightarrow\tau\in[0,1].
\]
Let $P$ be any set of rational primes such that every $p\in P$ splits
completely in every $K_i$.  Under RH for all $\zeta_{K_i}$,
\[
 \log\varrho\ge\log(8\pi)+\gamma+\frac\pi2\tau
          +2\sum_{p\in P}\frac{\log p}{\sqrt p-1}.
\tag{NT-H07}
\]
If $P$ is infinite, the inequality includes convergence of this
nonnegative series.  No normality assumption on $K_i/\mathbb Q$ is
needed.

To prove the claim, first take a subsequence with strictly increasing
degrees; its fields are pairwise nonisomorphic and its stated limits
are unchanged.  Note that
$g_i/n_i\to\tfrac12\log\varrho>0$.
For every fixed prime power $q$, the numbers $N_q(K_i)/g_i$ are
bounded since $N_q(K_i)\le n_i$.  Successive subsequence extraction,
followed by a diagonal subsequence over the countable prime powers,
makes all of them converge.  The archimedean ratios already converge;
on the resulting asymptotically exact subfamily they are
\[
 \phi_{\mathbb R}=\frac{2\tau}{\log\varrho},\qquad
 \phi_{\mathbb C}=\frac{1-\tau}{\log\varrho}.
\]
Complete splitting gives $N_p(K_i)=n_i$, hence
$\phi_p=2/\log\varrho$ for $p\in P$.
Apply (NT-basic) and discard the nonnegative terms from other norms.
For every finite $P_0\subset P$, multiplication by $\log\varrho$
yields
\[
 2\sum_{p\in P_0}\frac{\log p}{\sqrt p-1}
 +2\tau\alpha_{\mathbb R}+(1-\tau)\alpha_{\mathbb C}
 \le\log\varrho.
\]
The identity
$2\alpha_{\mathbb R}-\alpha_{\mathbb C}=\pi/2$
gives the claimed expression for finite $P_0$.  Taking the supremum
over finite subsets proves convergence and the same bound for all of
$P$.  The limiting quantities $\varrho,\tau$ are unchanged by the
subsequence, so this proves the original assertion.

\paragraph{NT-H08. Finiteness of a restricted arithmetic fundamental group.}
Fix a number field $K$ and a finite set $P$ of rational primes that
split completely in $K$.  Let $K^{\rm ur,P}$ be the compositum of all
finite extensions of $K$ unramified at every finite prime, in which
every prime of $K$ above $P$ splits completely.  Put
\[
 C(P)=8\pi\exp\left(\gamma+
                 2\sum_{p\in P}\frac{\log p}{\sqrt p-1}\right).
\]
If $\Delta_K^{1/[K:\mathbb Q]}<C(P)$, then $M$ implies
$[K^{\rm ur,P}:K]<\infty$.
If $K$ is totally real and only totally real extensions are allowed,
the corresponding conclusion holds under the weaker inequality
$\Delta_K^{1/[K:\mathbb Q]}<e^{\pi/2}C(P)$.
The inequalities are strict: no equality case is claimed.

Here is the full proof.  Being unramified at finite primes and splitting
completely at specified primes are both preserved by conjugation over
$K$ and by finite composita.  Thus the displayed compositum is Galois
over $K$, and all of its finite subextensions have these properties.
If it had infinite degree, it would contain a sequence of finite
Galois subextensions $K_i/K$ with degrees tending to infinity.
The relative discriminant formula and finite unramifiedness give
\[
 \Delta_{K_i}=\Delta_K^{[K_i:K]},\qquad
 \Delta_{K_i}^{1/[K_i:\mathbb Q]}=
                  \Delta_K^{1/[K:\mathbb Q]}=: \varrho.
\]
If $K=\mathbb Q$, this equality gives discriminant one for any such
$K_i$; Minkowski's discriminant bound already rules out a nontrivial
extension.  Otherwise $\varrho>1$.
After taking a subsequence, $r_1(K_i)/[K_i:\mathbb Q]$ tends to some
$\tau\in[0,1]$.  Every $p\in P$ splits completely in $K_i/\mathbb Q$
because it splits in $K$ and each prime above it splits in $K_i/K$.
$M$ supplies the Dedekind RH hypotheses of NT-H07.  That result yields
$\varrho\ge e^{\pi\tau/2}C(P)\ge C(P)$, a contradiction.
In the totally real version all $K_i$ are totally real, so $\tau=1$
and the same proof gives the second threshold.  This completes the proof.

For $P=\varnothing$, the thresholds are
$8\pi e^\gamma\simeq44.763232$ and
$8\pi e^{\gamma+\pi/2}\simeq215.332516$.
The group $\operatorname{Gal}(K^{\rm ur,P}/K)$ is the quotient of
$\pi_1^{\text{ét}}(\operatorname{Spec}\mathcal O_K)$ obtained by
killing the closed normal subgroup generated by Frobenius elements at
the primes above $P$.  Requiring total reality also kills the
archimedean complex-conjugation classes.  This is the finite-\'{e}tale
Galois-category dictionary, SGA~1, Expos\'e~V, \S7.  The discriminant
criterion therefore gives a concrete conditional finiteness statement
for this quotient.  It is a restriction on arithmetic towers in the
cote-149/GT comparison programme, not a claim of a new anabelian
reconstruction theorem.

\subsection{Prime races: zero location versus relations among ordinates}

\paragraph{NT-R01 and NT-R02. What $A$ alone supplies.}
Fix distinct reduced classes $a_1,\ldots,a_r$ modulo $q$, and set
\[
 E_{q;\boldsymbol a}(x)=\frac{\log x}{\sqrt x}
 \bigl(\varphi(q)\pi(x;q,a_1)-\pi(x),\ldots,
       \varphi(q)\pi(x;q,a_r)-\pi(x)\bigr).
\]
Rubinstein--Sarnak, \emph{Chebyshev's Bias} (1994), Theorem~1.1,
printed p.~174, proof \S2.1, gives under $A$ a probability measure
$\mu_{q;\boldsymbol a}$ such that, for every bounded continuous $f$,
\[
 \lim_{X\to\infty}\frac1{\log X}
       \int_2^X f(E_{q;\boldsymbol a}(x))\frac{dx}{x}
   =\int_{\mathbb R^r}f\,d\mu_{q;\boldsymbol a}.
\tag{NT-R01}
\]
Its normalization uses $\pi(x)$ in each coordinate.  Theorem~1.2,
printed p.~175, proof \S2.2, and the discussion following it imply
\[
 \underline\delta\{x:\pi(x;4,3)>\pi(x;4,1)\}>0,
 \qquad
 \underline\delta\{x:\pi(x;4,1)>\pi(x;4,3)\}>0,
\tag{NT-R02}
\]
where $\underline\delta S=\liminf_{X\to\infty}
(\log X)^{-1}\int_{S\cap[2,X]}dx/x$.
The same conclusion holds for the two classes modulo three.
These statements concern logarithmic sampling.  The existence of the
limiting distribution for continuous tests does not by itself justify
substituting an ordering-cone indicator: its boundary could have mass.

\paragraph{NT-R03 and NT-R04. The separate ordinate-independence input.}
Use the source's Grand Simplicity Hypothesis (GSH): the nonnegative
ordinates of zeros of all primitive Dirichlet functions, listed with
multiplicity, are linearly independent over $\mathbb Q$.  In particular
this excludes central zeros and repeated ordinates.  GSH remains an
additional hypothesis here.

Under $A$ and GSH, for $r\ge2$ every strict ordering race has a
logarithmic density in $(0,1)$; the source's formula (1.2) and
Remark~1.3, printed p.~176, give the requisite regularity and positive
mass.  For fixed $r$,
\[
 \max_{\substack{a_1,\ldots,a_r\pmod q\\
                   \text{distinct reduced classes}}}
 \left|\delta\{x:\pi(x;q,a_1)>\cdots>\pi(x;q,a_r)\}
                   -\frac1{r!}\right|\longrightarrow0
       \quad(q\to\infty).
\tag{NT-R04}
\]
This is Theorem~1.5, printed p.~177, proof \S3.2.
For the exact symmetry classification, put
$c(q,a)=-1+\#\{b\pmod q:b^2\equiv a\pmod q\}$.
Theorem~1.4, printed p.~177, proof Proposition~3.1, says that the
limiting law is invariant under every coordinate permutation precisely
in the following cases: $r=2$ and $c(q,a_1)=c(q,a_2)$; or $r=3$
and, after labelling, the classes are
$a_1,a_1\rho,a_1\rho^2$ with $\rho\ne1$ and $\rho^3=1\pmod q$.
This is symmetry of the full law.  It must not be conflated with a
converse assertion about one numerical race density.  The discussion
on p.~177 gives, in particular, the strict bias
$1/2<\delta\{x:\pi(x;4,3)>\pi(x;4,1)\}<1$ under these extra
hypotheses.  The relevant original pp.~174--178 were visually inspected
because their PDF text encoding is defective.

Geometrically these are races between Frobenius classes in cyclotomic
covers of $\operatorname{Spec}\mathbb Z[1/q]$.  Their limiting
distribution is an explicit-formula statistic of the corresponding
rank-one objects.  The additional GSH gate concerns relations among
zero ordinates, so the zero-location inputs $A$ and $M$ do not discharge
it in this argument.  A finite-field monodromy comparison likewise
does not remove that number-field hypothesis.

\subsection{Multiplicity, gaps, and the first zero: quantitative limits}

\paragraph{NT-Z13. An argument bound with its actual archimedean hypothesis.}
Carneiro--Chandee--Milinovich, \emph{A note on the zeros of zeta and
$L$-functions} (2015), author PDF Theorem~1, p.~2, proves under RH
\[
 |S(t)|\le \frac14\frac{\log t}{\log\log t}
       +O\left(\frac{\log t\log\log\log t}{(\log\log t)^2}\right).
\tag{NT-Z13a}
\]
Here $S(t)=\pi^{-1}\arg\zeta(1/2+it)$, with the argument continued
from $2$ through $2+it$ and the mean of its two one-sided values used
at a zero.  The proof is in \S3, pp.~3--5.

We record the precise Hecke specialization of their Theorem~5,
p.~7, equations (4.4)--(4.5), whose proof occupies pp.~7--12.
Let $\eta$ be a primitive finite-order Hecke character of a number
field $K$ of degree $n$.  Write its completion, up to a nonzero
constant independent of $s$, as
\[
 \Lambda(s,\eta)=G_\eta(s)L(s,\eta),\qquad
 G_\eta(s)=Q_\eta^{s/2}\prod_{j=1}^n\Gamma_{\mathbb R}(s+\mu_j),
 \qquad Q_\eta=\Delta_K\operatorname N\mathfrak f_\eta,
\]
where $\Gamma_{\mathbb R}(s)=\pi^{-s/2}\Gamma(s/2)$.
Each real place contributes $\mu_j=0$ or $1$ according to the sign
character.  Each complex place contributes the pair $0,1$.
Set
\[
 C_\eta=Q_\eta\prod_{j=1}^n(\mu_j+3),\quad
 C_\eta(t)=C_\eta(|t|+1)^n,\quad
 D_\eta(t)=\log\log\bigl(C_\eta(t)^{3/n}\bigr).
\]
For $t>0$ define
$S_\eta(t)=-\pi^{-1}\int_{1/2}^{\infty}
\operatorname{Im}(L'/L)(\sigma+it,\eta)\,d\sigma$ away from zero
ordinates, and use one-sided means at them.  Under RH for this
particular $L$-function,
\[
 |S_\eta(t)|\le\frac{\log C_\eta(t)}{4D_\eta(t)}
   +O\left(\frac{\log C_\eta(t)\log D_\eta(t)}{D_\eta(t)^2}\right).
\tag{NT-Z13b}
\]
The implied constant is absolute.  In particular, for fixed $K,\eta$,
\[
 |S_\eta(t)|\le\frac n4\frac{\log t}{\log\log t}
   +O_{K,\eta}\left(
       \frac{\log t\log\log\log t}{(\log\log t)^2}\right).
\tag{NT-Z13c}
\]

Here is the check of the source hypotheses, which matters for $M$.
After grouping prime ideals by their underlying rational prime, the
Hecke Euler factor is a product of at most $n$ linear reciprocal
factors with local parameters of modulus at most one: factor each
$1-\eta(\mathfrak p)p^{-fs}$ into its $f$ roots and pad missing
factors by zero parameters.  Thus the source's parameter $\vartheta$
is zero.  The completion has order one, the standard functional
equation with the conjugate character, and is entire except for the
trivial character's poles at $0,1$.  Its listed gamma parameters are
real and nonnegative.  These are precisely the analytic conditions
in \S4, p.~6, supplied by classical Hecke theory (RV1999,
Theorem~7-19).  $M$ supplies the remaining zero-location assumption;
$A$ suffices for Dirichlet characters over $\mathbb Q$.
In particular, the source explicitly requires gamma parameters to be
real or to occur in conjugate pairs.  We do not silently use it for
an arbitrary infinite-order unitary Hecke character.  The paired
$\pm it_k$ family in root MT-H01--H03 can be checked separately
against this same hypothesis.

\paragraph{NT-Z14. Full deduction of multiplicity and gap bounds.}
For each fixed $K,\eta$ as in NT-Z13, under its GRH, a zero
$1/2+i\gamma$ with $\gamma\to+\infty$ has multiplicity
\[
 m_\eta(\gamma)\le\frac n2\frac{\log\gamma}{\log\log\gamma}
   +O_{K,\eta}\left(
    \frac{\log\gamma\log\log\log\gamma}{(\log\log\gamma)^2}\right).
\tag{NT-mult}
\]
For consecutive distinct sufficiently large positive ordinates
$\gamma<\gamma'$,
\[
 \gamma'-\gamma\le\frac{\pi}{\log\log\gamma}
   +O_{K,\eta}\left(
    \frac{\log\log\log\gamma}{(\log\log\gamma)^2}\right).
\tag{NT-gap}
\]
Consequently every sufficiently large real $T$ lies within
\[
 \frac{\pi}{2\log\log T}
 +O_{K,\eta}\left(\frac{\log\log\log T}{(\log\log T)^2}\right)
\]
of a zero ordinate.  The analogous
negative-height statements follow by conjugating $\eta$.  These are
classical consequences in a fixed-character formulation; they assert
neither simplicity nor pair correlation.

We give the entire deduction, including the jump and endpoint issues.
In a sufficiently high horizontal region $L(s,\eta)$ has no pole;
GRH and the boundary nonvanishing theorem give no zero to the right
of the critical line.  Continue $\log L$ from its Euler-product
normalization there.  Near $\rho=1/2+i\gamma$ write
$L(s,\eta)=(s-\rho)^m U(s)$ with $U$ holomorphic and nonzero.
Approaching the critical line from its right, the argument of
$s-\rho$ changes from $-\pi/2$ to $\pi/2$ across the ordinate.
Hence
\[
 S_\eta(\gamma+)-S_\eta(\gamma-)=m.
\]
Taking one-sided limits in NT-Z13c bounds this jump by twice its
right-hand side and proves (NT-mult).

Let $\Theta_\eta(t)=\operatorname{Im}\log G_\eta(1/2+it)$, with
any continuous determination at large positive $t$.
The functional equation and conjugation give
$\Lambda(1/2+it,\eta)=\varepsilon_\eta
\overline{\Lambda(1/2+it,\eta)}$.
Therefore, on an interval containing no zero, the argument of the
completed function is constant modulo $\pi$, and
\[
 S_\eta'(t)=-\frac1\pi\Theta_\eta'(t),\qquad
 \frac1\pi\Theta_\eta'(t)
    =\frac n{2\pi}\log t+O_{K,\eta}(1).
\tag{NT-phase}
\]
The derivative formula is Stirling's formula applied to each of the
$n$ displayed $\Gamma_{\mathbb R}$ factors; the conductor term is
constant.  The same identity and NT-Z13c show that positive ordinates
are unbounded: if there were no zeros beyond a height, integration
would force $S_\eta(t)$ to have size comparable to $t\log t$.

For large $\gamma$, a zero-free gap cannot have length at least one.
Otherwise it contains an interval of length $1/2$ whose endpoints
are between $\gamma$ and $\gamma+1$.  Integrating (NT-phase) over it
gives a decrease of size $(n/(4\pi)+o(1))\log\gamma$, whereas
NT-Z13c bounds the difference of its endpoint values by
$O_{K,\eta}(\log\gamma/\log\log\gamma)$, a contradiction.
Thus $h=\gamma'-\gamma<1$ eventually.  Integrate (NT-phase) from
just above $\gamma$ to just below $\gamma'$ and then take endpoint
limits.  Uniformly on this interval,
\[
 \left(\frac n{2\pi}\log\gamma+O_{K,\eta}(1)\right)h
   =S_\eta(\gamma+)-S_\eta(\gamma'-)
   \le\frac n2\frac{\log\gamma}{\log\log\gamma}
      +O_{K,\eta}\left(
       \frac{\log\gamma\log\log\log\gamma}
            {(\log\log\gamma)^2}\right).
\]
Division proves (NT-gap).  Any $T$ between consecutive ordinates is
at distance at most half their gap from one endpoint.  Since that
gap is less than one, replacing $\gamma$ by $T$ in the slowly varying
bound changes only the error term.  This proves the last assertion
and completes the proof.  For $K=\mathbb Q$ and the trivial
character, $n=1$ gives the usual zeta multiplicity and gap bounds
under $A$ alone.

\paragraph{NT-Z15 and NT-Z16. Conductor aspect and the first ordinate.}
In the same finite-order Hecke setting put
$D_\eta=\log\log(C_\eta^{3/n})$.
Theorem~5, equation (4.5), gives
\[
 \operatorname{ord}_{s=1/2}L(s,\eta)
 \le\frac{\log C_\eta}{2D_\eta}
   +O\left(\frac{\log C_\eta\log D_\eta}{D_\eta^2}\right).
\tag{NT-Z15}
\]
There is no gamma pole or zero at $1/2$, so the order equals that
of the completion in the source.  Theorem~7, p.~13, with its full
proof on that page, gives at least one zero $1/2+i\gamma$ satisfying
\[
 |\gamma|\le\frac{\pi}{2D_\eta}
          +O\left(\frac{\log D_\eta}{D_\eta^2}\right)
 \quad\text{as } C_\eta^{3/n}\longrightarrow\infty.
\tag{NT-Z16}
\]
Both implied constants are universal.  A possible central zero is
included: this is not a bound for the first strictly positive
ordinate.  Keeping $n$ inside $D_\eta$ is essential when the degree
varies.  For primitive Dirichlet characters of conductor $q>2$, the
specializations are
\[
 \operatorname{ord}_{1/2}L(s,\chi)
       \le(1/2+o(1))\frac{\log q}{\log\log q},\qquad
 \min_\rho|\operatorname{Im}\rho|
       \le\frac{\pi/2+o(1)}{\log\log q}.
\]
Here $A$ supplies every required GRH hypothesis.  For finite-order
Hecke characters over general fields the corresponding input is
the finite-order part of $M$; it is not supplied by $A$ merely
because the base field happens to be abelian over $\mathbb Q$.

These bounds constrain the zero divisor of rank-one arithmetic
local systems and their Artin inductions.  In the Grothendieck
comparison language they are quantitative restrictions on the
archimedean explicit-formula spectrum.  They supply neither a
finite-field purity argument over $\operatorname{Spec}\mathbb Z$
nor a general geometric multiplicity-one theorem.  They refine the
specific open gates in NT-Z06 while leaving them open.

\subsection{Further prime constants and quadratic-family moments}

\paragraph{NT-P13. Explicit and asymptotic prime-gap constants.}
Under the RH part of $A$, every $x\ge4$ has a prime in
\[
 \left[x,x+\frac{22}{25}\sqrt x\log x\right].
\]
This is Carneiro--Milinovich--Soundararajan,
\emph{Fourier optimization and prime gaps} (2019), author v2
Theorem~5, p.~5; the explicit proof is \S6, pp.~21--26.
Consequently, for every consecutive pair with $p_n>3$,
$p_{n+1}-p_n\le(22/25)\sqrt{p_n}\log p_n$.
For clarity about the closed left endpoint, apply the interval
statement to $x=p_n+\varepsilon$ with
$0<\varepsilon<p_{n+1}-p_n$, then let $\varepsilon\downarrow0$.
Corollary~4 on the same page, following Theorem~3 and its proof in
\S5, strengthens the asymptotic constant:
\[
 \limsup_{n\to\infty}
       \frac{p_{n+1}-p_n}{\sqrt{p_n}\log p_n}<\frac{21}{25}.
\]
The explicit all-$x$ statement and the sharper asymptotic statement
have different quantifiers.  No claim is made that these are the
latest possible constants.

\paragraph{NT-P14--NT-P16. Fourier improvements in the conductor aspect.}
Carneiro--Milinovich--Quesada-Herrera--Ramos,
\emph{Fourier optimization, the least quadratic non-residue, and
the least prime in an arithmetic progression}, arXiv v2 dated
August 2025, Theorems~1 and 2, author p.~2, gives under $A$
\[
 \limsup_{p\to\infty}\frac{n_p}{(\log p)^2}<0.7615,
 \qquad
 \limsup_{p\to\infty}\frac{r_p}{(\log p)^2}<0.7615.
\tag{NT-P14}
\]
Here $p$ runs over primes, $n_p$ is the least positive quadratic
nonresidue modulo $p$, and $r_p$ is the least prime $r$ with
$(r/p)=1$.  The nonresidue is itself prime.  Theorem~1 follows from
Theorem~6, whose proof is \S4, pp.~13--17; the modification proving
Theorem~2 is \S4.4, pp.~16--17.  These bounds improve the
asymptotic constant; they do not replace NT-P07's explicit threshold
$p\ge5$ by an unstated numerical threshold.

More generally, for a nonprincipal character $\chi$ modulo $q$ of
fixed order $\ell\ge2$, define
$n_\chi=\min\{n\ge1:\chi(n)\notin\{0,1\}\}$.
Theorem~6, p.~5, and Corollary~7, p.~6, give
\[
 n_\chi\le(c_\ell+o_\ell(1))(\log q)^2,
 \qquad
 \begin{array}{c|rrrr}
   \ell&2&3&4&5\\ \hline
   c_\ell&0.7615&0.6707&0.6131&0.5765
 \end{array}
\tag{NT-P15}
\]
as $q\to\infty$.  For fixed $\ell\ge6$ the same corollary permits
\[
 c_\ell=\frac14\left(
  1-\frac{\frac{\ell}{\ell-1}\log\bigl((3\ell-4)/\ell\bigr)}
             {\log(\ell-1)}+\frac1{\ell-1}\right)^{-2}.
\]
The character need not be primitive.  Its values zero at primes
dividing $q$ are explicitly excluded in the definition.  If
$n_\chi$ were composite, every prime divisor would be smaller and
coprime to $q$, hence would have character value one by minimality;
multiplicativity would give $\chi(n_\chi)=1$, a contradiction.
Thus these are also bounds for the least prime outside the kernel
of the cyclotomic character.  The limit $c_\ell\to1/4$ is a
statement about these constants; we do not assert uniformity when
$\ell$ varies with $q$.

For the least prime $P(a,q)\equiv a\pmod q$, the same paper's
Theorem~3, author p.~3, proof \S5, pp.~18--22, gives under $A$
\[
 \limsup_{q\to\infty}
  \max_{(a,q)=1}\frac{P(a,q)}{(\varphi(q)\log q)^2}<\frac89.
\tag{NT-P16}
\]
The uniformity in the reduced class follows from the proof: its
errors in (5.15)--(5.16) depend on the fixed Fourier test function
and not on $a$.  The numerical optimization certification in the
paper is cited, not rerun here.  These results concern closed points
and selected Frobenius classes of cyclotomic covers.  They sharpen
the effective detection bounds in the SGA~1 Galois-category
dictionary; they do not enlarge the class of $L$-functions furnished
by $A$.

\paragraph{NT-L11. Quadratic central moments with the full family specified.}
For each fixed real $k\ge0$, $A$ and Harper,
\emph{Sharp conditional bounds for moments of the Riemann zeta
function}, Theorem~2, author p.~15, give
\[
 \sum_{\substack{X\le |d|\le2X\\d\ {
m fundamental}}}
       |L(1/2,\chi_d)|^k
     \ll_k X(\log X)^{k(k+1)/2}.
\tag{NT-L11}
\]
The family includes positive and negative fundamental discriminants;
$\chi_d$ is the associated primitive quadratic character.  The
argument is in \S5, pp.~15--17, using Propositions~3 and 4.
The input is only GRH for the quadratic functions in question.
This is an upper bound of the predicted order, not an asymptotic
formula with its leading constant.  In particular it asserts no
bound for negative moments or reciprocals of central values.

\paragraph{NT-L12. Full fixed-threshold large-value deductions.}
For every fixed $\alpha>0$, $A$ implies
\[
 \operatorname{meas}\{t\in[T,2T]:
            |\zeta(1/2+it)|\ge(\log T)^\alpha\}
       \ll_\alpha \frac{T}{(\log T)^{\alpha^2}}.
\tag{NT-L12a}
\]
For every fixed $\alpha>1/2$ it also implies
\[
 \#\{d\text{ fundamental}:X\le|d|\le2X,
          |L(1/2,\chi_d)|\ge(\log X)^\alpha\}
       \ll_\alpha
          \frac{X}{(\log X)^{(\alpha-1/2)^2/2}}.
\tag{NT-L12b}
\]
Here is the complete proof.  For any fixed $k>0$, Markov's
inequality and Harper's Theorem~1 (NT-L04) bound the first measure by
\[
 (\log T)^{-2k\alpha}
       \int_T^{2T}|\zeta(1/2+it)|^{2k}\,dt
       \ll_k T(\log T)^{k^2-2k\alpha}.
\]
Choose the fixed value $k=\alpha$ to obtain (NT-L12a).
Similarly NT-L11 bounds the second count by
$O_k(X(\log X)^{k(k+1)/2-k\alpha})$.
For $\alpha>1/2$, choose the positive fixed value
$k=\alpha-1/2$; completing the square gives
$k(k+1)/2-k\alpha=-(\alpha-1/2)^2/2$.
This proves (NT-L12b).  All constants may depend on the fixed
threshold exponent; no choice of $k$ growing with $T$ or $X$ has
been made.  Central zeros cause no difficulty because only positive
moments occur.  The result is a direct optimization of the cited
moment bounds, not a new moment theorem.

The quadratic family consists of rank-one Artin objects coming from
quadratic covers with growing discriminant.  Its central-value
moment exponent is an arithmetic family statement.  A geometric
symmetry-type comparison is useful context, but is not substituted
for Harper's analytic proof or the exact GRH input.

\subsection{Closed-point Frobenius relations in unramified towers}

\paragraph{NT-H09. A relative splitting-degree budget.}
Let $K\ne\mathbb Q$ be a number field of degree $n$, discriminant
$\Delta_K$, and root discriminant $\varrho=\Delta_K^{1/n}$.
Let $S$ be a finite set of prime ideals of $K$, with
$q_{\mathfrak p}=\operatorname N\mathfrak p$, and assign to each
$\mathfrak p\in S$ a positive integer $m_{\mathfrak p}$.
Suppose that $K_i/K$ are finite Galois extensions unramified at all
finite primes, $d_i=[K_i:K]\to\infty$, and the common residue
degree of primes above $\mathfrak p$ in $K_i/K$ is at most
$m_{\mathfrak p}$ for every $i$.  Assume
$r_1(K_i)/[K_i:\mathbb Q]\to\tau$ and RH for every
$\zeta_{K_i}$.  Then
\[
 \log\varrho\ge\log(8\pi)+\gamma+\frac\pi2\tau
  +\frac2n\sum_{\mathfrak p\in S}
       \frac{\log q_{\mathfrak p}}
            {q_{\mathfrak p}^{m_{\mathfrak p}/2}-1}.
\tag{NT-H09}
\]
Thus $M$ supplies the needed zero hypotheses.  No nesting of the
extensions is required.  Complete splitting is the case
$m_{\mathfrak p}=1$ and does not require the rational prime below
$\mathfrak p$ to split in $K$.

Here is the full proof.  Pass to strictly increasing degrees so that
the fields are pairwise nonisomorphic.  Finite unramifiedness gives
$\Delta_{K_i}=\Delta_K^{d_i}$ and
$g_i=\tfrac12 d_i\log\Delta_K$.
There are finitely many possible residue degrees at each point of
the finite set $S$, so pass to a subsequence on which these degrees
are fixed, say $f_{\mathfrak p}\le m_{\mathfrak p}$.
As in NT-H07, boundedness of $N_Q(K_i)/g_i$ and a diagonal
subsequence over rational prime powers $Q$ give an asymptotically
exact family.  Neither the limiting signature nor the displayed
residue degrees change.  Since the extension is Galois and
unramified, precisely $d_i/f_{\mathfrak p}$ primes lie above
$\mathfrak p$, each of norm $q_{\mathfrak p}^{f_{\mathfrak p}}$.
Different base prime ideals give disjoint sets of primes upstairs,
even when their absolute norms coincide.  Therefore
\[
 \phi_Q\ge\frac2{\log\Delta_K}
       \sum_{\substack{\mathfrak p\in S\\
                    q_{\mathfrak p}^{f_{\mathfrak p}}=Q}}
           \frac1{f_{\mathfrak p}}.
\]
The archimedean invariants are
$\phi_{\mathbb R}=2n\tau/\log\Delta_K$ and
$\phi_{\mathbb C}=n(1-\tau)/\log\Delta_K$.
Substitute these lower bounds in the TV inequality (NT-basic),
discard the other nonnegative terms, and multiply by
$\log\Delta_K/n$.  Since
$\log(q_{\mathfrak p}^{f_{\mathfrak p}})/f_{\mathfrak p}
=\log q_{\mathfrak p}$, the finite-place contribution is at least
\[
 \frac2n\sum_{\mathfrak p\in S}
      \frac{\log q_{\mathfrak p}}
           {q_{\mathfrak p}^{f_{\mathfrak p}/2}-1}
 \ge\frac2n\sum_{\mathfrak p\in S}
      \frac{\log q_{\mathfrak p}}
           {q_{\mathfrak p}^{m_{\mathfrak p}/2}-1}.
\]
The archimedean part simplifies using
$2\alpha_{\mathbb R}-\alpha_{\mathbb C}=\pi/2$ exactly as in
NT-H07.  This proves (NT-H09), including the possibility of
coincident prime-power norms.

\paragraph{NT-H10. A finite quotient from bounded-order Frobenius relations.}
Let
$G_K=\pi_1^{\text{ét}}(\operatorname{Spec}\mathcal O_K)$, and
for the data $S,\boldsymbol m$ above let $N$ be the closed normal
subgroup generated by
$\operatorname{Frob}_{\mathfrak p}^{m_{\mathfrak p}}$ for
$\mathfrak p\in S$.  The normal subgroup is independent of the
choices inside the Frobenius conjugacy classes.  Put
\[
 C_K(S,\boldsymbol m)=8\pi\exp\left(
       \gamma+\frac2n\sum_{\mathfrak p\in S}
       \frac{\log q_{\mathfrak p}}
            {q_{\mathfrak p}^{m_{\mathfrak p}/2}-1}\right).
\]
Under $M$, if $\varrho<C_K(S,\boldsymbol m)$, then $G_K/N$ is
finite; equivalently, $N$ is open.  For totally real $K$, first
quotient $G_K$ by the closed normal subgroup generated by complex
conjugations.  In that totally real quotient the weaker condition
$\varrho<e^{\pi/2}C_K(S,\boldsymbol m)$ suffices.  No equality
case is asserted.

To prove this, suppose the stated quotient is infinite.  Its finite
Galois subextensions $K_i/K$ can then be chosen with degrees tending
to infinity.  They are unramified at every finite prime.  In each
finite Galois group the order of
$\operatorname{Frob}_{\mathfrak p}$ divides
$m_{\mathfrak p}$, and that order equals its residue degree.
Choose a subsequence whose signature ratios tend to
$\tau\in[0,1]$.  NT-H09 gives
$\varrho\ge e^{\pi\tau/2}C_K(S,\boldsymbol m)
\ge C_K(S,\boldsymbol m)$, contradicting the strict inequality.
In the totally real quotient every $K_i$ is totally real, so
$\tau=1$ and the second conclusion follows.  If $K=\mathbb Q$,
Minkowski's discriminant bound already makes the whole unramified
group trivial, so this omitted case satisfies the conclusion
unconditionally.  This completes the proof.

For $m_{\mathfrak p}=1$ this is the unrestricted closed-point
version of NT-H08.  If $S$ consists of all primes above rational
primes that split completely in $K$, then each such rational prime
contributes $n$ equal summands and NT-H08's constant is recovered.
For larger $m_{\mathfrak p}$ the imposed relations only bound the
orders of Frobenius elements.  The conclusion concerns the index
of their closed normal closure; it does not assert that this normal
subgroup is generated by finitely many elements as an ordinary
topological group.  This is a quantitative Galois-category
consequence of the TV explicit-formula inequality, with its precise
arithmetic input recorded, rather than a general anabelian
reconstruction statement.

\subsection{Symmetrization removes the infinity-type restriction for zero bounds}

\paragraph{NT-Z17. Uniform zero bounds for arbitrary unitary Hecke infinity type.}
Let $\eta$ be a primitive unitary Hecke character of a number field
$K$, $n=[K:\mathbb Q]$, which is not a pure norm character.
Write its standard completion as
\[
 \Lambda(s,\eta)=Q_\eta^{s/2}
         \prod_{j=1}^n\Gamma_{\mathbb R}(s+\mu_j)L(s,\eta),
 \qquad Q_\eta=\Delta_K\operatorname N\mathfrak f_\eta,
\]
up to a harmless nonzero constant.  Unitarity gives
$\operatorname{Re}\mu_j\ge0$, but the $\mu_j$ need not occur
in conjugate pairs.  For any real $T$, put
\[
 \mathscr C_\eta(T)=Q_\eta\prod_{j=1}^n(|\mu_j+iT|+3),
 \qquad
 R_\eta(T)=\log\log\bigl(\mathscr C_\eta(T)^{3/n}\bigr).
\]
Under GRH for $L(s,\eta)$, in particular under $M$,
\[
 \operatorname{ord}_{s=1/2+iT}L(s,\eta)
 \le\frac{\log\mathscr C_\eta(T)}{2R_\eta(T)}
  +O\left(\frac{\log\mathscr C_\eta(T)\log R_\eta(T)}
                 {R_\eta(T)^2}\right),
\tag{NT-Z17a}
\]
with an absolute implied constant.  Also, as
$\mathscr C_\eta(T)^{3/n}\to\infty$, there is a nontrivial zero
$1/2+i\gamma$ of this same $L$-function such that
\[
 |\gamma-T|\le\frac{\pi}{2R_\eta(T)}
                 +O\left(\frac{\log R_\eta(T)}{R_\eta(T)^2}\right),
\tag{NT-Z17b}
\]
again with an absolute implied constant.  The degree-normalized
conductor is retained, so these are uniform formulations even when
$K,\eta,T$, and $n$ vary within the stated class.

Here is the complete transfer proof.  Define the product and its
completion by
\[
 F_T(s)=L(s+iT,\eta)L(s-iT,\overline\eta),\qquad
 \mathcal F_T(s)=\Lambda(s+iT,\eta)
                         \Lambda(s-iT,\overline\eta).
\]
Classical Hecke theory makes both completed factors entire of order
one precisely because the pure norm exception has been excluded.
The gamma parameters of the product are
$\mu_j+iT$ and $\overline\mu_j-iT$, in conjugate pairs.
The factors $Q_\eta^{iT/2}$ and $Q_\eta^{-iT/2}$ cancel, so its
finite conductor is $Q_\eta^2$ and its degree is $2n$.
After grouping prime ideals by rational primes, the Euler parameters
of either shifted factor have modulus at most one.  Thus the product
has the parameter $\vartheta=0$ in CCM2015, \S4.  Conjugation and
the Hecke functional equation show that it is self-dual with root
number one.  GRH for $\eta$ implies GRH for its conjugate and for
both shifts, hence for the product.  It therefore satisfies all
the hypotheses of CCM2015 Theorems~5 and 7, including the actual
gamma-pairing condition.

Its analytic conductor in that source is exactly
\[
 C(\mathcal F_T)=Q_\eta^2
       \prod_{j=1}^n(|\mu_j+iT|+3)^2
       =\mathscr C_\eta(T)^2.
\]
Consequently its denominator is
\[
 \log\log\bigl(C(\mathcal F_T)^{3/(2n)}\bigr)
                  =R_\eta(T),
\]
with no alteration of degree dependence.
Conjugation gives
$\operatorname{ord}_{1/2-iT}L(s,\overline\eta)
=\operatorname{ord}_{1/2+iT}L(s,\eta)$.
The gamma factors are finite and nonzero at the points under
consideration, since their arguments have positive real part.
Thus the central order of $\mathcal F_T$ is twice the order on
the left of (NT-Z17a).  Apply CCM Theorem~5, equation (4.5): the
factor two in $\log C(\mathcal F_T)$ is cancelled by division
by this doubled order.  This gives (NT-Z17a), with a universal
error constant as stated.

By CCM Theorem~7 there is a zero $1/2+iu$ of $\mathcal F_T$
with the bound on $|u|$ appearing in (NT-Z17b).  If its first
factor vanishes, $L(s,\eta)$ has ordinate $T+u$.
If its second factor vanishes, conjugating that zero gives an
ordinate $T-u$ for $L(s,\eta)$.  Either case proves (NT-Z17b).
This argument also permits $u=0$, as it should.

For each fixed non-pure-norm $K,\eta$,
$\log\mathscr C_\eta(T)=n\log|T|+O_{K,\eta}(1)$ as
$|T|\to\infty$.  Substitution in the two bounds gives exactly
the multiplicity and centered covering-radius estimates of
NT-Z14 for this arbitrary infinity type.  Applying the covering
bound at the midpoint of a large gap first shows that its length
is less than one, and then gives
\[
 \gamma'-\gamma\le\frac{\pi}{\log\log\gamma}
    +O_{K,\eta}\left(
        \frac{\log\log\log\gamma}{(\log\log\gamma)^2}\right)
\]
for consecutive large positive ordinates.  If $\eta$ is a pure
norm character, its $L$-function is $\zeta_K(s+i\tau)$ for a
fixed real $\tau$.  Translate the finite-order trivial-character
case of NT-Z14 to obtain the same fixed-character asymptotic
conclusions.  We do not apply CCM verbatim to its two shifted
completed zeta factors: they have poles away from $0,1$ and would
violate that theorem's stated hypotheses.

This proves a zero-divisor transfer, not an individual argument
bound for $S_\eta(T)$ with unpaired gamma parameters.  The product
is the conjugate-dual pair of a rank-one arithmetic object and its
norm twist.  It provides a precise analytic use of that symmetry
without presupposing an algebraic motive for the infinite-order
character; compare root MT-H01 and the geometry programme's
self-dual determinant constructions.
When $L(s,\eta)=L(s+i\tau,\chi)$ for a primitive Dirichlet
character $\chi$, $A$ already supplies the needed zero-location
input by translation.

\paragraph{NT-Z18. At least half of the prime-conductor character family is nonvanishing.}
Let $q\to\infty$ run over primes and average over all $q-2$
nonprincipal characters modulo $q$.  Under $A$,
\[
 \frac1{q-2}\sum_{\chi\ne\chi_0}
             \operatorname{ord}_{s=1/2}L(s,\chi)
       \le\frac12+O\left(\frac1{\log q}\right),
\]
and therefore
\[
 \#\{\chi\ne\chi_0\pmod q:L(1/2,\chi)\ne0\}
       \ge\left(\frac12-O\left(\frac1{\log q}\right)\right)(q-2).
\tag{NT-Z18}
\]
For the full proof take
$f(x)=(\sin(2\pi x)/(2\pi x))^2$, with $f(0)=1$.
This is real, even, nonnegative, entire, and decays as
$O((1+|x|)^{-2})$ on the real line.  Its Fourier transform is
$\widehat f(u)=\tfrac12(1-|u|/2)_+$, supported on $[-2,2]$,
and $\int f=1/2$.  It meets Hughes--Rudnick2003 Definition~2.1,
author p.~4, including the analytic-strip condition in \S2.1.
Their Theorem~3.1, p.~5, permits this closed support and gives
average $W_f=1/2+O(1/\log q)$.
Under $A$ all the scaled ordinates in $W_f$ are real, so every
summand is nonnegative; the central zeros alone contribute their
multiplicity because $f(0)=1$.  This proves the mean-order bound.
The indicator of central vanishing is at most that multiplicity,
which proves (NT-Z18) after summing.  No parity assumption on
central multiplicity is used for this non-self-dual family.

The exact family and the positivity use of $A$ distinguish this
unitary-family corollary from the quadratic-family statement
NT-Z12.  In geometric language it concerns the central zero divisor
of all nontrivial rank-one objects of the prime cyclotomic cover;
it does not assert nonvanishing for each object.

\subsection{Classwise prime counts and the average conductor}

\paragraph{NT-C06. Explicit prime counts in a ray class.}
Let $K$ have degree $n$ and absolute discriminant $\Delta_K$.
Let $\mathfrak m=\mathfrak m_0\mathfrak m_\infty$ be a modulus,
$G=\operatorname{Cl}_{\mathfrak m}(K)$ its ray class group, and
$h=|G|$.  Write $H/K$ for its ray class field.  For
$\chi\in\widehat G$, let $\mathfrak f_\chi$ be the finite part
of the conductor of its primitive Hecke character.  Define
\[
 a_{\mathfrak m}=\log\Delta_K+
           \frac1h\sum_{\chi\in\widehat G}
                    \log\operatorname N\mathfrak f_\chi.
\]
The following unconditional relations will be used:
\[
 a_{\mathfrak m}=\frac{\log\Delta_H}{h},\qquad
 0\le a_{\mathfrak m}\le
 \log\Delta_K+\left(1-\frac1h\right)
                     \log\operatorname N\mathfrak m_0.
\tag{NT-ray-conductor}
\]
The lower inequality is only a convenient bound; the exact expression
retains the mean of the individual conductor logarithms.

For a class $c\in G$, put
\begin{align*}
 \psi_c^{\mathfrak m}(x)
 &=\sum_{\substack{\mathfrak p\nmid\mathfrak m_0,\ j\ge1\\
           (\operatorname N\mathfrak p)^j\le x,\ [\mathfrak p]^j=c}}
                  \log\operatorname N\mathfrak p,\\
 \pi_c^{\mathfrak m}(x)
 &=\#\{\mathfrak p\nmid\mathfrak m_0:
             \operatorname N\mathfrak p\le x,\ [\mathfrak p]=c\}.
\end{align*}
Let $e_{\mathfrak m}$ be the number of prime ideals dividing
$\mathfrak m_0$ which are unramified in $H/K$; in particular,
$e_{\mathfrak m}\le\omega_K(\mathfrak m_0)$.
Under RH for $\zeta_H$, and hence under $M$, one has
\begin{align*}
 \left|\psi_c^{\mathfrak m}(x)-\frac{x}{h}\right|
 &\le\sqrt x\left[
       \left(\frac{\log x}{2\pi}+2\right)a_{\mathfrak m}
        +\left(\frac{\log^2x}{8\pi}+2\right)n\right]
       +e_{\mathfrak m}\log x,
       &&x\ge1,                                      \tag{NT-C06a}\\
 \left|\pi_c^{\mathfrak m}(x)-\frac{\operatorname{Li}_2(x)}h\right|
 &\le\sqrt x\left[
       \left(\frac1{2\pi}+\frac3{\log x}\right)a_{\mathfrak m}
       +\left(\frac{\log x}{8\pi}+\frac1{4\pi}
                         +\frac6{\log x}\right)n\right]
       +e_{\mathfrak m},
       &&x\ge2.                                      \tag{NT-C06b}
\end{align*}
Here $\operatorname{Li}_2(x)=\int_2^xdt/\log t$.
Both estimates are uniform in the class $c$.
Their errors contain no multiplicative ray class number: its role
is retained in the main term and in the average conductor.
Replacing $a_{\mathfrak m}$ by the upper bound in
(NT-ray-conductor) gives a version using only
$K,\mathfrak m_0,h$.

We give the full deduction, with the exact classical inputs identified.
Class-field reciprocity and existence identify $G$ with
$\operatorname{Gal}(H/K)$ and the class of each
$\mathfrak p\nmid\mathfrak m_0$ with its Frobenius;
see Milne, \emph{Class Field Theory}, v4.03, V, Theorems~3.5--3.6,
printed p.~158/PDF p.~167.  The finite conductors of the associated
characters divide $\mathfrak m_0$, and the trivial character has
conductor one.  The conductor--discriminant formula gives
\[
 \Delta_H=\Delta_K^h
              \prod_{\chi\in\widehat G}
                      \operatorname N\mathfrak f_\chi.
\]
Its proof locator is Neukirch, \emph{Algebraic Number Theory},
VII, (11.7)--(11.9), printed pp.~533--534/PDF pp.~541--542;
(11.10), printed p.~535/PDF p.~543, identifies the degree-one
Artin and class-field conductors.  The complete supplied proof of
the conductor--discriminant formula on those pages was read on
the original images.  Taking logarithms proves the equality in
(NT-ray-conductor).  Bounding each of the $h-1$ nontrivial
conductors by $\mathfrak m_0$ proves its upper bound.

Every conjugacy class of $\operatorname{Gal}(H/K)$ is a singleton,
so its density is $1/h$.  Moreover $[H:\mathbb Q]=nh$.
Apply NT-C01/C02, the precise Greni\'e--Molteni2019
Theorem~1.1 and Corollary~1.2, to that singleton.  Dividing their
errors by $h$ replaces $\log\Delta_H$ by
$a_{\mathfrak m}$ and $[H:\mathbb Q]$ by $n$, giving exactly
the bracketed expressions above.

Those source counts include every prime unramified in $H/K$.
All primes outside $\mathfrak m_0$ are unramified, but a prime
inside $\mathfrak m_0$ may be unramified as well.  Each such prime
contributes at most
\[
 \sum_{\substack{j\ge1\\(\operatorname N\mathfrak p)^j\le x}}
          \log\operatorname N\mathfrak p\le\log x
\]
to a prime-power count, and at most one to an ordinary prime
count.  Removing these $e_{\mathfrak m}$ possible primes and
using the triangle inequality proves (NT-C06a)--(NT-C06b), including
degenerate moduli.  This completes the proof.  For example, the
prime $2$ is unramified in the ray class field $\mathbb Q$ for
modulus $2\infty$ and is correctly removed by this boundary term.

The input is Dedekind RH for this one field $H$; no Artin
holomorphy conjecture is needed.  When $H/\mathbb Q$ is abelian,
$A$ suffices by Dirichlet factorization.  In particular, $A$
supplies every base-field-$\mathbb Q$ ray-class specialization.
For general $K$, abelianity of $K/\mathbb Q$ alone does not assert
abelianity of $H/\mathbb Q$ and is not substituted for this input.
Geometrically, the count takes place in the finite abelian
quotient of $\pi_1^{\text{ét}}(\operatorname{Spec}
\mathcal O_K[1/\mathfrak m_0])$ specified by the modulus and the
real-place conditions.  The quantity $a_{\mathfrak m}$ is the
average Artin conductor cost of its rank-one objects.

\paragraph{NT-C07. Ideal classes and degree-one prime representatives.}
Let $h_K=|\operatorname{Cl}(K)|$ and assume RH for
$\zeta_{\mathbb Q}$ and the Hilbert class-field Dedekind function.  Then every ordinary ideal class
contains a prime ideal of residue degree one over $\mathbb Q$
whose norm is at most
\[
 \left[h_K\bigl(4\log\Delta_K+2.5n\bigr)+5\right]^2.
\tag{NT-C07}
\]
For the narrow ideal class group the same conclusion holds with
$h_K^+$ in place of $h_K$, assuming RH for the corresponding
narrow class field.  $M$ supplies either hypothesis.
For either group, the classwise estimates (NT-C06a)--(NT-C06b)
hold with $a_{\mathfrak m}=\log\Delta_K$ and
$e_{\mathfrak m}=0$.

Indeed, use the modulus with finite part one and with empty real
part for the ordinary group, or all real places for the narrow
group.  Its class field $H/K$ is unramified at all finite primes;
consequently $\Delta_H=\Delta_K^{[H:K]}$.
The prime-count assertions follow at once from NT-C06.
For $K\ne\mathbb Q$ one has $H\ne\mathbb Q$, so
Bach--Sorenson1996 Theorem~5.1, printed p.~1729/PDF p.~13,
applies to each singleton of $\operatorname{Gal}(H/K)$.
Its degree-one prime has norm at most
$(4\log\Delta_H+2.5[H:\mathbb Q]+5)^2$.
Substitute $[H:K]=h_K$ or $h_K^+$ and the discriminant equality
to obtain (NT-C07).  The source requires RH for $\zeta_H$
and $\zeta_{\mathbb Q}$; $M$ provides both, and $A$ already
provides the latter.  If $K=\mathbb Q$, both groups are trivial
and the prime $2$ directly proves the displayed bound, without
applying the source theorem to its excluded field.
This proves the assertion in all cases.

The ordinary class group is
$\operatorname{Pic}(\operatorname{Spec}\mathcal O_K)$; the narrow
version adds the real sign conditions.  These are prime
representatives of each individual class, distinct from a small
set of primes generating the whole group.  No arbitrary extra
forbidden-prime set is included in (NT-C07); that additional
condition can instead be tested by the explicit count in NT-C06.

\paragraph{NT-C08. Certified ray-class Mertens constants and Euler products.}
Retain the notation and exact RH hypothesis of NT-C06, and let
$\delta=1/h$.  For each ray class $c$ there are real constants
$B_c$ and $A_c>0$ such that, for every $x\ge2$,
\begin{align*}
 \left|\sum_{\substack{\mathfrak p\nmid\mathfrak m_0\\
          [\mathfrak p]=c,\ \operatorname N\mathfrak p\le x}}
        \frac1{\operatorname N\mathfrak p}
        -\delta\log\log x-B_c\right|
       &\le U_{\mathfrak m}(x),                       \tag{NT-C08a}\\
 \left|\log\prod_{\substack{\mathfrak p\nmid\mathfrak m_0\\
          [\mathfrak p]=c,\ \operatorname N\mathfrak p\le x}}
       (1-(\operatorname N\mathfrak p)^{-1})^{-1}
       -\delta\log\log x-\log A_c\right|
       &\le U_{\mathfrak m}(x)+\frac{n}{2\lfloor x\rfloor},
                                                        \tag{NT-C08b}
\end{align*}
where the explicit error is
\[
 U_{\mathfrak m}(x)=\frac1{\sqrt x}
 \left[
   \left(\frac3{2\pi}+\frac9{\log x}\right)a_{\mathfrak m}
   +\left(\frac{3\log x}{8\pi}+\frac5{4\pi}
                    +\frac{18}{\log x}\right)n\right]
       +\frac{2e_{\mathfrak m}}x.
\tag{NT-C08-error}
\]
In particular the product divided by $(\log x)^{1/h}$ tends
to $A_c$.  The two finite expressions in (NT-C08a)--(NT-C08b),
with these errors, give certified intervals for $B_c$ and
$\log A_c$ once the relevant primes have been enumerated.
No uncomputed numerical value of either constant is asserted.

Here is the complete proof.  Abbreviate
$\pi_c^{\mathfrak m}(t)=\delta\operatorname{Li}_2(t)+R_c(t)$,
and let $S_c(x)$ be the reciprocal-prime sum in (NT-C08a).
Partial summation, including possible primes of norm $2$, gives
\[
 S_c(x)=\frac{\pi_c^{\mathfrak m}(x)}x+
                  \int_2^x\frac{\pi_c^{\mathfrak m}(t)}{t^2}\,dt.
\]
The lower boundary uses $\pi_c^{\mathfrak m}(2^-)=0$;
at $x=2$ the displayed formula is still valid.
Since $\operatorname{Li}_2(2)=0$, integration by parts yields
\[
 \frac{\operatorname{Li}_2(x)}x+
          \int_2^x\frac{\operatorname{Li}_2(t)}{t^2}\,dt
               =\log\log x-\log\log2.
\]
The estimate in NT-C06 makes $R_c(t)/t^2$ absolutely integrable
on $[2,\infty)$.  Define
\[
 B_c=-\delta\log\log2+
                       \int_2^\infty\frac{R_c(t)}{t^2}\,dt.
\]
These identities give the exact error formula
\[
 S_c(x)-\delta\log\log x-B_c
          =\frac{R_c(x)}x-
                        \int_x^\infty\frac{R_c(t)}{t^2}\,dt.
\tag{NT-C08-tail}
\]
Apply NT-C06b to both terms.  The elementary integrals needed are
\[
 \int_x^\infty t^{-3/2}\,dt=\frac2{\sqrt x},\qquad
 \int_x^\infty t^{-3/2}\log t\,dt
                      =\frac{2\log x+4}{\sqrt x},\qquad
 \int_x^\infty\frac{t^{-3/2}}{\log t}\,dt
                      \le\frac2{\sqrt x\log x}.
\]
For the $a_{\mathfrak m}$ terms the point and integral bounds
sum to $3/(2\pi)+9/\log x$ after division by $\sqrt x$.
For the $n$ terms they sum to
$3\log x/(8\pi)+5/(4\pi)+18/\log x$.
The boundary-prime contribution is
$e_{\mathfrak m}/x+\int_x^\infty e_{\mathfrak m}t^{-2}dt
=2e_{\mathfrak m}/x$.  This proves (NT-C08a) with exactly
the error displayed.

To pass to the product, define the absolutely convergent series
\[
 D_c=\sum_{\substack{\mathfrak p\nmid\mathfrak m_0\\
                         [\mathfrak p]=c}}
                 \sum_{k\ge2}\frac1{k(\operatorname N\mathfrak p)^k},
 \qquad A_c=\exp(B_c+D_c).
\]
For $q\ge2$,
$\sum_{k\ge2}1/(kq^k)\le1/(2q(q-1))$.
There are at most $n$ prime ideals of $K$ of any given
integer norm: if that norm is a rational prime power $p^f$,
there are at most $n/f\le n$ primes of that residue degree,
by the degree formula over $p$, and otherwise there are none.
Consequently
\[
 0\le\sum_{\substack{\mathfrak p\nmid\mathfrak m_0\\
              [\mathfrak p]=c,\ \operatorname N\mathfrak p>x}}
            \sum_{k\ge2}\frac1{k(\operatorname N\mathfrak p)^k}
 \le\frac n2\sum_{j>\lfloor x\rfloor}\frac1{j(j-1)}
 =\frac{n}{2\lfloor x\rfloor}.
\]
The logarithm of the finite Euler product is $S_c(x)$ plus
the corresponding truncated $k\ge2$ series.  Subtract
$\delta\log\log x+B_c+D_c$, use (NT-C08a) and this last
tail estimate, and obtain (NT-C08b).  Both errors tend to zero;
exponentiation proves the positive limiting constant and the
claimed interval certification.  This completes the proof.

The same deductions apply to ordinary and narrow ideal classes
with $a_{\mathfrak m}=\log\Delta_K$ and no boundary term.
For $K=\mathbb Q$ they are consequences of $A$ for every ray
class modulus.  The geometric object is the reciprocal-norm
measure on the closed points with a prescribed image in the
abelian fundamental-group quotient; its regularized Euler
product has a quantitatively controlled real constant.
This extends the finite prime counts to a convergent
regularization, without identifying that constant with a
particular motivic period.

\subsection{A conductor-uniform critical-line bound}

\paragraph{NT-L13. Uniform Lindel\"of bounds for non-pure-norm unitary Hecke characters.}
Let $\eta$ be a primitive unitary Hecke character of $K$ which is
not a pure norm character, and let $n=[K:\mathbb Q]$.
Use the completion and the sharp conductor already defined in
NT-Z17:
\[
 \mathscr C_\eta(T)=Q_\eta
                  \prod_{j=1}^n(|\mu_j+iT|+3),
 \qquad y=\log\log\mathscr C_\eta(T).
\]
Assume GRH for $L(s,\eta)$ and RH for $\zeta_{\mathbb Q}$.
Thus $M$ supplies both inputs; $A$ suffices for a norm shift
of a nonprincipal primitive Dirichlet character.
If $y\ge10$, then
\[
 |L(1/2+iT,\eta)|\le
  \exp\left[
    \left(\frac3{8y}+\frac{23n}{25y^2}\right)
                      \log\mathscr C_\eta(T)\right].
\tag{NT-L13}
\]
All constants are explicit and independent of $K,\eta,T,n$.
In particular, for every $\epsilon>0$ one has
$|L(1/2+iT,\eta)|\le\mathscr C_\eta(T)^\epsilon$ whenever
\[
 y\ge\max\left\{10,\frac3{4\epsilon},
                         \sqrt{\frac{46n}{25\epsilon}}\right\}.
\tag{NT-L13-threshold}
\]
This gives the conductor-uniform Lindel\"of statement for fixed
degree.  If the degree varies, the same conclusion in the form
$|L|\le\mathscr C^{o(1)}$ follows along any family for which
$y\to\infty$ and $n=o(y^2)$.  The dependence on degree is thus
specified rather than absorbed in an implied constant.

We verify the full specialization and the threshold deduction.
The classical input is Chandee2009, \emph{Explicit upper bounds
for $L$-functions on the critical line}, Corollary~1.1,
author p.~3, whose axioms and conductor are on p.~2,
equations (2)--(6).  Its proof is on p.~7, equations
(14)--(15); the underlying logarithmic inequality is proved
in Theorem~2.1, pp.~4--6.
Apply it to $s\mapsto L(s+iT,\eta)$.
The gamma parameters are $\mu_j+iT$ and have nonnegative
real part; this source does not impose the gamma-pairing
condition encountered in CCM2015.  Standard primitive Hecke
continuation and the functional equation give its analytic
axioms, and non-pure-norm unitary Hecke functions have no
completed pole or zero on the boundary lines.  The latter
nonvanishing is also part of the adopted GRH zero-location
input; its unconditional classical form is recorded in
Ramakrishnan--Valenza, Theorem~7-28, printed p.~289.

Grouping the prime-ideal Euler factors over each rational
prime produces $n$ parameters of modulus at most one,
padding with zero parameters at ramified primes when needed.
Indeed a factor $1-\eta(\mathfrak p)p^{-fs}$ factors into
$f$ linear factors in $p^{-s}$ whose roots have modulus one;
the degree formula gives total degree at most $n$.
The height shift multiplies those roots by $p^{-iT}$ and
preserves their moduli.  This proves the Ramanujan hypothesis
used in the corollary.  It also proves that the logarithmic
Euler coefficient at $p^k$ has modulus at most $n\log p$.
The source's proof of its explicit prime-power estimate uses
RH for $\zeta_{\mathbb Q}$, which is why that input has been
stated separately above.

In the source's normalization the conductor is
\[
 C_{\mathrm{Ch}}=
 \frac{Q_\eta}{\pi^n}
        \prod_{j=1}^n\left|\frac14+
                              \frac{\mu_j+iT}{2}\right|
 =\frac{Q_\eta}{(2\pi)^n}
        \prod_{j=1}^n|1/2+\mu_j+iT|
 \le\mathscr C_\eta(T).
\]
It is legitimate to use this larger conductor even if only
its double logarithm exceeds ten.  To see this directly in
the cited proof, retain the free parameter $x$ in equation
(14), replace $\log C_{\mathrm{Ch}}$ by
$\log\mathscr C_\eta(T)$ in its term with positive
coefficient $3/(4\log x)$, and then choose
$x=(\log\mathscr C_\eta(T))^2$.
The estimate (15) and the remaining terms depend only on
$x$ and $n$, and $\log x=2y\ge20$ meets their threshold.
The same collection of constants that proves Corollary~1.1
therefore gives exactly (NT-L13).  This explains the conductor
normalization and does not silently replace it by an
asymptotically equivalent one when the degree varies.
If the central value is zero, the nonlogarithmic bound is
immediate; otherwise the logarithmic argument applies.

Finally the two conditions
$y\ge3/(4\epsilon)$ and
$y^2\ge46n/(25\epsilon)$ make the two terms
$3/(8y)$ and $23n/(25y^2)$ at most $\epsilon/2$ each.
Exponentiation proves (NT-L13-threshold).
The varying-degree assertion follows from the same two
terms tending to zero.  This completes the deduction.

Pure norm characters have shifted completed zeta poles and
are not inserted into these pole-free axioms; their fixed-function
Lindel\"of consequences remain in NT-L03.  Geometrically,
(NT-L13) bounds the critical value of a rank-one unitary
arithmetic object uniformly in its finite and archimedean
conductors.  No algebraicity of an infinite-order infinity
type, negative-moment estimate, or nonzero lower bound at the
critical point is required or inferred.

\paragraph{NT-L14. Explicit two-sided bounds on the one-line.}
Let $\eta$ be a primitive unitary non-pure-norm Hecke character
over a degree-$n$ field, and set
\[
 C_1(\eta,T)=\frac{Q_\eta}{(2\pi)^n}
                   \prod_{j=1}^n|1+\mu_j+iT|.
\]
Here the exact normalization is retained.
For $\log C\ge23n$ put $Y=\log\log C-\log(2n)>0$ and define
\begin{align*}
 k_n&=2.31+\frac{22.59}{n}(e^{0.31n}-1-0.31n),\\
 j_n&=2+\frac{4.18}{n}(e^{0.69n}-1-0.69n),\\
 j_n'&=9+\frac{16.74}{n}(e^{0.69n}-1-0.69n),\\
 U_n(C)&=(2e^\gamma)^n
       \left(Y^n+\frac n2Y^{n-1}+\frac{nk_n}4Y^{n-2}\right),\\
 V_n(C)&=\left(\frac{12e^\gamma}{\pi^2}\right)^n
       \left(Y^n+\frac n2Y^{n-1}+\frac{nj_n}4Y^{n-2}
                         +\frac{n^2j_n'Y^n}{\log C}\right).
\end{align*}
Under GRH for $L(s,\eta)$ and RH for $\zeta_{\mathbb Q}$,
\[
  V_n(C_1(\eta,T))^{-1}
        \le |L(1+iT,\eta)|\le U_n(C_1(\eta,T))
        \qquad\text{if }\log C_1(\eta,T)\ge23n.
\tag{NT-L14}
\]
Thus $M$ supplies the hypotheses, while $A$ supplies them for
norm shifts of nonprincipal primitive Dirichlet characters.
For fixed $n$ the leading upper order is
$(2e^\gamma\log\log C)^n$, and the reciprocal leading order
is $(12e^\gamma\log\log C/\pi^2)^n$.
These are two-sided inequalities, not a claim that every value
has either extremal order.

This is Lumley2018, \emph{Explicit bounds for $L$-functions
on the edge of the critical strip}, Theorems~2.1--2.2,
author pp.~3--4, and Corollary~2.3, p.~5.
The definitions and axioms are in \S1.1, pp.~2--3;
the proofs are \S4.1, pp.~11--13, and \S4.2, pp.~14--16.
The exact constants (2.1)--(2.3) and both theorem formulas
were inspected on the original p.~4 image.
The rational-prime Euler-factor and analytic-axiom verification
is exactly the one given in NT-L13, with parameters shifted
by $iT$.  The source allows complex infinity parameters and
requires an entire completion, which explains the non-pure-norm
condition.  Its Lemmas~3.6--3.7, p.~10, explicitly use RH for
$\zeta_{\mathbb Q}$; that input is retained here.
Its quantity $J_1(n)$ is at most $j_n$, while $J_2(n)=j_n'$.
Since $Y>0$, replacing $J_1(n)$ by its displayed upper bound
preserves the reciprocal inequality.  This verifies every
specialization and normalization in (NT-L14).
The numerical constants remain the cited classical theorem;
no new numerical optimization is being claimed.

\paragraph{NT-H11. Relative class-number--regulator bounds over any base field.}
Let $E/K$ be a finite abelian extension with group $G$, and
write $n=[K:\mathbb Q]$.
For every $\chi\in\widehat G\setminus\{1\}$, let
$L(s,\chi)$ mean the primitive Hecke function attached to that
Artin character, including its correct ramified local factors.
Let $\mathfrak f_\chi$ be its finite conductor and let
$a_\chi$ count the real places of $K$ where its sign is odd.
Its conductor for NT-L14 at height zero is explicitly
\[
 C_\chi=
 \frac{\Delta_K\operatorname N\mathfrak f_\chi}
                  {(2\pi)^n}\,2^{a_\chi+r_2(K)}.
\tag{NT-relative-conductor}
\]
For a number field $F$, let $h_F$ be the ordinary ideal class
number, $R_F$ its full unit regulator, $w_F$ its number of roots
of unity, and $r_1(F),r_2(F)$ its signature.  Put
\[
 D(E/K)=\frac{w_E}{w_K}
       \sqrt{\frac{\Delta_E}{\Delta_K}}\,
       2^{r_1(K)-r_1(E)}(2\pi)^{r_2(K)-r_2(E)}.
\]
Assume RH for $\zeta_E$ and $\zeta_{\mathbb Q}$; in particular,
$M$ suffices.  If every nontrivial character satisfies
$\log C_\chi\ge23n$, then
\[
 D(E/K)\prod_{\chi\ne1}V_n(C_\chi)^{-1}
 \le\frac{h_ER_E}{h_KR_K}
 \le D(E/K)\prod_{\chi\ne1}U_n(C_\chi).
\tag{NT-H11}
\]
If some conductors fall below this threshold, let
$\mathcal S$ be precisely that finite set of characters and
put $F_{\mathrm{small}}=\prod_{\chi\in\mathcal S}|L(1,\chi)|$.
The valid bounds are then (NT-H11) with the products restricted
to $\chi\notin\mathcal S\cup\{1\}$ and both endpoints
multiplied by $F_{\mathrm{small}}$.
Thus no small-conductor value has been silently set to one.
A bound on $h_E/h_K$ alone additionally requires information
on $R_E/R_K$.

Here is the complete proof, including the ramified factors.
At a prime $\mathfrak p$ of $K$ put $q=\operatorname N\mathfrak p$.
Let $I\subset G$ be its inertia group, of order $e$, and let
$g\in G/I$ be Frobenius of order $f$.
There are $|G|/(ef)$ primes of $E$ above $\mathfrak p$,
all of norm $q^f$; the local Dedekind factor of $E$ is therefore
$(1-q^{-fs})^{-|G|/(ef)}$.
Only the characters trivial on $I$ contribute a local Hecke
factor, namely $(1-\chi(g)q^{-s})^{-1}$.
On restricting the characters of $G/I$ to its cyclic subgroup
$\langle g\rangle$, each of the $f$ roots of unity occurs
$|G|/(ef)$ times.  Consequently the product of these local
Hecke factors is the same Dedekind factor.
This proves on the absolutely convergent half-plane, and hence
by continuation,
\[
 \zeta_E(s)=\zeta_K(s)\prod_{\chi\ne1}L(s,\chi).
\tag{NT-relative-factorization}
\]
This uses the classical class-field identification of finite-order
Hecke characters with the characters of $G$, as in Milne V,
Theorems~3.5--3.6; the local factor computation above proves
the actual Euler-product identity used here.

All the nontrivial Hecke factors are entire, and the only pole
of $\zeta_K$ is at $1$.  A nontrivial zero off the critical
line of any factor would therefore be a nontrivial zero of
$\zeta_E$.  Thus RH for $\zeta_E$ supplies the Hecke GRH
hypotheses required by NT-L14.  The remaining RH for
$\zeta_{\mathbb Q}$ has been stated explicitly.
At each real place the gamma parameter of $\chi$ is $0$ or $1$;
at a complex place the two real-gamma parameters are $0,1$.
Substitution into the definition of $C_1$ proves
(NT-relative-conductor), with no asymptotic normalization.

Let $\kappa_F=\operatorname{Res}_{s=1}\zeta_F(s)>0$.
Taking residues in (NT-relative-factorization) and absolute
values gives
\[
 \frac{\kappa_E}{\kappa_K}
                  =\prod_{\chi\ne1}|L(1,\chi)|.
\]
The analytic class-number formula in Milne VI,
Corollary~2.12 and Example~2.13, printed p.~191/PDF p.~200,
is
\[
 \kappa_F=\frac{2^{r_1(F)}(2\pi)^{r_2(F)}h_FR_F}
                         {w_F\sqrt{\Delta_F}}.
\]
Its ratio therefore yields the exact identity
\[
 \frac{h_ER_E}{h_KR_K}
             =D(E/K)\prod_{\chi\ne1}|L(1,\chi)|.
\]
Apply NT-L14 to each factor above its threshold and leave the
remaining factors equal to their actual values.  This proves
both (NT-H11) and the small-conductor version.
Empty products give the correct identity when $E=K$.
The proof is complete.

For $E/\mathbb Q$ abelian, $A$ supplies RH for $\zeta_E$,
recovering and sharpening the characterwise framework of NT-H04.
For a general abelian extension $E/K$, the needed Dedekind
input is genuinely that of $E$, not just that of $K$.
The geometric comparison is between the Picard groups of the
two arithmetic curves and the rank-one character decomposition
of their finite abelian cover.  Its regulator and real-place
factors are integral parts of the comparison, not omitted
normalization terms.

\subsection*{NT-L15. A mean-square criterion for simple Hecke zeros and reciprocal derivatives}
Let $\eta$ be a fixed primitive unitary Hecke character of a number field $K$,
including a pure norm character. Define coefficients $b_\eta(\mathfrak a)$ by
\[
 L(s,\eta)^{-1}=\sum_{\mathfrak a\ne0}
       b_\eta(\mathfrak a)(N\mathfrak a)^{-s}
                    \qquad(\Re s>1),
 \qquad
 \mathcal M_\eta(x)=\sum_{N\mathfrak a\le x}b_\eta(\mathfrak a).
\]
Thus $b_\eta(\mathfrak a)=\mu_K(\mathfrak a)\eta(\mathfrak a)$
on ideals prime to the conductor and is zero on the other ideals.
Assume the additional weak mean-square hypothesis
\[
 B_\eta:=\limsup_{X\to\infty}\frac1{\log X}
       \int_1^X\left|\frac{\mathcal M_\eta(x)}x\right|^2dx<\infty.
                                                        \tag{NT-L15a}
\]
Then GRH holds for $L(s,\eta)$, every nontrivial zero is simple, and
\[
 \sum_{\rho=1/2+i\gamma}\frac1{|\rho L'(\rho,\eta)|^2}
                          \le B_\eta.                   \tag{NT-L15b}
\]
In particular $B_\eta>0$ and
\[
 |L'(\rho,\eta)|^{-1}\le\sqrt{B_\eta}\,|\rho|,
 \qquad |L'(\rho,\eta)|^{-1}=o(|\rho|),
                                                        \tag{NT-L15c}
\]
where the little-$o$ is along zeros of unbounded height. Also
\[
 \sum_{|\gamma|\le T}|L'(1/2+i\gamma,\eta)|^{-2}=o(T^2).
                                                        \tag{NT-L15d}
\]
No use of $A$ or $M$ is required once (NT-L15a) is assumed. Neither $A$ nor
$M$ is asserted to imply this additional mean-square hypothesis.
For $K=\mathbb Q$, $\eta=1$, the classical source is Titchmarsh,
\emph{The Theory of the Riemann Zeta-function}, 2nd ed.,
\S14.29: hypothesis (14.29.1) and Theorem14.29(A), pp.376--377;
reciprocal derivative estimate (14.29.3), p.377;
Theorem14.29(B)/(14.29.4), p.377, with proof ending p.380.
The following proof gives the Hecke formulation and its precise limsup constant.

\begin{proof}
Put $F(y)=e^{-y/2}\mathcal M_\eta(e^y)$ for $y\ge0$ and
$A(Y)=\int_0^Y|F(y)|^2dy$. Changing variables $x=e^y$ gives
$B_\eta=\limsup_{Y\to\infty}A(Y)/Y$.
For every $\epsilon>0$ there is $C_\epsilon$ such that
\[
 A(Y)\le (B_\eta+\epsilon)Y+C_\epsilon\qquad(Y\ge0).
                                                        \tag{NT-L15e}
\]
In particular the Laplace integral
\[
 H(w)=\int_0^\infty F(y)e^{-wy}\,dy
\]
is holomorphic for $\Re w>0$. To check this, integrate $|F|^2$ against
$e^{-\sigma y}$ by parts using (NT-L15e), then apply Cauchy--Schwarz to
$F(y)e^{-wy}$. The same argument with any fixed factor $y^j$ proves
uniform convergence of every derivative on each compact subset of that
half-plane. For $\Re w>1/2$, absolute convergence of the inverse Euler
series and partial summation give
\[
 H(w)=\int_1^\infty\mathcal M_\eta(x)x^{-(w+1/2)-1}\,dx
      =\frac1{(w+1/2)L(w+1/2,\eta)}.                    \tag{NT-L15f}
\]
The meromorphic identity theorem extends (NT-L15f) throughout $\Re w>0$.
Since $H$ is holomorphic there and $w+1/2\ne0$, $L(s,\eta)$ has no zeros
with $\Re s>1/2$. The conjugate character has coefficients
$b_{\bar\eta}(\mathfrak a)=\overline{b_\eta(\mathfrak a)}$, hence the
same mean-square hypothesis and the same zero-free half-plane.
The Hecke functional equation now excludes nontrivial zeros of
$L(s,\eta)$ with $\Re s<1/2$. This proves GRH, including non-real
characters; possible pure-norm poles only give zeros of $H$ and cause
no difficulty.

For $\sigma>0$ introduce the positive inner product
\[
 \langle f,g\rangle_\sigma
       =\sigma\int_0^\infty f(y)\overline{g(y)}e^{-\sigma y}dy.
\]
An integration by parts and (NT-L15e) show that
\[
 \|F\|_\sigma^2
   =\sigma^2\int_0^\infty A(y)e^{-\sigma y}dy,
 \qquad \limsup_{\sigma\downarrow0}\|F\|_\sigma^2\le B_\eta.
                                                        \tag{NT-L15g}
\]
Each $e_t(y)=e^{ity}$ has norm one. Consequently
\[
 |\sigma H(\sigma+it)|
       =|\langle F,e_t\rangle_\sigma|
       \le\|F\|_\sigma.
                                                        \tag{NT-L15h}
\]
At a zero $\rho=1/2+i\gamma$, (NT-L15f) has a pole of the same order
as that zero, since $\rho\ne0$. The bound (NT-L15h) as
$\sigma\downarrow0$ rules out order greater than one. Thus every zero
is simple and
\[
 \lim_{\sigma\downarrow0}\langle F,e_\gamma\rangle_\sigma
       =a_\gamma:=\frac1{\rho L'(\rho,\eta)}.
                                                        \tag{NT-L15i}
\]
For two distinct ordinates,
\[
 \langle e_\gamma,e_\nu\rangle_\sigma
       =\frac\sigma{\sigma-i(\gamma-\nu)}\longrightarrow0;
\]
the same inner product is one when $\gamma=\nu$.
Take an arbitrary finite set $S$ of zero ordinates and put
$P_S=\sum_{\gamma\in S}a_\gamma e_\gamma$.
Expanding the nonnegative number $\|F-P_S\|_\sigma^2$, using
(NT-L15g)--(NT-L15i), and taking the upper limit gives
\[
 0\le\limsup_{\sigma\downarrow0}\|F-P_S\|_\sigma^2
       \le B_\eta-\sum_{\gamma\in S}|a_\gamma|^2.
\]
The supremum over finite $S$ proves (NT-L15b). Nontrivial Hecke zeros
exist, by the classical zero-count asymptotic, so $B_\eta>0$.
Every summand is at most $B_\eta$, proving the first assertion of
(NT-L15c). Square summability makes $a_\gamma\to0$ along unbounded
ordinates and proves the second assertion. Finally, for $T\ge1$,
\[
 T^{-2}\sum_{|\gamma|\le T}|L'(\rho,\eta)|^{-2}
   =\sum_\gamma {\bf1}_{|\gamma|\le T}
           \frac{1/4+\gamma^2}{T^2}|a_\gamma|^2.
\]
Each factor before $|a_\gamma|^2$ tends to zero and is at most $5/4$.
Dominated convergence for the summable series proves (NT-L15d).
\end{proof}

The finite Bessel argument proves an inequality. It does not assert an
exact mean-square limit, completeness of an infinite Fourier expansion
of $F$, or a limiting distribution for $F$. Such assertions require
additional approximation estimates.
Geometrically, the coefficients are the inverse Euler weights of the
rank-one arithmetic object attached to $\eta$; squarefree ideals are
squarefree effective divisors on $\operatorname{Spec}\mathcal O_K$.
The new hypothesis is analytic control of these weighted divisors and
is not a consequence of a finite-field purity theorem established here.

\subsection*{NT-Z19. Weak mean square supplies the canonical arithmetic basis gate}
Use the canonical completed, pole-removed $\Xi_\eta$, inverse Mellin
function $\psi_\eta$, and closed quotient
\[
 Q_\eta=\mathcal E/\overline{\psi_\eta*\mathcal E},\qquad
 \mathcal E=\{f\in C^\infty(\mathbb R):
      \sup_x e^{N|x|}|f^{(j)}(x)|<\infty\text{ for all }N,j\ge0\}
\]
of MT-M12/MT-M16. Under (NT-L15a), all zeros are simple and
\[
 \overline{\psi_\eta*\mathcal E}
   =\{f:\mathcal Mf(\rho)=0\text{ for every nontrivial zero }\rho\},
 \qquad Q_\eta\simeq s(Z).
                                                        \tag{NT-Z19a}
\]
Here $\mathcal Mf(s)=\int_\mathbb Rf(x)e^{sx}dx$,
and $s(Z)$ is the space of sequences decreasing faster than every
power of $1+|\operatorname{Im}\rho|$. The quotient has an absolute
Schauder basis with these Mellin evaluations as coefficient functionals.
Convolution by every $h\in\mathcal E$ is nuclear. In the diagonal
Hilbert realization $\ell^2(Z)$ it is trace class, with Hilbert trace
\[
 \operatorname{tr}_{\ell^2(Z)}T_h=\sum_\rho\mathcal Mh(\rho).
\]
Writing $h^\#(x)=e^{-x}\overline{h(-x)}$, the operator
$T_{h*h^\#}$ is positive and
\[
 \operatorname{tr}_{\ell^2(Z)}T_{h*h^\#}
                     =\sum_\rho|\mathcal Mh(\rho)|^2\ge0.
                                                        \tag{NT-Z19b}
\]
\begin{proof}
NT-L15 proves GRH, simplicity, and
$|L'(\rho,\eta)|\ge B_\eta^{-1/2}|\rho|^{-1}$.
Since $|\rho|\le1+|\operatorname{Im}\rho|$, this is precisely the
polynomial derivative lower bound of MT-M12, or condition (1) of
MT-M16 with every multiplicity equal to one. The entire proofs of
those results give (NT-Z19a) and the absolute basis. MT-M13/MT-M17
then give nuclearity and the trace-class Hilbert realization. For
completeness, on each coordinate the convolution eigenvalue is
$\mathcal Mh(\rho)$, and direct substitution gives
$\mathcal M h^\#(s)=\overline{\mathcal Mh(1-\bar s)}$.
GRH implies $1-\bar\rho=\rho$, so the product eigenvalue is
$|\mathcal Mh(\rho)|^2$. Rapid Mellin decay and NT-Z10's local zero
count make both displayed trace series absolutely convergent. This
proves positivity and (NT-Z19b).
\end{proof}
This is a conditional analytic realization of the zero divisor with an
explicit arithmetic gate. It does not identify the quotient with an
independently constructed motivic or étale cohomology group.
For the Riemann zeta function the sole additional arithmetic condition
is the classical weak Mertens conjecture; $A$ is not needed separately.

\subsection*{NT-H12. Class-group torsion from the field's own Dedekind RH}
Fix integers $n\ge2$ and $\ell\ge2$. For every degree-$n$ number field
$K$ with absolute discriminant $D_K$ whose Dedekind zeta function satisfies
RH, and every $\epsilon>0$,
\[
 |\operatorname{Cl}_K[\ell]|
       \ll_{n,\ell,\epsilon}
 D_K^{\,1/2-1/(2\ell(n-1))+\epsilon}.                   \tag{NT-H12a}
\]
Only RH for $\zeta_K$ is needed. In particular $M$ supplies the hypothesis
conditionally, and adopted $A$ supplies it when $K/\mathbb Q$ is abelian.
No RH for a normal closure or Hilbert class field is required by the
argument below. The exponent is the classical Ellenberg--Venkatesh
bound, not a new class-group exponent. Our point is the exact
Dedekind-only input reduction.

The precise cited algebraic input is Ellenberg--Venkatesh,
\emph{Reflection principles and bounds for class group torsion},
author preprint, Lemma3 and its proof, p.6: if
$0<\delta<1/(2\ell(n-1))$ and there are $N$ prime ideals of $K$
of norm at most $D_K^\delta$, unramified over $\mathbb Q$ and not
extensions of prime ideals of any proper intermediate field, then
\[
 |\operatorname{Cl}_K[\ell]|
       \ll_{n,\ell,\epsilon,\delta}D_K^{1/2+\epsilon}/N.
                                                        \tag{NT-H12b}
\]
The source's Proposition1, p.7, applies this lemma using completely
split primes. We instead count degree-one prime ideals. The implicit
constant in (NT-H12a) is not assigned an unchecked numerical value.

\begin{proof}
Let $\pi_K(x)$ count all prime ideals of $K$ of norm at most $x$.
Apply NT-C02 to the trivial Galois extension $K/K$. Its conjugacy density
is one, there are no relative ramified primes to omit, and RH for
$\zeta_K$ gives, for $x\ge2$,
\[
 |\pi_K(x)-\operatorname{Li}_2(x)|\le B_K(x),
\]
where the completely explicit error is
\[
 B_K(x)=\sqrt x\left[
 \left(\frac1{2\pi}+\frac3{\log x}\right)\log D_K+
 \left(\frac{\log x}{8\pi}+\frac1{4\pi}+\frac6{\log x}\right)n
                  \right].                            \tag{NT-H12c}
\]
Every prime ideal of residue degree at least two above its rational
prime has underlying prime at most $\sqrt x$. There are at most $n$
prime ideals above each rational prime, so such ideals number at most
$n\pi(\sqrt x)\le n\sqrt x$. There are at most $n\omega(D_K)$
prime ideals above rational primes ramified in $K$. Therefore the
number $N_K^{(1),\mathrm{ur}}(x)$ of unramified degree-one prime ideals
satisfies the useful finite inequality
\[
 N_K^{(1),\mathrm{ur}}(x)
   \ge\operatorname{Li}_2(x)-B_K(x)-n\sqrt x-n\omega(D_K).
                                                        \tag{NT-H12d}
\]
Here $\omega$ counts distinct rational prime divisors.

These degree-one ideals meet the remaining condition in (NT-H12b).
Indeed, if a prime ideal $\mathfrak p$ of norm $p$ were
$\mathfrak q\mathcal O_K$ for a prime ideal of a proper subfield $F$,
then multiplicativity of the norm under ideal extension would give
\[
 p=N_K(\mathfrak p)=(N_F\mathfrak q)^{[K:F]}.
\]
The right side is a proper positive integral power of an integer
at least two, and cannot be the rational prime $p$.

Fix $\delta>0$ and put $x=D_K^\delta$. For large $D_K$, the bound
$\operatorname{Li}_2(x)\ge x/(2\log x)$ follows by integrating on
$[x/2,x]$. The error in (NT-H12c) is
$O_{n,\delta}(D_K^{\delta/2}\log D_K)$, and
$\omega(D_K)\le\log D_K/\log2$. Consequently
\[
 N_K^{(1),\mathrm{ur}}(D_K^\delta)
                  \gg_{n,\delta}D_K^\delta/\log D_K.   \tag{NT-H12e}
\]
Write $\alpha=1/(2\ell(n-1))$. It suffices to prove (NT-H12a) for
$0<\epsilon<\alpha$, since a result with any smaller positive epsilon
implies the assertion with a larger one. Choose
$\delta=\alpha-\epsilon/3$ and apply (NT-H12b) with its epsilon
parameter equal to $\epsilon/3$. Equations (NT-H12b) and (NT-H12e)
give
\[
 |\operatorname{Cl}_K[\ell]|
   \ll D_K^{1/2-\alpha+2\epsilon/3}\log D_K
   \ll_{n,\ell,\epsilon}D_K^{1/2-\alpha+\epsilon}.
\]
The bounded-discriminant cases are absorbed into the constant by
Hermite's finiteness theorem (Milne, \emph{Algebraic Number Theory}
v3.08, Theorem8.43, printed pp.151--152). This proves the claim.
\end{proof}
The geometry is the torsion subgroup of
$\operatorname{Pic}(\operatorname{Spec}\mathcal O_K)$.
The source's Arakelov argument incorporates all archimedean places;
our analytic step counts rational degree-one closed points without
requiring a completely split Frobenius class.

\subsection*{NT-H13. Bounds for unramified cyclic covers and Kummer torsors}
Retain $n\ge2$, $K$, and RH for $\zeta_K$ from NT-H12.
For a prime $\ell$, the number of $K$-isomorphism classes of cyclic
extensions $E/K$ of degree $\ell$, unramified at every finite place
and split at every real place, is exactly
\[
 c_\ell(K)=\frac{|\operatorname{Cl}_K[\ell]|-1}{\ell-1}
       \ll_{n,\ell,\epsilon}
       D_K^{\,1/2-1/(2\ell(n-1))+\epsilon}.             \tag{NT-H13a}
\]
For odd $\ell$ the real-place condition is automatic. For $\ell=2$
it is essential to this ordinary-class-group formula.

For any integer $\ell\ge2$, let $S$ be the set of primes of $K$
above divisors of $\ell$, let $U=\operatorname{Spec}\mathcal O_K[1/\ell]$,
and let $w_K$ be the number of roots of unity in $K$. The number of
isomorphism classes of $\mu_\ell$-torsors over $U$ is
\[
 \begin{split}
 |H^1_{\text{ét}}(U,\mu_\ell)|
  &=\ell^{r_1+r_2-1+|S|}\gcd(\ell,w_K)
                             |\operatorname{Pic}(U)[\ell]|\\
  &\le\ell^{r_1+r_2+|S|}|\operatorname{Cl}_K[\ell]|\\
  &\ll_{n,\ell,\epsilon}
                 D_K^{\,1/2-1/(2\ell(n-1))+\epsilon}.
 \end{split}                                           \tag{NT-H13b}
\]
The finite coefficient $\ell^{r_1+r_2+|S|}$ is at most
$\ell^{n(1+\omega(\ell))}$. These torsors need not be connected;
$\mu_\ell$ need not be a constant group scheme over $K$.

\begin{proof}
Global class field theory identifies the ordinary Hilbert class field
Galois group with $\operatorname{Cl}_K$, with real places split.
The convention is stated explicitly in Milne, \emph{Algebraic Number
Theory} v3.08, Remark4.11 and its footnote, printed p.72; the general
reciprocity/existence input is Milne, \emph{Class Field Theory}, V3.5--3.6.
Degree-$\ell$ cyclic subextensions correspond to kernels of surjective
homomorphisms $\operatorname{Cl}_K\to\mathbb Z/\ell\mathbb Z$.
For prime $\ell$ every nonzero homomorphism is surjective, and
\[
 |\operatorname{Hom}(\operatorname{Cl}_K,\mathbb Z/\ell\mathbb Z)|
      =|\operatorname{Cl}_K[\ell]|.
\]
Exactly $\ell-1$ such homomorphisms have each kernel, by postcomposition
with automorphisms of the target. This proves the equality in
(NT-H13a), and NT-H12 gives the bound. A real-place inertia group has
order at most two, so it is trivial in an extension of odd prime degree.

Put $E_{K,S}=\mathcal O_K[1/\ell]^\times$. Since $\ell$ is invertible
on $U$, the étale Kummer sequence and
$H^1_{\text{ét}}(U,\mathbb G_m)=\operatorname{Pic}(U)$ give
\[
 0\longrightarrow
 E_{K,S}/E_{K,S}^{\ell}
 \longrightarrow H^1_{\text{ét}}(U,\mu_\ell)
 \longrightarrow\operatorname{Pic}(U)[\ell]
 \longrightarrow0.                                    \tag{NT-H13c}
\]
A precise source is Stacks Project, Section59.28, Tag03PK,
Lemma59.28.1 and its following cohomology sequence; Lemma59.28.5,
Tag040Q, identifies the classes as line bundles with an $\ell$th-power
trivialization. Thus this cohomology group is not being confused with
Picard torsion alone.

The $S$-unit theorem gives
$\mathcal O_K[1/\ell]^\times\simeq\mu(K)\times
\mathbb Z^{r_1+r_2-1+|S|}$ as abelian groups. Here $\mu(K)$ is cyclic
of order $w_K$. See Milne, \emph{Algebraic Number Theory} v3.08,
Theorem5.11 and its complete proof, printed p.90. Taking the quotient
by $\ell$th powers has cardinality
$\ell^{r_1+r_2-1+|S|}\gcd(\ell,w_K)$, proving the first equality in
(NT-H13b).

Localization makes $\operatorname{Pic}(U)$ the quotient of
$\operatorname{Cl}_K$ by the classes of primes in $S$. For any
surjection $A\to B$ of finite abelian groups, the map
$A/\ell A\to B/\ell B$ is surjective and
$|A[\ell]|=|A/\ell A|$, $|B[\ell]|=|B/\ell B|$.
It follows that $|\operatorname{Pic}(U)[\ell]|
\le|\operatorname{Cl}_K[\ell]|$; no injection on torsion induced by
localization is asserted. Apply NT-H12 and
$\gcd(\ell,w_K)\le\ell$, $r_1+r_2\le n$,
$|S|\le n\omega(\ell)$ to finish the proof.
\end{proof}
This is a literal finite étale/Kummer interpretation of a conditional
class-group bound: cyclic quotients of the arithmetic fundamental
group with the stated real conditions, and torsors under the finite
étale group scheme $\mu_\ell$ over $U$. It gives no bound for arbitrary
nonabelian covers or for all Selmer groups without further reduction.

\subsection*{NT-L16. Zero multiplicity forces logarithmic mean-square growth}
For a fixed primitive unitary Hecke function, use the inverse coefficients
and summatory function $\mathcal M_\eta$ of NT-L15, and put
\[
 I_\eta(X)=\int_1^X|\mathcal M_\eta(x)/x|^2dx.
\]
If $\rho=1/2+i\gamma$ is a zero of multiplicity $m\ge1$, write
$c_\rho=L^{(m)}(\rho,\eta)/m!$. Then, without assuming GRH,
\[
 \limsup_{X\to\infty}\frac{I_\eta(X)}{(\log X)^{2m-1}}
       \ge\frac1{(2m-1)((m-1)!)^2\,|\rho c_\rho|^2}>0.
                                                        \tag{NT-L16a}
\]
More strongly, the reciprocal leading-coefficient terms may be summed
over \emph{all distinct critical-line zeros of exact multiplicity $m$}:
\[
 \limsup_{X\to\infty}\frac{I_\eta(X)}{(\log X)^{2m-1}}
 \ge\frac1{(2m-1)((m-1)!)^2}
   \sum_{\substack{\Re\rho=1/2\\ \operatorname{ord}_\rho L=m}}
            \frac1{|\rho c_\rho|^2}.                    \tag{NT-L16a'}
\]
This includes the possibility that the sum and the limsup are infinite.
Consequently, for every $0\le a<2m-1$,
\[
 \limsup_{X\to\infty}I_\eta(X)/(\log X)^a=+\infty.
                                                        \tag{NT-L16b}
\]
If instead there is a zero with real part greater than $1/2$, then
(NT-L16b) holds for every finite real $a\ge0$.
The critical-line location of the zero in (NT-L16a) is an explicit
hypothesis; $M$ supplies it for every nontrivial zero, but is not
used in the proof of this implication.

\begin{proof}
Set $F(y)=e^{-y/2}\mathcal M_\eta(e^y)$ and
$A(Y)=\int_0^Y|F(y)|^2dy=I_\eta(e^Y)$.
Let $b=2m-1$ and
$B=\limsup_{Y\to\infty}A(Y)/Y^b$.
If $B=+\infty$, both asserted lower bounds are immediate.
Otherwise, for every $\epsilon>0$,
$A(Y)\le(B+\epsilon)Y^b+C_\epsilon$.
As in NT-L15, Cauchy--Schwarz with an exponential weight gives a
holomorphic Laplace transform on $\Re w>0$, and partial summation
in its initial absolute-convergence half-plane gives
\[
 H(w)=\int_0^\infty F(y)e^{-wy}dy
      =\frac1{(w+1/2)L(w+1/2,\eta)}.
                                                        \tag{NT-L16c}
\]
The identity extends meromorphically through $\Re w>0$.
Use again the inner product
\[
 \langle f,g\rangle_\sigma=\sigma\int_0^\infty
       f(y)\overline{g(y)}e^{-\sigma y}dy.
\]
Integration by parts yields
\[
 \begin{split}
 \|F\|_\sigma^2
 &=\sigma^2\int_0^\infty A(y)e^{-\sigma y}dy\\
 &\le (B+\epsilon)\Gamma(b+1)\sigma^{1-b}
                               +C_\epsilon\sigma.
 \end{split}                                           \tag{NT-L16d}
\]
Put $k=m-1$ and use the normalized polynomial waves
\[
 u_{\gamma,\sigma}(y)=
       \frac{\sigma^k y^k e^{i\gamma y}}{\sqrt{(2k)!}},
 \qquad v_\sigma(y)=\sigma^k F(y).
\]
Their Gram matrix has entries
\[
 \langle u_{\gamma,\sigma},u_{\nu,\sigma}\rangle_\sigma
       =\left(\frac{\sigma}{\sigma-i(\gamma-\nu)}\right)^{2k+1}.
                                                        \tag{NT-L16e}
\]
In particular each vector has norm one, and for any fixed finite set
of distinct ordinates its Gram matrix tends to the identity as
$\sigma\downarrow0$. Polynomial growth of $A$ permits differentiation
under the Laplace integral on every compact subset of $\Re w>0$,
by weighted Cauchy--Schwarz. Hence
\[
 \langle v_\sigma,u_{\gamma,\sigma}\rangle_\sigma
 =\frac{\sigma^{2k+1}}{\sqrt{(2k)!}}
                 (-1)^k H^{(k)}(\sigma+i\gamma).
\]
At a zero $\rho=1/2+i\gamma$ of exact order $m=k+1$,
the meromorphic identity (NT-L16c) gives
$H(w)=(\rho c_\rho)^{-1}(w-i\gamma)^{-(k+1)}$
plus terms of smaller pole order. Differentiating this finite
principal part yields
\[
 \lim_{\sigma\downarrow0}
 \langle v_\sigma,u_{\gamma,\sigma}\rangle_\sigma
       =\frac{\sqrt{(2k)!}}{k!\,\rho c_\rho}.           \tag{NT-L16f}
\]
For a fixed finite set $S$ of the specified zeros, finite-dimensional
orthogonal projection onto the span of these waves gives the Gram
matrix version of Bessel's inequality. By (NT-L16e), letting
$\sigma\downarrow0$ in this inequality and using (NT-L16d) gives
\[
 \frac{(2k)!}{(k!)^2}\sum_{\rho\in S}|\rho c_\rho|^{-2}
   \le \limsup_{\sigma\downarrow0}\|v_\sigma\|_\sigma^2
   \le (B+\epsilon)(2k+1)!.
\]
Let $\epsilon\downarrow0$ and take the supremum over finite $S$.
Since the terms are nonnegative, this proves (NT-L16a') and
therefore (NT-L16a). For $m=1$ it is exactly the Abel--Bessel
normalization of NT-L15. The one-frequency model
$F(y)=a y^{m-1}e^{i\gamma y}/(m-1)!$ has
$A(Y)=|a|^2Y^{2m-1}/((2m-1)((m-1)!)^2)$, so this constant is
sharp for the single-pole Laplace argument; no attainment by an
arithmetic Hecke function is asserted.
The positive limsup in (NT-L16a) gives (NT-L16b) along an unbounded sequence
when $a<2m-1$.

Finally suppose a zero has real part greater than $1/2$.
A finite limsup in (NT-L16b), for any fixed $a\ge0$, would imply
$A(Y)\le C(1+Y^a)$. The same weighted Cauchy--Schwarz argument
makes $H$ holomorphic on $\Re w>0$. Equation (NT-L16c) would then
exclude every zero of $L$ in $\Re s>1/2$, a contradiction.
This proves the final assertion.
\end{proof}
This is a direct multiplicity-sensitive version of the weak-Mertens
Laplace argument. It is a necessary analytic growth condition, not
an asymptotic formula for $I_\eta(X)$.

\subsection*{NT-L17. Actual Kummer fields obstruct universal weak Mertens and full-jet positivity}
Let $r\ge3$ be an odd prime and let $p\ne r$ be a rational prime. Put
\[
 K=\mathbb Q(\zeta_r),\qquad E=\mathbb Q(\zeta_r,p^{1/r}).
\]
There is a nontrivial primitive finite-order Hecke character $\theta$
of $K$ whose entire finite-part $L$-function $H(s)=L_K(s,\theta)$
satisfies the all-prime factorization
\[
                  \zeta_E(s)=\zeta_K(s)H(s)^{r-1}.     \tag{NT-L17a}
\]
The function $H$ has infinitely many nontrivial zeros. Hence $\zeta_E$
has infinitely many nontrivial zeros of multiplicity at least $r-1$,
\emph{unconditionally}.

Let $\mathcal M_E(x)=\sum_{N\mathfrak a\le x}\mu_E(\mathfrak a)$
and $I_E(X)=\int_1^X|\mathcal M_E(x)/x|^2dx$. Then
\[
            \limsup_{X\to\infty}I_E(X)/\log X=+\infty.
                                                        \tag{NT-L17b}
\]
Thus a universal weak-Mertens hypothesis for all Dedekind zeta
functions is false. This statement makes no assertion about an
individual Dirichlet function's weak-Mertens conjecture.

If RH for $\zeta_E$ is additionally assumed, choose any zero $\rho$
of $H$. It is then on the critical line. With
$m=\operatorname{ord}_\rho\zeta_E\ge r-1$, NT-L16 gives the explicit
nonzero lower bound
\[
 \limsup_{X\to\infty}\frac{I_E(X)}{(\log X)^{2m-1}}
 \ge
 \frac1{(2m-1)((m-1)!)^2\,
             |\rho\,\zeta_E^{(m)}(\rho)/m!|^2}.
                                                        \tag{NT-L17c}
\]
In particular $\limsup I_E(X)/(\log X)^{2r-3}>0$; it is infinite
if the chosen multiplicity exceeds $r-1$. If instead $\zeta_E$
has an off-critical zero, its functional equation supplies a
right-half-strip zero, and every polynomial-logarithmic mean bound
fails by NT-L16. This gives the same positive-limsup conclusion
unconditionally by a dichotomy, while (NT-L17c) retains its stated
critical-zero/RH hypothesis.

There is also a concrete consequence for MT-M20: the canonical
\emph{full-jet} quotient for $\zeta_E$ has no positive-definite
continuous inner product for which every convolution extends boundedly
and every $C_{h*h^\sharp}$ is positive. Under $M$, the Hilbert traces
in MT-M19 are nevertheless nonnegative. Trace positivity and a faithful
positive operator realization are different assertions here.

\begin{proof}
The elementary Galois calculation in CX-ALG-024 gives
$[E:K]=r$ and
$\operatorname{Gal}(E/\mathbb Q)=\mathbb F_r\rtimes\mathbb F_r^\times$.
For clarity, $p$ is unramified in the cyclotomic field $K$, and
$X^r-p$ is Eisenstein at every prime above $p$, proving the degree.
The translations send $p^{1/r}$ to $\zeta_r^a p^{1/r}$.
The automorphism fixing $p^{1/r}$ and sending $\zeta_r$ to
$\zeta_r^b$ conjugates a translation by $a$ to one by $ba$.
Thus the $r-1$ nontrivial characters of
$\operatorname{Gal}(E/K)=\mathbb F_r$ form a single orbit under
$\operatorname{Gal}(K/\mathbb Q)$.

Class field theory identifies these characters with primitive
finite-order Hecke characters $\theta^a$ of $K$.
Conjugating ideals by an automorphism of $K/\mathbb Q$ preserves
absolute norms and carries Frobenius and inertia by conjugation.
Consequently it permutes every finite Euler factor, including the
ramified factors, and all $L_K(s,\theta^a)$ with $a\ne0$ are the
same function $H(s)$. The unconditional all-prime abelian factorization
proved within NT-H11 therefore gives (NT-L17a).
Equivalently, CX-ALG-024's permutation-character identity gives
$H=\zeta_{\mathbb Q(p^{1/r})}/\zeta_{\mathbb Q}$.
No general Artin-holomorphy conjecture is being invoked: $H$ is an
actual nontrivial Hecke $L$-function, so it is entire by the classical
Hecke theorem (Ramakrishnan--Valenza, Theorem7-19, printed p.279).

Greni\'e--Molteni, AppendixA TheoremA.1, author p.26, applies
unconditionally to primitive finite-order Hecke characters. Its
zero count has main term
$([K:\mathbb Q]/\pi)T\log T+O_{K,\theta}(T)$.
Hence $H$ has infinitely many nontrivial zeros in $0<\Re s<1$.
At each such zero $\zeta_K$ is holomorphic, because its only pole
is at $1$. Equation (NT-L17a) therefore gives multiplicity at
least $r-1$ in $\zeta_E$.
If the limsup in (NT-L17b) were finite, NT-L15 applied to the
trivial primitive Hecke character over $E$ would make all these
zeros simple, a contradiction. This proves (NT-L17b).
The quantitative claims follow exactly from NT-L16 and the
functional equation of $\zeta_E$, as described above.
Finally MT-M20 proves that a faithful positive completion with the
stated convolution properties requires simple zeros. The actual
multiple zeros just proved exclude it. MT-M19 supplies the separate
trace-positive conclusion under RH for $\zeta_E$, which $M$ supplies
conditionally.
\end{proof}
For example $r=3$, $p=2$ gives the degree-six $S_3$ field
$E=\mathbb Q(\sqrt[3]{2},\sqrt{-3})$ and at least double nontrivial
zeros. The argument given here concerns these explicitly monomial
Galois fields; this proof alone does not address every nonabelian
Galois field.
Geometrically this is the positive multiplicity of a single
rank-one Hecke realization in the zeta realization of an affine
finite Galois cover. It is compatible with GRH and with the
unconditional adapted-block synthesis of MT-M18, but prevents
a faithful positive Hilbert realization of the full jet algebra,
even when its traces are positive under $M$.


\subsection*{NT-E01. Cyclic reductions from precisely the division-field zeta functions}
Let $E/K$ be an elliptic curve over any number field, with or without
complex multiplication, and put $K_m=K(E[m])$ in a fixed algebraic closure.
Write $S_{E/K}(x)$ for the number of good prime ideals $\mathfrak p$ of
$K$ with $N\mathfrak p\le x$ for which $E_{\mathfrak p}(k_{\mathfrak p})$
is cyclic. The classical conditional density theorem gives
\[
 S_{E/K}(x)=\delta_{E/K}\frac{x}{\log x}
                 +o_{E,K}\!\left(\frac{x}{\log x}\right),\qquad
 \delta_{E/K}=\sum_{m\ge1}\frac{\mu(m)}{[K_m:K]}\ge0.       \tag{NT-E01a}
\]
The series is absolutely convergent. A sufficient analytic input is
\emph{RH for $\zeta_{K_m}$ for every squarefree $m$, including $m=1$}.
Thus $M$ supplies this consequence conditionally. Neither GRH for
$L(s,E)$ nor Artin holomorphy for the division-field Galois groups is an
additional hypothesis. No assertion of a positive density is implicit
in (NT-E01a).

The exact cited locators are Murty, \emph{On Artin's conjecture}
(1983), Theorem1, printed p.153, with its proof on pp.153--155;
and Campagna--Stevenhagen, \emph{Cyclic reduction densities for
elliptic curves} (2023), \S2, Lemma2.1, Corollary2.2 and
Propositions2.3--2.4 (author v2 pp.6--9). The latter verifies the
geometric degree growth, root-discriminant bound and large-prime
tail needed by Murty's theorem for arbitrary $E/K$.

\paragraph{Exact input identification.}
For a squarefree $m$, the compositum of the fields $K_\ell$ for
$\ell\mid m$ is $K_m$: the Chinese remainder decomposition of
$E[m]$ gives both inclusions. These are therefore exactly the fields
called $L_k$ in Murty's Theorem1. Its stated RH assumptions concern
their Dedekind zeta functions, rather than unspecified Artin factors.
The geometric criterion is that, at a good prime of characteristic
different from $\ell$, the full rank-two $\ell$-torsion is rational
over the residue field precisely when Frobenius acts trivially on
$E[\ell]$. Such full torsion prevents cyclicity. The characteristic
$p$ geometric $p$-torsion has rank at most one, so contributes no
additional rank-two obstruction. Corollary2.2 removes the resulting
finite set of residue-characteristic and model exceptions. This
identifies the inclusion--exclusion in (NT-E01a) with the stated
prime count. The limiting-density argument itself is the cited
classical theorem, not a new proof here.

Geometrically, each $K_m/K$ is the finite monodromy cover trivializing
the locally constant rank-two torsion sheaf of the elliptic scheme
over the appropriate arithmetic open. This uses finite quotients of
$\pi_1^{\text{ét}}$, and does not assert the Riemann hypothesis for
the infinite-image rank-two elliptic $L$-function. Work/run:
\texttt{GRH\_NIGHT/number\_theory}.

\subsection*{NT-E02. A finite Frobenius certificate and density-zero obstructions}
Suppose $E/K$ is non-CM. Let $T$ consist of the rational primes
dividing $30D_K$, those below a prime of bad reduction, and those for
which the image on $E[\ell]$ is not $\mathrm{GL}_2(\mathbb F_\ell)$.
This is a finite set, by the open image theorem. Put
$N=\prod_{\ell\in T}\ell$, $G_N=\operatorname{Gal}(K_N/K)$, and
\[
 C_N=\{g\in G_N:g|_{K_\ell}\ne1\text{ for every }\ell\mid N\}.
\]
Campagna--Stevenhagen Theorem1.1, author p.3, gives
\[
 \delta_{E/K}=\frac{|C_N|}{|G_N|}
   \prod_{\ell\nmid N}
       \left(1-\frac1{(\ell^2-1)(\ell^2-\ell)}\right).       \tag{NT-E02a}
\]
The product is strictly positive. Consequently $C_N=\varnothing$
is an unconditional finite certificate that all but finitely many
good reductions are noncyclic. If $C_N\ne\varnothing$, the hypothesis
of NT-E01 implies a positive density of cyclic reductions.

\paragraph{Complete finite-group deduction.}
For $m\mid N$, the event that a uniform $g\in G_N$ is trivial on
$K_m$ has probability $[K_m:K]^{-1}$. Intersections of these events
for primes dividing $m$ equal the $K_m$ event, by the compositum
identity proved above. Finite inclusion--exclusion therefore yields
\[
 \frac{|C_N|}{|G_N|}=\sum_{m\mid N}\frac{\mu(m)}{[K_m:K]}.
\]
The cited factorization gives (NT-E02a). For $\ell\ge2$, the
reciprocal of $(\ell^2-1)(\ell^2-\ell)$ is at most a constant
times $\ell^{-4}$ and is strictly less than one. The convergence
of its sum implies positivity of the product, by applying the
convergent logarithm series to all sufficiently large factors and
retaining the finitely many others.

If $C_N$ is empty, every element of $G_N$ lies in the kernel of at
least one restriction to $K_\ell$, $\ell\mid N$. Every good prime
outside the finite primes above $N$ and those ramified in $K_N$ has
such a Frobenius. Its reduction then contains
$(\mathbb Z/\ell\mathbb Z)^2$ for one such $\ell$, and is noncyclic.
This proves the finite-obstruction assertion without RH. Conversely,
$C_N\ne\varnothing$ makes the density expression positive, so NT-E01
gives infinitely many cyclic primes under its exact RH input.

For any $E/K$, including CM curves, the simpler obstruction
$E[\ell]\subset E(K)$ prevents cyclic reduction at every good prime
away from $\ell$, by the same reduction argument. Full rational
two-torsion is thus an immediate obstruction. Over a general $K$
the absence of this particular obstruction is not a positivity
criterion: the finite kernels in (NT-E02a) may jointly cover $G_N$
without any one being all of $G_N$. Actual elliptic examples of this
phenomenon are supplied by Campagna--Stevenhagen, Theorem5.1 and
\S5 (author pp.14--16); their construction uses three quadratic
division fields with Klein-four compositum. No independence of
division fields is silently assumed.

This is a finite monodromy-cover certificate, not a criterion for
a prescribed rational point to generate every cyclic reduction.
It also makes no claim about prime values of $\#E(\mathbb F_p)$.
Work/run: \texttt{GRH\_NIGHT/number\_theory}.

\subsection*{NT-E03. Lang--Trotter upper bounds without an extra holomorphy conjecture}
For a fixed non-CM elliptic curve $E/\mathbb Q$, set
$a_p=p+1-\#E(\mathbb F_p)$ at good primes, and
$P_{E,a}(x)=\#\{p\le x:p\nmid N_E,\ a_p=a\}$.
Zywina, \emph{Bounds for the Lang--Trotter conjectures},
Theorem1.2, author p.1, proves
\[
 P_{E,a}(x)\ll_E\frac{x^{4/5}}{(\log x)^{3/5}},\qquad
 P_{E,0}(x)\ll_E\frac{x^{3/4}}{(\log x)^{1/2}}.            \tag{NT-E03a}
\]
The first constant is uniform over integers $a$. These are upper
bounds, not the Lang--Trotter asymptotic. The source defines its GRH
on author p.3 as Dedekind RH for number fields. The proof permits the
following sufficient smaller input: adopted $A$, together with RH
for $\zeta_{\mathbb Q(E[\ell])}$ for every sufficiently large prime
$\ell$. Hence $M$ supplies the conclusion. Non-CM $\mathrm{GL}_2$
GRH is not being inferred.

\paragraph{Complete holomorphy/input audit.}
For large $\ell$ the mod-$\ell$ image is $G=\mathrm{GL}_2(\mathbb F_\ell)$.
Write $L=\mathbb Q(E[\ell])$, let $B$ be its upper triangular
subgroup, and let $U$ be its unipotent subgroup. Zywina's complete
\S4 proof, author pp.10--11, applies its Theorem2.3 to
$L^U/L^B$. Its group $B/U$ is abelian, so its Artin functions are
Hecke functions and the holomorphy assumption of Theorem2.3 is
known. Moreover $L/L^U$ has the cyclic group $U$, so the all-prime
abelian factorization NT-H11 gives
$\zeta_L=\zeta_{L^U}\prod_{\chi\ne1}L_{L^U}(s,\chi)$ with every
other factor entire. Thus RH for $\zeta_L$ implies RH for
$\zeta_{L^U}$, exactly what Theorem2.3 requires.

For $a=0$ the proof replaces $U$ by the subgroup $H$ of triangular
matrices with equal diagonal entries. This group is abelian:
writing a matrix as $aI+b\left(\begin{smallmatrix}0&1\\0&0\end{smallmatrix}\right)$
makes commutativity explicit. The same factorization transfers RH
from $\zeta_L$ to $\zeta_{L^H}$, while $B/H$ is again abelian.
The separate step choosing a prime $\ell$ split in
$\mathbb Q(\sqrt{a_p^2-4p})$ uses only RH for quadratic Dedekind
zeta functions and for $\zeta_{\mathbb Q}$: see Murty--Murty--Saradha
(1988), Lemma4.4 and its entire proof, printed pp.269--270.
These are supplied by $A$. This proves the stated sufficient input
reduction. The quantitative estimates in (NT-E03a) remain the
cited classical theorem.

The geometric data used here are finite mod-$\ell$ monodromy of
$H^1_{\text{ét}}(E_{\overline{\mathbb Q}},\mathbb F_\ell)$, with
the usual dual convention immaterial to the division field, and
abelian covers of intermediate arithmetic bases. An arbitrary
nonabelian use of Zywina's Theorem2.3 would still require its stated
Artin holomorphy hypothesis. Work/run:
\texttt{GRH\_NIGHT/number\_theory}.

\subsection*{NT-E04. Uniform counts with a prescribed Frobenius field}
Let $E/\mathbb Q$ be non-CM and let $k$ be an imaginary quadratic
field with class number $h_k$ and absolute discriminant $d_k$.
Put
\[
 P_{E,k}(x)=\#\{p\le x:p\nmid N_E,\
            \mathbb Q(\sqrt{a_p^2-4p})\cong k\}.
\]
Zywina Theorem1.3(i), author p.2, gives under Dedekind GRH
\[
 P_{E,k}(x)\ll_E
   h_k^{-3/5}\frac{x^{4/5}}{(\log x)^{3/5}}
                   +x^{1/2}(\log x)^3,                    \tag{NT-E04a}
\]
with a constant independent of $k$. Its proof is \S5.1,
author pp.13--14, after the algebraic reductions on pp.11--13.
Lemma5.5, author p.13, also gives $P_{E,k}(x)=0$ when $d_k>4x$.

For precision, this field family is larger than the division
fields in NT-E03. Let $H_k$ be the Hilbert class field of $k$,
and for split primes $\ell\ge5$ let $L_{k,\ell}/H_k$ be the
field of the joint representation formed from $E[\ell]$ and
the modulus-$\ell$ ray-class representation
$\psi_{k,\ell}:G_{H_k}\to
(\mathcal O_k/\ell\mathcal O_k)^*/\mathcal O_k^*$.
Zywina \S3.2--3.3, pp.7--10, defines it. The proof uses RH for
the Dedekind zeta functions of the auxiliary fields
$L_{k,\ell}^{\mathcal U}$ appearing in \S5, together with RH
for $\zeta_k$ and $\zeta_{\mathbb Q}$ in the split-prime choice.
Here $\mathcal B$ is the preimage of the triangular subgroup and
$\mathcal U$ is the image of the pairs $(aI+bN,a)$; the quotient
$\mathcal B/\mathcal U$ is abelian, so holomorphy is known.
This is an exact sufficient family specified by the cited proof;
$M$ supplies all of it. We do not replace it by Dirichlet GRH alone.

These counts distinguish endomorphism fields of reductions using
finite torsion and ray-class monodromy. They do not establish the
conjectured $x^{1/2}/\log x$ asymptotic. Work/run:
\texttt{GRH\_NIGHT/number\_theory}.

\subsection*{NT-E05. A refined lower bound for distinct Frobenius fields}
Let $\mathcal D_E(x)$ be the set of imaginary quadratic fields
arising as $\mathbb Q(\sqrt{a_p^2-4p})$ at good primes $p\le x$
for a fixed non-CM $E/\mathbb Q$. Under the exact input of NT-E04
for every imaginary quadratic $k$, together with adopted $A$,
\[
 |\mathcal D_E(x)|\gg_E
 \frac{x^{2/7}}
      {(\log x)^{4/7}(\log\log x)^{6/7}}
 \qquad(x\longrightarrow\infty).                         \tag{NT-E05a}
\]
In particular $M$ supplies these inputs. This is a newly proved
formulation in this run, not a claim of literature novelty.
Zywina's Corollary1.4, author p.2 with proof p.15, gives
$x^{2/7}/(\log x)^{10/7}$ using a weaker class-number estimate;
the proof below retains the sharper checked NT-H02 input.

\paragraph{Complete proof.}
Put $\Lambda(t)=\log\log(e^e t)$ for $t\ge1$. The explicit
NT-H02 bound, together with its finitely many excluded fundamental
discriminants, gives an absolute $c>0$ such that
\[
 h_k\ge c\,\frac{\sqrt{d_k}}{\Lambda(d_k)}
 \quad\hbox{for every imaginary quadratic }k.             \tag{NT-E05b}
\]
Indeed the displayed $V(d)$ in NT-H01 is $O(\log\log d)$ as
$d\to\infty$; lowering $c$ handles the finitely many smaller
discriminants, all of whose class numbers are positive.

Let $D=|\mathcal D_E(x)|$ and take $x$ sufficiently large that
$D\ge1$. Every such field has $d_k\le4x$ by NT-E04. Distinct
imaginary quadratic fields have distinct positive absolute
fundamental discriminants, so, ordering these discriminants,
the $j$th is at least $j$. It follows that
\begin{align*}
 \sum_{k\in\mathcal D_E(x)}h_k^{-3/5}
 &\ll \Lambda(4x)^{3/5}
                \sum_{k\in\mathcal D_E(x)}d_k^{-3/10}\\
 &\ll \Lambda(x)^{3/5}\sum_{j=1}^D j^{-3/10}
 \ll \Lambda(x)^{3/5}D^{7/10}.
\end{align*}
The last estimate uses
$\sum_{j=1}^D j^{-3/10}\le(10/7)D^{7/10}$, directly by integration of
$t^{-3/10}$ on $[1,D]$; its implicit constant is independent of
$E$, $x$ and the chosen fields.

Every good prime contributes to exactly one of these fields.
The prime number theorem, removal of the finitely many primes
dividing $N_E$, and (NT-E04a) therefore give
\[
 \frac{x}{\log x}\ll_E
    D^{7/10}\frac{x^{4/5}\Lambda(x)^{3/5}}{(\log x)^{3/5}}
                      +D x^{1/2}(\log x)^3.               \tag{NT-E05c}
\]
At least one right-hand term is at least a fixed positive
multiple of $x/\log x$. If it is the first, solving the
inequality gives
\[
 D\gg_E x^{2/7}(\log x)^{-4/7}\Lambda(x)^{-6/7}.
\]
If it is the second, then $D\gg_E x^{1/2}/(\log x)^4$.
The ratio of this latter lower bound to the former tends to
infinity, since it is
\[
 x^{3/14}(\log x)^{-24/7}\Lambda(x)^{6/7}\longrightarrow\infty.
\]
Thus the first lower bound follows in both cases for all
sufficiently large $x$. Finally
$\Lambda(x)\sim\log\log x$, which proves (NT-E05a).

This measures the diversity of Frobenius endomorphism fields
among closed fibres of an elliptic scheme. It does not prove
prime-order reductions or the Koblitz conjecture. Work/run:
\texttt{GRH\_NIGHT/number\_theory}.

\subsection*{NT-R05--NT-R07. Frobenius error laws with only Dedekind RH}

Fix a finite Galois extension $L/K$ of number fields, with group $G$.
For a real class function $t$ on $G$, put
\[
 \delta_t=\frac1{|G|}\sum_{g\in G}t(g),\qquad
 \pi_t(x)=\sum_{\substack{N\mathfrak p\le x\\\mathfrak p\ \text{unramified in }L}}
 t(\operatorname{Frob}_{\mathfrak p}),\qquad
 E_t(y)=ye^{-y/2}\bigl(\pi_t(e^y)-\delta_t\operatorname{Li}_2(e^y)\bigr).
\]
All assertions below require only RH for $\zeta_L$; hence $M$ suffices,
and $A$ suffices when $L/\mathbb Q$ is abelian. No assertion is uniform
as the extension varies.

\paragraph{The cited analytic input and its exact field requirement (NT-R05).}
Fiorilli--Jouve, \emph{Distribution of Frobenius elements in families of
Galois extensions}, author arXiv:2001.05428v2, Proposition3.18,
pp.32--34, proves the limiting-distribution theorem by a
Besicovitch $B^2$ approximation. Their Corollary3.17 and
Lemmata3.14--3.15, pp.30--32, supply the cyclic-induction version
which avoids Artin holomorphy. We use its specialization to
$\beta=1/2$, for real $t$. This gives the $B^2$ expansion below;
the full analytic almost-periodicity theorem is a cited classical
input. The precise specialization does not require $L/\mathbb Q$
to be Galois.

Here are all coefficients, with the conjugacy-class normalization
made explicit. For each conjugacy class $C$ choose $g_C\in C$,
put $H_C=\langle g_C\rangle$ and $F_C=L^{H_C}$, and define
\[
 a_t(\rho)=\sum_{C\subset G}\frac{|C|}{|G|}t(C)
       \sum_{\chi\in\widehat H_C}\overline{\chi(g_C)}
             \operatorname{ord}_{s=\rho}L(s,\chi,L/F_C).
 \tag{NT-R05a}
\]
The sum over $C$ means a sum over conjugacy classes. All functions
inside this formula are finite-order Hecke functions, with their
full local factors. The coefficient is independent of the choices
and is equivalently the residue at $\rho$ of the formal logarithmic
derivative attached to $t$.

Indeed, character orthogonality in the cyclic group and induction give
\[
 1_C=\frac{|C|}{|G|}
       \sum_{\chi\in\widehat H_C}\overline{\chi(g_C)}
                \operatorname{Ind}_{H_C}^G\chi.
\]
Artin induction for Euler factors, including inertia invariants,
then proves (NT-R05a). Moreover
$\zeta_L=\prod_{\chi\in\widehat H_C}L(s,\chi,L/F_C)$.
Every factor is holomorphic in the open critical strip. Thus RH for
$\zeta_L$ forces the nontrivial zeros of every factor in this formula
onto the critical line, without invoking holomorphy of any
nonabelian irreducible Artin function. This is the required analytic
input reduction in the cited cyclic argument.

Write $\Gamma$ for the set of distinct positive ordinates of zeros
of $\zeta_L$, and set
\[
 q_t=\frac1{|G|}\sum_{g\in G}t(g^2),\qquad
 b_t=-q_t-2a_t(1/2).
\]
Reality of $t$ gives $a_t(\overline\rho)=\overline{a_t(\rho)}$,
so $b_t$ is real. The cited expansion is
\[
 E_t(y)\ \sim_{B^2}\ b_t-
       2\Re\sum_{\gamma\in\Gamma}
        \frac{a_t(1/2+i\gamma)}{1/2+i\gamma}e^{i\gamma y}.
 \tag{NT-R05b}
\]
The notation means convergence of symmetric height truncations for
the seminorm
\[
 \|f\|_{B^2}^2=\limsup_{Y\to\infty}\frac1Y
                     \int_{\log2}^Y|f(y)|^2\,dy.
\]
The constant $-q_t$ is the prime-square term; the central zero
contributes $-2a_t(1/2)$. The local zero count NT-Z10 implies
\[
 \sum_{T<\gamma\le T+1}|a_t(1/2+i\gamma)|^2
       \ll_{L,K,t}\log^2(T+3),\qquad
 \sum_{\gamma\in\Gamma}
 \frac{|a_t(1/2+i\gamma)|^2}{1/4+\gamma^2}<\infty.
 \tag{NT-R05c}
\]
To check this directly, each nonnegative factor order in
(NT-R05a) is at most the corresponding order of $\zeta_L$.
The sum of these latter orders in a unit interval is $O_L(\log(T+3))$.
Taking absolute values in the finite class-character sum and then
squaring proves the first estimate; summing over integer intervals
proves the second. Multiplicities are retained throughout.

\paragraph{Complete joint-law and covariance deduction (NT-R06).}
For any finite list of real class functions $t_1,\ldots,t_d$, the
vector $(E_{t_1}(y),\ldots,E_{t_d}(y))$ has a limiting probability
law $\nu$ on $\mathbb R^d$. Precisely, the averages of every bounded
continuous function of this vector converge to its $\nu$ integral.
The mean vector is $(b_{t_j})_j$, and its covariance matrix is
\[
 \Sigma_{jk}=2\Re\sum_{\gamma\in\Gamma}
       \frac{a_{t_j}(1/2+i\gamma)
             \overline{a_{t_k}(1/2+i\gamma)}}{1/4+\gamma^2}.
 \tag{NT-R06a}
\]
Every series here converges absolutely. This is a full new
formulation of the joint deduction from the cited scalar expansion,
not a claim of literature novelty.

To prove it, take the closure $\mathcal T$ of the one-parameter
subgroup
$y\mapsto(e^{i\gamma y})_{\gamma\in\Gamma}$ in the countable
product of circles and give it Haar probability measure. Every
continuous character of a finite-dimensional projection has the
form $e^{i\lambda y}$ along this subgroup. Its Cesaro mean is zero
unless $\lambda=0$, in which case it is one. Stone--Weierstrass
therefore proves that the time averages of every continuous
function of finitely many coordinates equal the corresponding
Haar integrals. Distinct positive frequencies are orthogonal in
$L^2(\mathcal T)$; so are a positive frequency and the conjugate of
another positive frequency.

By (NT-R05c) the series
\[
 Z_j=b_{t_j}-2\Re\sum_{\gamma\in\Gamma}
       \frac{a_{t_j}(1/2+i\gamma)}{1/2+i\gamma}z_\gamma
\]
converges in $L^2(\mathcal T)$. Let $\nu$ be the law of this vector.
For any bounded Lipschitz test function, Cauchy--Schwarz bounds the
difference between its time average on $E$ and on a common finite
truncation by the Euclidean sum of the $B^2$ approximation errors.
The same bound holds for its Haar integrals on $Z$ and that
truncation. The finite-truncation averages converge by the preceding
character argument; letting the truncation grow proves convergence
for bounded Lipschitz tests. These tests determine weak convergence
of probability measures on $\mathbb R^d$, which gives the stated
bounded-continuous formulation. Taking Haar means and second moments
in the $L^2$ series proves (NT-R06a), including the factor two.
Cauchy--Schwarz and (NT-R05c) prove absolute convergence of its
mixed series.

No independence of the coordinates $z_\gamma$ was used: all rational
relations between the ordinates remain encoded in $\mathcal T$.
In particular this argument proves neither a Gaussian law nor
independence of zeros. A zero covariance in some direction means
exactly that its noncentral coefficients $a_t(1/2+i\gamma)$ all
vanish; its limiting law is then the point mass at $b_t$.

\paragraph{Finite etale fibres have a negative mean and a quantitative
one-sided lower density (NT-R07).}
Let $X\to U\subset\operatorname{Spec}\mathcal O_K$ be a nonempty
finite etale scheme, with generic algebra $\prod_{i=1}^r K_i$ and
Galois splitting field $L/K$. Put
\[
 \Pi_X(x)=\sum_{\substack{\mathfrak p\in U\\N\mathfrak p\le x}}
                         \#X(\mathcal O_K/\mathfrak p),\qquad
 E_X(y)=ye^{-y/2}\bigl(\Pi_X(e^y)-r\operatorname{Li}_2(e^y)\bigr).
\]
Apply NT-R05 to the permutation character $t$ of the generic
geometric fibre $S$. Omitting the finitely many additional primes
outside $U$ changes $\pi_t(x)$ by a bounded amount; its normalized
contribution tends to zero and has $B^2$ seminorm zero. Burnside's
lemma gives $\delta_t=r$. The full Artin permutation identity gives
\[
 a_t(1/2)=\sum_{i=1}^r\operatorname{ord}_{s=1/2}\zeta_{K_i}(s)
       \ge0.
\]
To evaluate the sign of $q_t$, decompose $S=\coprod_iG/H_i$.
Counting the pairs $(g,aH_i)$ fixed by $g^2$ gives
\[
 q_t=\sum_{i=1}^r
      \frac{\#\{g\in G:g^2\in H_i\}}{|H_i|}\ge r,
\]
because every element of $H_i$ belongs to the set being counted.
Thus the mean is $-B_X$, where
\[
 B_X=q_t+2\sum_i\operatorname{ord}_{1/2}\zeta_{K_i}\ge r>0,
 \qquad
 V_X=2\sum_{\gamma\in\Gamma}
          \frac{|a_t(1/2+i\gamma)|^2}{1/4+\gamma^2}<\infty.
\]
Consequently
\[
 \liminf_{X_0\to\infty}\frac1{\log X_0}
 \int_2^{X_0}
 1_{\{\Pi_X(x)<r\operatorname{Li}_2(x)\}}\,\frac{dx}{x}
       \ge\frac{B_X^2}{B_X^2+V_X}>0.
 \tag{NT-R07a}
\]
Here is the probability step in full. If $Z$ has the limiting law,
then $W=Z+B_X$ has mean zero and variance $V_X$. For $u\ge0$,
Markov's inequality gives
$\Pr(W\ge B_X)\le(V_X+u^2)/(B_X+u)^2$.
For $V_X>0$, choose $u=V_X/B_X$; for $V_X=0$ the variable
$Z=-B_X$ almost surely. In both cases
$\Pr(Z<0)\ge B_X^2/(B_X^2+V_X)$.
Weak convergence and the open set $(-\infty,0)$ imply the same
lower bound for the liminf of the time proportions. The substitution
$x=e^y$ gives (NT-R07a).

This is a positive lower logarithmic density. The existence of the
sign density, its being greater than one half, and densities of all
strict orderings are not asserted. They require additional control
of the limiting measure, notably its mass on the relevant boundary.
Geometrically these statements concern the trace of Frobenius on
finite-monodromy local systems and, for NT-R07, the permutation
realization of $H^0$ of a finite etale arithmetic cover. They use
ordinary Hecke/Dedekind realizations and do not invoke GRH for an
infinite-image geometric representation.

\subsection*{NT-R08. A discriminant-uniform lower logarithmic density}

Retain the finite etale setting of NT-R07 and assume only RH for
$\zeta_L$. There is an absolute constant $C>0$ such that
\[
 V_X\le C r^2\log^2(\Delta_L+2),\qquad
 \underline\delta\{\Pi_X(x)<r\operatorname{Li}_2(x)\}
 \ge \frac1{1+C\log^2(\Delta_L+2)}.
 \tag{NT-R08a}
\]
Here $\underline\delta$ denotes the lower logarithmic density
defined by the left side of (NT-R07a). The same constant works
for every number field $K$, every finite Galois splitting field
$L/K$, and every nonempty finite etale cover with this splitting
field. No lower bound on zero spacing or upper bound on individual
multiplicities independent of $\Delta_L$ is assumed. The
almost-periodic limiting process is still taken separately for
each fixed extension; no uniform rate of convergence to the law
is asserted.

\paragraph{Complete proof.}
Write $n=[L:\mathbb Q]$ and $D=\Delta_L$. Let
$m_L(\gamma)=\operatorname{ord}_{s=1/2+i\gamma}\zeta_L(s)$.
For every $\gamma>0$, the cyclic-character formula (NT-R05a)
and nonnegativity of the permutation character $t$ give
\[
 |a_t(1/2+i\gamma)|
 \le\sum_C\frac{|C|}{|G|}t(C)
       \sum_{\chi\in\widehat H_C}
          \operatorname{ord}_{1/2+i\gamma}L(s,\chi,L/F_C)
 =r\,m_L(\gamma).
 \tag{NT-R08b}
\]
This is an exact inequality: the inner sum equals the order
of $\zeta_L$, and the outer sum equals $r$ by Burnside.

The unconditional zero-count formula in Grenie--Molteni,
AppendixA TheoremA.1, author p.26, specialized to the trivial
character of $L$, implies with an absolute implied constant
\[
 U_j:=\sum_{\substack{\gamma\ {\rm distinct}\\j<\gamma\le j+1}}
                  m_L(\gamma)
       \ll \log D+n\log(j+3)\qquad(j=0,1,2,\ldots).
 \tag{NT-R08c}
\]
For clarity, the source counts both signs and includes
multiplicities. Its main term is
$T\{\log D+n\log(T/(2\pi e))\}/\pi$, with error
$O(\log D+n\log T+n)$ for $T\ge1$, and an additional
$O(n)$ constant term. For $j\ge2$ subtract its values at
$j-1$ and $j+1$; the difference of the main terms has the
displayed size by differentiation, and the sum of the two
error bounds has that size as well. This bounds the positive
ordinates in $(j,j+1]$. For $j=0,1$ use the source bound for
the total count at $T=3$. The absolute constants in its table
therefore give (NT-R08c) uniformly in the field.

Within a unit interval the orders are nonnegative, so the sum
of their squares is at most $U_j^2$. Equations (NT-R06a) and
(NT-R08b) yield
\[
 \begin{split}
 V_X
 &\le 2r^2\sum_{\gamma>0}
                  \frac{m_L(\gamma)^2}{1/4+\gamma^2}\\
 &\ll r^2\sum_{j=0}^{\infty}
          \frac{(\log D+n\log(j+3))^2}{1/4+j^2}\\
 &\ll r^2(\log D+n)^2.
 \end{split}
 \tag{NT-R08d}
\]
The final constant is absolute, since the series with numerator
$1$, $\log(j+3)$ or $\log^2(j+3)$ all converge against this
denominator.

We include the elementary removal of the degree parameter.
Minkowski's discriminant bound gives
\[
 \sqrt D\ge\frac{n^n}{n!}\left(\frac{\pi}{4}\right)^{r_2}.
\]
An exact reference is Milne, \emph{Algebraic Number Theory},
v3.08, Theorem4.3 printed p.70 and the proof of Theorem4.9
printed p.72. For $n\ge2$, the arithmetic-geometric mean
inequality gives $n!\le((n+1)/2)^n$, and $r_2\le n/2$.
Consequently
\[
 \sqrt D\ge
 \left(\frac{2n}{n+1}\sqrt{\frac{\pi}{4}}\right)^n
 \ge\left(\frac{2\sqrt\pi}{3}\right)^n.
\]
Since $2\sqrt\pi/3>1$, this proves
$n\ll\log D$ with an absolute constant for $n\ge2$.
The case $n=1,D=1$ is handled by $\log(D+2)$.
Substitution in (NT-R08d) proves the first assertion of
(NT-R08a). Finally $B_X\ge r$ in NT-R07 gives
\[
 \frac{B_X^2}{B_X^2+V_X}
 \ge \frac1{1+C\log^2(D+2)},
\]
which proves the lower-density assertion.

The variance estimate is consistent with the squared
$\|t\|_1^2$ bound in Fiorilli--Jouve Proposition4.6(87),
author p.40 and its proof p.41. The proof above derives the
coarse bound directly from the uniform zero count and gives
the finite etale consequence without invoking a bounded
multiplicity hypothesis. No numerical value for the absolute
constant is asserted.

In the Grothendieck language of NT-F02 and NT-R07, this is a
uniform lower-frequency statement for the Frobenius trace on
the permutation realization of $H^0$ of a finite etale cover.
The generic rank of that realization may grow; the final lower
bound depends on the splitting-field discriminant, not that
rank. It does not give a uniform threshold in $x$ or a uniform
error term for the approach to the limiting law.

\subsection*{NT-E06. Uniform CM angle counts from one Hecke-power family}
Let $E/k$ be an elliptic curve whose complex multiplication is defined
over the number field $k$. Put
\[
 a_{\mathfrak p}=N\mathfrak p+1-\#E_{\mathfrak p}(k_{\mathfrak p}),
 \qquad t_{\mathfrak p}=\frac{a_{\mathfrak p}}{\sqrt{N\mathfrak p}}.
\]
Only primes of good reduction are counted. The classical CM character
theorem supplies a unitary Hecke character $\eta$ of $k$ for which
\[
 t_{\mathfrak p}=\eta(\mathfrak p)+\overline{\eta(\mathfrak p)}
                =2\cos(2\pi\theta_{\mathfrak p}),
 \qquad \theta_{\mathfrak p}\in\mathbb R/\mathbb Z.       \tag{NT-E06a}
\]
Assume RH for $L_k(s,\eta^m)$ for every integer $m$, including
$m=0$, where the function is $\zeta_k$. This is a precisely specified
subfamily of $M$; Dirichlet GRH alone does not supply its nonzero
infinity types. Then, uniformly over all arcs $J\subset\mathbb R/\mathbb Z$
and all intervals $I\subset[-2,2]$, for $x\ge2$,
\begin{align}
 \#\{\mathfrak p:N\mathfrak p\le x,\ \theta_{\mathfrak p}\in J\}
   &= |J|\operatorname{Li}_2(x)
       +O_{E,k}\bigl(\sqrt x\log^2(2x)\bigr), \tag{NT-E06b}\\
 \#\{\mathfrak p:N\mathfrak p\le x,\ t_{\mathfrak p}\in I\}
   &= \mu(I)\operatorname{Li}_2(x)
       +O_{E,k}\bigl(\sqrt x\log^2(2x)\bigr), \tag{NT-E06c}\\
 d\mu(t)&=\frac{dt}{\pi\sqrt{4-t^2}}.\notag
\end{align}
The constants are independent of the interval and its endpoint choices.
In particular the intervals may vary with $x$. No numerical implicit
constant or uniformity in the curve or base field is asserted.

\paragraph{Classical input and complete deduction.}
We use the Hecke prime-sum estimate recorded by Bucur--Fit\'e--Kedlaya,
\emph{Effective Sato--Tate conjecture for abelian varieties and
applications}, Proposition2.5, equation(2.3), published p.1720
(author v2 p.6). Its CM interpretation is explicitly identified in
footnote2 to Proposition4.7, published p.1740 (author p.20).
For the fixed curve and field it gives, uniformly in $m\ne0$ and $x\ge2$,
\begin{align}
 S_m(x):=\sum_{N\mathfrak p\le x}\eta(\mathfrak p)^m
       &=O_{E,k}\bigl(\sqrt x\log(2x+2|m|)\bigr),\notag\\
 S_0(x)&=\operatorname{Li}_2(x)
       +O_{E,k}\bigl(\sqrt x\log(2x)\bigr).               \tag{NT-E06d}
\end{align}
Their sums omit the primes over the absolute conductor; restoring any
good primes omitted by that convention changes a sum by $O_{E,k}(1)$.
The representation dimension in this specialization is one, while the
motivic weight is $|m|$: the distinction is the point of their
Remark2.7. The finite conductor of $\eta^m$ divides that of $\eta$,
and its archimedean parameters grow at most linearly in $|m|$.
Every nonzero power has a nonzero angular infinity type, so is not
a pure norm twist. Consequently the required entire continuation,
functional equation, and boundary nonvanishing are classical Hecke
facts; the only conjectural analytic input in (NT-E06d) is the
stated RH family. The standard Hecke completion over $k$ has
degree $[k:\mathbb Q]$, which is fixed here.

For clarity, the elementary harmonic-analysis input is
Erd\H{o}s--Tur\'an, \emph{On a problem in the theory of uniform
distribution.I} (1948), TheoremIII, printed p.374 (parallel
proceedings pagination p.1150). For $P$ points on the circle and an
integer $H\ge1$ its discrepancy bound is
\[
 \left|\#\{\theta_j\in J\}-|J|P\right|
 \le C\left(\frac P{H+1}
       +\sum_{m=1}^{H}\frac1m
                    \left|\sum_{j=1}^{P}e^{2\pi i m\theta_j}\right|\right),
                                                               \tag{NT-E06e}
\]
where $C$ is absolute. The original statement includes closed
endpoints. The same bound, after changing $C$, holds with any endpoint
convention: shrink or enlarge an arc, use the same bound, and pass to
the limit for the finite set of points. Singletons are handled by this
argument as well.

Apply (NT-E06e) with $H=\lceil x\rceil$ and the angles in (NT-E06a).
The first term is $O_k(1)$, since the number of prime ideals of norm at
most $x$ is at most $[k:\mathbb Q]x$. By (NT-E06d) the sum is at most
\[
 C_{E,k}\sqrt x\log(4x+2)\sum_{m\le\lceil x\rceil}\frac1m
     =O_{E,k}\bigl(\sqrt x\log^2(2x)\bigr).
\]
Replacing $P$ by $\operatorname{Li}_2(x)$ using $m=0$ proves
(NT-E06b). The inverse image of a trace interval under
$\theta\mapsto2\cos(2\pi\theta)$ is the union of at most two arcs.
The pushforward of normalized Lebesgue measure is exactly $\mu$;
adding the two estimates proves (NT-E06c), with singleton endpoints
already accounted for. This proves all the stated uniformities.

\paragraph{Geometric meaning and locator.}
The local system $R^1f_*\mathbb Q_\ell$ of a smooth elliptic scheme over
an arithmetic open has rank two and weight one. CM splits its scalar
extension into two conjugate rank-one characters; their unitary
normalizations give the circle above. Its tensor powers are the
specific Hecke family used in the proof. Thus this is a Frobenius
distribution statement for a CM torus in the motivic monodromy group,
not a deduction of GRH for an arbitrary non-CM elliptic $L$-function.
Work/run: \texttt{GRH\_NIGHT/number\_theory},
\texttt{cm\_prime\_angles.tex}, NT-E06; sources S-NT-BFK and S-NT-ET.

\subsection*{NT-E07. A least-prime bound in any positive-measure CM interval}
Under the exact hypotheses of NT-E06, for any trace interval $I$ with
$\mu=\mu(I)>0$, there is a good prime ideal with $t_{\mathfrak p}\in I$
and
\[
 N\mathfrak p\le C_{E,k}\,\mu^{-2}\log^6(e/\mu).           \tag{NT-E07}
\]
The analogous conclusion holds for an arc of length $\mu$. This is a
newly proved formulation of the classical discrepancy consequence;
no claim of literature novelty or numerical value of $C_{E,k}$ is made.

\paragraph{Complete proof.}
Let $A_{E,k}$ be a common error constant in (NT-E06b)--(NT-E06c).
For $x\ge4$,
$\operatorname{Li}_2(x)\ge x/(2\log x)$, by integrating over $[x/2,x]$.
Write $L=\log(e/\mu)\ge1$ and set $x=C\mu^{-2}L^6$, with $C\ge4$.
Since $\log L\le L$ and $\log(1/\mu)\le L$, we have
\[
 \log(2x)\le(\log(2C)+8)L,\qquad
 \mu\sqrt x=\sqrt C\,L^3.
\]
The prime count is positive provided
$\mu\sqrt x>2A_{E,k}\log x\,\log^2(2x)$.
It is therefore enough to choose the fixed $C$ so that
\[
 \sqrt C>2A_{E,k}\bigl(\log(2C)+8\bigr)^3.
\]
Such a $C$ exists because $\sqrt C/\log^3(2C)\to\infty$.
This choice is independent of $I$, proves the bound, and also covers
intervals with endpoints included or excluded. Geometrically it
detects a prescribed positive-measure region in the compact CM
monodromy torus by one closed arithmetic fibre. Work/run:
\texttt{GRH\_NIGHT/number\_theory}, \texttt{cm\_prime\_angles.tex}, NT-E07.

\subsection*{NT-E08. An extremal CM trace asymptotic with a power-saving error}
Let now $E/k$ be any elliptic curve with potential CM by the imaginary
quadratic field $K$. Put $F=kK$ and $d=[F:k]\in\{1,2\}$. Let $\eta$
be a CM unitary Hecke character of $E_F$. Assume RH for the functions
$L_F(s,\eta^m)$ for every integer $m$, including $m=0$.
For either choice of sign, define
\[
 M^\pm_{E/k}(x)
  =\#\{\mathfrak p\text{ good}:N\mathfrak p\le x,\
               a_{\mathfrak p}=\pm\lfloor2\sqrt{N\mathfrak p}\rfloor\}.
\]
Then
\begin{align}
 M^\pm_{E/k}(x)
   &=\frac1{\pi d}\int_2^x\frac{t^{-1/4}}{\log t}\,dt
       +O_{E,k}\bigl(x^{5/8}\log^2(2x)\bigr),             \tag{NT-E08a}\\
 M^\pm_{E/k}(x)
   &\sim \frac4{3\pi d}\frac{x^{3/4}}{\log x}.
                                                               \tag{NT-E08b}
\end{align}
In particular $M$ supplies the result conditionally. The factor is
$4/(3\pi)$ when all CM is defined over $k$, and $2/(3\pi)$ in the
quadratic descent case. The integral in (NT-E08a) is retained because
replacing it by its first asymptotic term would introduce a larger
logarithmic error. This is a fully proved formulation in this run,
not an assertion that the asymptotic or its constant is new in the
literature.

\paragraph{Complete proof when $d=1$.}
It suffices first to consider the positive sign. If $a$ is an integer
with $a\le y$, then $a=\lfloor y\rfloor$ if and only if $a>y-1$.
Thus the Hasse bound gives the exact identity
\[
 a_{\mathfrak p}=\lfloor2\sqrt{N\mathfrak p}\rfloor
 \quad\Longleftrightarrow\quad
 t_{\mathfrak p}\in I_{N\mathfrak p},\qquad
 I_u=(2-u^{-1/2},2].                                    \tag{NT-E08c}
\]
The strict lower endpoint is essential and makes the identity valid
also when $N\mathfrak p$ is a square. No fractional-part factor occurs.
Direct integration of the arcsine density gives, uniformly for $u\ge2$,
\[
 \mu(I_u)=\frac1\pi\arccos\!\left(1-\frac1{2\sqrt u}\right)
         =\frac1\pi u^{-1/4}+O(u^{-3/4}).                 \tag{NT-E08d}
\]
For example the last estimate follows by setting $v=u^{-1/2}$ and
integrating $(4z-z^2)^{-1/2}=(2\sqrt z)^{-1}(1+O(z))$
from $0$ to $v$; the constants are absolute on $0\le v\le2^{-1/2}$.

For $x\ge4$ set $h=x^{-1/8}$ and partition $[2,x]$ into intervals
$(u_{j+1},u_j]$, with $u_0=x$, successive ratios $u_j/u_{j+1}=1+h$,
and the final lower endpoint truncated to $2$. Every ratio is at most
$1+h$. Since $I_u$ decreases with $u$, the contribution on this
interval is bracketed by the counts using the two fixed intervals
$I_{u_j}$ and $I_{u_{j+1}}$. Apply (NT-E06c) at both endpoints.
After summing, the total error in each bracket is at most
\[
 C_{E,k}\sum_j\sqrt{u_j}\log^2(2u_j)
       \ll_{E,k}h^{-1}\sqrt x\log^2(2x),                 \tag{NT-E08e}
\]
because $\sum_{j\ge0}(1+h)^{-j/2}\ll h^{-1}$ for $0<h\le1$.
Primes of norm exactly $2$ contribute $O_k(1)$.

The main brackets are the upper and lower Riemann sums of the
decreasing function $\mu(I_u)$ against the positive measure
$d\operatorname{Li}_2(u)=du/\log u$. Replacing $\mu(I_u)$ using
(NT-E08d) introduces an error bounded by a constant times
\[
 \int_2^x\frac{u^{-3/4}}{\log u}\,du+1
      \ll x^{1/4}+1.                                   \tag{NT-E08f}
\]
Indeed values of $u^{-3/4}$ at the two endpoints differ by a factor
at most $(1+h)^{3/4}$. For $f(u)=u^{-1/4}$ the corresponding
upper and lower sums differ by at most
$((1+h)^{1/4}-1)$ times the lower sum. They therefore bracket
$\int_2^x u^{-1/4}\,du/\log u$ with error
\[
 \ll h\int_2^x\frac{u^{-1/4}}{\log u}\,du
 \ll h x^{3/4}.                                       \tag{NT-E08g}
\]
Combining (NT-E08e)--(NT-E08g), then substituting $h=x^{-1/8}$,
proves (NT-E08a) when $d=1$. Applying exactly the same argument to
$-t_{\mathfrak p}$ proves the negative sign, since $\mu$ is symmetric.
Finally l'H\^opital's rule, or a single integration by parts, gives
\[
 \int_2^x\frac{u^{-1/4}}{\log u}\,du
      \sim \frac43\frac{x^{3/4}}{\log x},
\]
and $x^{5/8}\log^2(2x)=o(x^{3/4}/\log x)$.

\paragraph{Complete quadratic descent.}
The field defining all endomorphisms is $F=kK$. One can see this
directly on invariant differentials: in characteristic zero the
differential embeds $\operatorname{End}(E_{\bar k})\otimes\mathbb Q=K$
faithfully in $\bar k$, compatibly with $G_k$. Its kernel of Galois
action is therefore $G_{kK}$. If $d=2$, at an unramified good prime
inert in $F/k$, Frobenius interchanges the two CM eigenspaces of the
Tate module. Its matrix has zero diagonal, so $a_{\mathfrak p}=0$.
Such a prime is never extremal because $\lfloor2\sqrt{N\mathfrak p}\rfloor
\ge2$.

Every split good prime of $k$ has two primes of $F$ of the same norm
and the same Frobenius trace. A prime of $F$ above an inert prime has
norm $(N\mathfrak p)^2$; consequently the number of such primes of
norm at most $x$ is $O_k(\sqrt x)$, even by the elementary bound for
prime ideals above rational primes. Ramified primes and primes where
good reduction changes under base extension form a fixed finite set.
It follows that
\[
 M^\pm_{E/F}(x)=2M^\pm_{E/k}(x)+O_{E,k}(\sqrt x).
\]
Apply the already proved connected-CM case over $F$ and divide by
two. This proves both formulas for $d=2$, with the same stated
error scale.

\paragraph{Literature comparison, geometry, and proof location.}
Bucur--Fit\'e--Kedlaya Proposition4.9 and Corollary4.10, published
pp.1741--1742, give the order $x^{3/4}/\log x$ and the quadratic
base-change mechanism under this same Hecke-power input.
James--Pollack, \emph{Extremal primes for elliptic curves with
complex multiplication} (2017), Theorem1 (author p.2), states the
$2/(3\pi)$ asymptotic over $\mathbb Q$ unconditionally; its separate
erratum replaces the narrow-sector analytic input. That published
result and its erratum are comparison references, not inputs to
the proof above. This run does not certify the erratum's sector
estimate independently.

The arithmetic content is the number of closed fibres for which the
trace on $H^1$ reaches the integral Hasse endpoint. The CM torus and
its rank-one Hecke characters supply all the analytic hypotheses;
there is no implication here for extremal traces of a non-CM curve.
Work/run: \texttt{GRH\_NIGHT/number\_theory},
\texttt{cm\_prime\_angles.tex}, NT-E08; complete deduction from NT-E06,
with source inputs S-NT-BFK and S-NT-ET and comparison S-NT-JP.

\subsection*{NT-E09. Shrinking absolute deficits from the CM Hasse endpoint}
Retain $E/k$, its imaginary quadratic CM field $K$, $F=kK$, and
$d=[F:k]\in\{1,2\}$ from NT-E08. Fix real numbers
\[
 c>0,\qquad -\frac12<\lambda<\frac12,\qquad
 \alpha=\frac14-\frac{\lambda}{2}\in(0,\tfrac12).
\]
Set $u_0=\max\{2,c^{1/(1/2-\lambda)}\}$, so that
$c u_0^{\lambda-1/2}\le1$, and put, for $u\ge u_0$,
\[
 g_{c,\lambda}(u)=\frac1\pi
       \arccos\!\left(1-\frac c2u^{\lambda-1/2}\right).
                                                               \tag{NT-E09a}
\]
For either sign, let
\[
 D^\pm_{E/k;c,\lambda}(x)
 =\#\{\mathfrak p\text{ good}:N\mathfrak p\le x,\
              0\le2\sqrt{N\mathfrak p}\mp a_{\mathfrak p}
                           <c(N\mathfrak p)^\lambda\}.
                                                               \tag{NT-E09b}
\]
Assume RH for all integral powers of the unitary CM Hecke character
of $E_F$, including its zeroth power $\zeta_F$. Then, as $x\to\infty$,
\begin{align}
 D^\pm_{E/k;c,\lambda}(x)
   &=\frac1d\int_{u_0}^{x}\frac{g_{c,\lambda}(u)}{\log u}\,du
       +O_{E,k,c,\lambda}
           \bigl(x^{5/8+\lambda/4}\log^2(2x)\bigr),       \tag{NT-E09c}\\
 D^\pm_{E/k;c,\lambda}(x)
   &\sim\frac{\sqrt c}{\pi d(3/4+\lambda/2)}
                     \frac{x^{3/4+\lambda/2}}{\log x}.   \tag{NT-E09d}
\end{align}
The same statements hold if only prime ideals of absolute residue
degree one are counted. All parameters $E,k,c,\lambda$ are fixed;
the constants and threshold are not asserted uniform as
$\lambda$ approaches either endpoint.
For $-1/2<\lambda<0$ this gives infinitely many primes with an
absolute Hasse deficit tending to zero. This is a new complete
formulation in the run, without a claim of literature priority.

\paragraph{Complete proof of the moving-target estimate.}
First suppose $d=1$. For the positive sign the condition in
(NT-E09b), at norm $q\ge u_0$, is precisely
\[
 t_{\mathfrak p}\in
  J_q=(2-cq^{\lambda-1/2},2].
\]
The arcsine measure of $J_u$ is exactly $g_{c,\lambda}(u)$.
Let $v=cu^{-2\alpha}$, so $0<v\le1$. Elementary differentiation
and integration give
\[
 g_{c,\lambda}(u)\asymp_{c,\lambda}u^{-\alpha},
 \qquad
 -u\,g'_{c,\lambda}(u)
       =\frac{\alpha\sqrt v}{\pi\sqrt{1-v/4}}.            \tag{NT-E09e}
\]
In fact $g(u)\ge\sqrt v/\pi$, because
$g(u)=\pi^{-1}\int_0^v(4z-z^2)^{-1/2}\,dz$.
Since $v\le1$, (NT-E09e) consequently yields the explicit bound
\[
 0\le-\frac{u g'(u)}{g(u)}\le\frac{2\alpha}{\sqrt3}.
                                                               \tag{NT-E09f}
\]
Thus $g$ is decreasing and its values on a multiplicative interval
of ratio at most $1+h$ differ by a factor at most
$(1+h)^{2\alpha/\sqrt3}=1+O_\lambda(h)$ for $0<h\le1$.

Set $h=x^{-1/4+\alpha/2}$, which tends to zero because $\alpha<1/2$.
Partition $[u_0,x]$ into successive multiplicative intervals of
ratio $1+h$, truncating the last one at $u_0$. On each interval,
the moving interval $J_q$ lies between its two endpoint intervals.
Applying NT-E06 at both norm endpoints and adding the errors gives
\[
 O_{E,k}\!\left(\sqrt x\,h^{-1}\log^2(2x)\right),
                                                               \tag{NT-E09g}
\]
by the geometric series already evaluated in (NT-E08e).
The two main terms are Riemann sums for
$\int_{u_0}^x g(u)\,du/\log u$. By (NT-E09f) their difference is
bounded by a constant times
\[
 h\int_{u_0}^x\frac{g(u)}{\log u}\,du
    \ll_{c,\lambda}h x^{1-\alpha}.                      \tag{NT-E09h}
\]
The finite range below $u_0$ adds $O_{E,k,c,\lambda}(1)$.
Both powers in (NT-E09g)--(NT-E09h) equal
$x^{3/4-\alpha/2}=x^{5/8+\lambda/4}$. This proves (NT-E09c)
for $d=1$. Reflection of the trace interval proves the negative
sign with the identical mass and constants.

To obtain (NT-E09d), (NT-E09a) gives
$g(u)\sim(\sqrt c/\pi)u^{-\alpha}$. Comparison of the positive
integrals, followed by differentiation of
$x^{1-\alpha}/((1-\alpha)\log x)$, yields
\[
 \int_{u_0}^x\frac{g(u)}{\log u}\,du
     \sim\frac{\sqrt c}{\pi(1-\alpha)}
                           \frac{x^{1-\alpha}}{\log x}.
\]
The ratio of the error in (NT-E09c) to this main term is
$O(x^{-1/4+\alpha/2}\log^3(2x))=o(1)$, exactly because
$\alpha<1/2$.

If $d=2$, the nonsplit good primes have $a_{\mathfrak p}=0$ by
the CM eigenspace argument in NT-E08. They cannot satisfy
(NT-E09b) once $q$ exceeds a fixed threshold, since
$2\sqrt q>c q^\lambda$ for large $q$. Split primes again have two
lifts of equal norm and trace, while nonsplit primes upstairs have
square norms. The same $O_{E,k}(\sqrt x)$ correction therefore gives
the factor $1/d$. Finally, the number of prime ideals of any fixed
number field with absolute residue degree at least two and norm at
most $x$ is $O_k(\sqrt x)$: their underlying rational prime is at
most $\sqrt x$, and there are at most $[k:\mathbb Q]$ primes above it.
Because $5/8+\lambda/4>1/2$, deleting them does not change either
formula.

\paragraph{Exact range, analytic input, and geometric meaning.}
The small-angle expansion is used only to identify the asymptotic
constant. The exact arcsine integral is retained in the
power-saving statement. In particular an error of size
$\int^{x}u^{-3\alpha}\,du/\log u$ from replacing it by the first
term need not fit within (NT-E09c); for example it is larger when
$\lambda>3/10$. No such replacement is made.

At $\lambda=-1/2$ the predicted exponent becomes $1/2$ and this
argument no longer gives a smaller error. No endpoint assertion,
quadratic-polynomial prime theorem, or corresponding non-CM
statement is claimed. At $\lambda=0,c=1$ the strict inequality
recovers the exact integer extremal event in NT-E08. The result
describes shrinking neighbourhoods of a scalar Frobenius in the
compact CM monodromy torus of $H^1$, while retaining the actual
integer trace and absolute norm of the arithmetic fibre.

The only analytic hypothesis is the same power family as NT-E06;
$M$ supplies it conditionally. Sources S-NT-BFK and S-NT-ET supply
the classical prime-sum/discrepancy premises already audited there.
Work/run: \texttt{GRH\_NIGHT/number\_theory},
\texttt{cm\_shrinking\_deficits.tex}, NT-E09.

\subsection*{NT-A04. General Hecke angle discrepancy and its norm-twist exception}
Let $K$ be a number field and let $\eta$ be a unitary Hecke character
which is \emph{not} a finite-order character times a unitary power
of the norm. For every prime ideal unramified for $\eta$, define
$\theta_{\mathfrak p}\in\mathbb R/\mathbb Z$ by
$\eta(\mathfrak p)=e^{2\pi i\theta_{\mathfrak p}}$.
Assume RH for $L_K(s,\eta^m)$ for all positive integers $m$, and RH
for $\zeta_K$. Then, uniformly over all circle arcs $J$,
\[
 \#\{\mathfrak p\text{ unramified}:N\mathfrak p\le x,\
                          \theta_{\mathfrak p}\in J\}
 =|J|\operatorname{Li}_2(x)
       +O_{K,\eta}\bigl(\sqrt x\log^2(2x)\bigr)
 \quad(x\ge2).                                         \tag{NT-A04a}
\]
Any endpoint convention is allowed. In particular, uniformly for
$z\in\mathbb C$,
\[
 \#\{\mathfrak p\text{ unramified}:N\mathfrak p\le x,\
                                  \eta(\mathfrak p)=z\}
       \ll_{K,\eta}\sqrt x\log^2(2x).                    \tag{NT-A04b}
\]
The hypothesis on the character is essential for the angular law;
infinite order by itself is insufficient.

\paragraph{Exact classical locator and deductions.}
Rajan, \emph{Distribution of values of Hecke characters of infinite
order}, Acta Arithmetica85 (1998), Proposition4, printed
pp.289--290, gives (NT-A04a) with $\pi_K(x)$ in place of
$\operatorname{Li}_2(x)$. The character restriction is the one fixed
in Theorem1, printed p.286: in that paper a \emph{Dirichlet character}
means a finite-order Hecke character of $K$, and equivalence means
equality up to a power of the norm. The arguments use normalized
angles modulo one. Proposition4 has a constant depending only on
$K$ above a character-dependent threshold; allowing the constant
to depend on the fixed $\eta$ gives the displayed range $x\ge2$.
The finite ramification set is omitted throughout. Its analytic
premises are ordinary Hecke continuation and the precisely stated
power RH family; $M$ supplies them conditionally.

The prime-ideal estimate NT-C02 for the trivial extension $K/K$
replaces $\pi_K(x)$ by $\operatorname{Li}_2(x)$ with an error
$O_K(\sqrt x\log(2x))$. Shrinking or enlarging an arc and passing
to the limit for the finite collection of primes handles all
endpoint conventions, exactly as in NT-E06. For $|z|=1$, shrink
arcs to the singleton $z$; their main terms tend to zero while the
uniform error is fixed. This proves (NT-A04b). For $|z|\ne1$ the
set is empty.

\paragraph{Complete check of the norm-twist obstruction.}
Suppose instead that $\eta(\mathfrak p)=
\chi(\mathfrak p)(N\mathfrak p)^{it}$, where $\chi$ has finite order
$h$ and $t\ne0$. Set $\tau=ht\ne0$. Then the $h$th Fourier moment
of its prime angles is
\[
 \frac1{\pi_K(x)}\sum_{N\mathfrak p\le x}\eta(\mathfrak p)^h
       =\frac{x^{i\tau}}{1+i\tau}+o_{K,\chi,t}(1).       \tag{NT-A04c}
\]
Omitting the fixed ramification set does not affect this formula.
Here is a direct proof under the same available input $M$.
By NT-C02,
$\pi_K(u)=\operatorname{Li}_2(u)+O_K(\sqrt u\log(2u))$.
Abel summation therefore gives
\[
 \sum_{N\mathfrak p\le x}(N\mathfrak p)^{i\tau}
 =\int_2^x\frac{u^{i\tau}}{\log u}\,du
      +O_{K,\tau}(\sqrt x\log(2x)+1).
\]
Integration by parts with primitive
$u^{1+i\tau}/(1+i\tau)$ shows that the integral is
\[
 \frac{x^{1+i\tau}}{(1+i\tau)\log x}
           +O_\tau\!\left(\frac{x}{\log^2 x}+1\right).
\]
For the remainder, bound
$\int_2^xdu/\log^2u\ll x/\log^2x$, splitting at $\sqrt x$ if
desired. Since $\pi_K(x)\sim x/\log x$, (NT-A04c) follows.
Its absolute value tends to $(1+\tau^2)^{-1/2}>0$, whereas a
uniform angular distribution would force every nonzero Fourier
moment to tend to zero. Thus even under $M$ these infinite-order
characters do not obey (NT-A04a). The source exclusion prevents
the norm-twist poles from being silently treated as entire
nonprincipal factors.

\paragraph{Geometric dictionary and location.}
The general statement concerns rank-one unitary characters of the
idele class group. If the character is the unitary normalization
of an algebraic Hecke character, this is the Frobenius-value
distribution of the associated rank-one CM motive and its tensor
powers. A general unitary infinity type is not automatically
asserted to be an algebraic or \'{e}tale motive.
Work/run: \texttt{GRH\_NIGHT/number\_theory},
\texttt{hecke\_shrinking\_angles.tex}, NT-A04; primary source
S-NT-RAJAN, with the complete exception proof above.

\subsection*{NT-A05. Normalized shrinking arcs for infinite-order Hecke values}
Under the exact assumptions of NT-A04, fix
$c>0$ and $0<\alpha<1/2$, and set
$u_0=\max\{2,c^{1/\alpha}\}$.
For a center $\vartheta\in\mathbb R/\mathbb Z$ and $u\ge u_0$,
let $J_u(\vartheta)$ be the arc centered at $\vartheta$ of length
$cu^{-\alpha}$, with any fixed endpoint convention. Then, as $x\to\infty$,
\begin{align}
 &\#\{\mathfrak p\text{ unramified}:u_0<N\mathfrak p\le x,\
               \theta_{\mathfrak p}\in J_{N\mathfrak p}(\vartheta)\}
       \notag\\
 &\hspace{12mm}
   =c\int_{u_0}^x\frac{u^{-\alpha}}{\log u}\,du
       +O_{K,\eta,c,\alpha}
          \bigl(x^{3/4-\alpha/2}\log^2(2x)\bigr)
       \notag\\
 &\hspace{12mm}\sim\frac{c}{1-\alpha}\frac{x^{1-\alpha}}{\log x}.
                                                               \tag{NT-A05}
\end{align}
The error constant is uniform in the chosen center $\vartheta$.
All other parameters are fixed, and no assertion at $\alpha=1/2$
is made.

\paragraph{Complete proof.}
Set $h=x^{-1/4+\alpha/2}$. Partition $[u_0,x]$ into
successive multiplicative intervals of ratio $1+h$, with the
last endpoint truncated at $u_0$. The arcs are nested for the
fixed center. On each norm interval the moving prime count is
bounded above and below by the counts for its two endpoint arcs.
Apply (NT-A04a) at both norm endpoints, with a common constant
independent of $\vartheta$. The sum of all endpoint errors is
\[
 \ll_{K,\eta}\log^2(2x)
      \sum_{j\ge0}\sqrt{x(1+h)^{-j}}
 \ll_{K,\eta}\sqrt x\,h^{-1}\log^2(2x).
\]
The main brackets are the upper and lower Riemann sums for
$c\int_{u_0}^x u^{-\alpha}\,du/\log u$.
On an interval of ratio at most $1+h$, the endpoint values of
$u^{-\alpha}$ differ by a factor at most
$(1+h)^\alpha=1+O_\alpha(h)$. The difference of these Riemann
sums is consequently
$O_{c,\alpha}(h x^{1-\alpha})$.
With the chosen $h$, both power terms are $x^{3/4-\alpha/2}$.
This proves the equality in (NT-A05), with uniformity in the
center and the stated endpoint conventions.

Finally, integration by parts gives
$\int_{u_0}^x u^{-\alpha}\,du/\log u
\sim x^{1-\alpha}/((1-\alpha)\log x)$.
The ratio of the error to that main term is
$O(x^{-1/4+\alpha/2}\log^3(2x))=o(1)$.
This proves the asymptotic and the exact open range
$0<\alpha<1/2$. The window is normalized in the compact rank-one
torus; it is an arithmetic prime-counting assertion, and it does
not identify zeros of the Hecke function with independent random
angles. Work/run: \texttt{GRH\_NIGHT/number\_theory},
\texttt{hecke\_shrinking\_angles.tex}, NT-A05.

\subsection*{NT-E10. The degree-one endpoint reduces to finitely many CM norm forms}
Let $E/k$ have complex multiplication by the imaginary quadratic
field $K$, let $F=kK$, and put $D=|\operatorname{disc}K|$.
Fix $c>0$. This statement is unconditional. Consider good prime
ideals $\mathfrak p$ of $k$ of absolute residue degree one, so that
$N\mathfrak p=p$ is a rational prime, and fix either sign
$a=\pm a_{\mathfrak p}$. Write
\[
 \delta=2\sqrt p-a,\qquad B=4p-a^2 .
\]
Outside the finite ramification set of $F/k$, a prime satisfying
$0\le\delta<c/\sqrt p$ and $p>c/2$ must split in $F/k$.
At every such split prime there is a nonzero integer $b$ with
\[
 B=Db^2,\qquad
 \pi=\frac{a+b\sqrt{\operatorname{disc}K}}2\in\mathcal O_K,\qquad
 N_{K/\mathbb Q}\pi=p.                                  \tag{NT-E10a}
\]
For each fixed $c$, at all sufficiently large split good
degree-one primes the endpoint event is precisely
\[
 0\le2\sqrt p-a<\frac c{\sqrt p}
 \quad\Longleftrightarrow\quad
 a>0\ \hbox{and}\quad 0<Db^2<4c.                        \tag{NT-E10b}
\]
The equality $Db^2=4c$ never qualifies. Consequently, if $4c\le D$,
there are only finitely many degree-one good primes satisfying
either endpoint event. Outside the ramification set of $F/k$,
every possible exception then has $p\le c/2$.

\paragraph{Complete algebraic reduction.}
At a split good prime, all geometric CM endomorphisms specialize
to endomorphisms defined over $\mathbb F_p$. Their rational algebra
$K$ embeds in $\operatorname{End}^0(E_{\overline{\mathbb F}_p})$.
Here specialization is faithful: for $\ell\ne p$, good reduction
identifies the Tate modules, the action of an endomorphism on them
is compatible with specialization, and a nonzero endomorphism
cannot kill all $\ell$-power torsion. Indeed a nonzero endomorphism
of an elliptic curve is an isogeny with finite kernel.

The centralizer of this $K$ inside the rational endomorphism
algebra is $K$ itself. To check this without a classification of
the reduction, use the faithful action
\[
 \operatorname{End}^0(E_{\overline{\mathbb F}_p})
      \otimes_{\mathbb Q}\mathbb Q_\ell
       \hookrightarrow\operatorname{End}_{\mathbb Q_\ell}(V_\ell E)
\]
on the two-dimensional Tate module. The usual integral
injectivity follows from the fact that an endomorphism killing
all $\ell^n$-torsion factors through $[\ell^n]$; passing to the
$\ell$-adic completion and then inverting $\ell$ gives the displayed
map. The centralizer of
$K\otimes\mathbb Q_\ell$ there has dimension two: if the quadratic
algebra is a field the module has rank one over it, and if it is
split its two faithful eigenspaces each have dimension one.
Thus the rational centralizer has dimension at most two and
already contains the two-dimensional algebra $K$.

The $p$-power Frobenius endomorphism commutes with the specialized
CM endomorphisms, because they are defined over $\mathbb F_p$.
It therefore belongs to $K$. Its characteristic polynomial is
$X^2-a_{\mathfrak p}X+p$, so it is an algebraic integer of norm
$p$ and trace $a_{\mathfrak p}$. For the opposite sign replace
Frobenius by its negative, which still belongs to $\mathcal O_K$
and has norm $p$. Writing an algebraic integer in a quadratic
field in the usual integral basis gives
\[
 \pi=(a+b\sqrt{\operatorname{disc}K})/2,\qquad
 a\equiv b\,\operatorname{disc}K\pmod2,\qquad
 4p=a^2+Db^2.
                                                               \tag{NT-E10c}
\]
The integer $b$ is nonzero: otherwise $\pi$ would be a rational
algebraic integer and $p=\pi^2$ would not be prime.
This proves (NT-E10a). The argument uses only the Frobenius
characteristic polynomial, good-reduction specialization and
faithful Tate-module actions; no CM reduction-density theorem
or GRH is required.

\paragraph{Complete analytic inequality at the boundary.}
For a rational prime $p$, $2\sqrt p$ is not integral, so the Hasse
bound is strict and $B>0$. The identity
\[
 \delta\sqrt p=\frac{B}{2+a/\sqrt p}                    \tag{NT-E10d}
\]
holds, with a positive denominator smaller than $4$. Hence
$\delta\sqrt p<c$ implies $B<4c$ strictly.
If $p>c/2$, the same inequality forces $a>0$.
Conversely, suppose $a>0$, $p>c/2$, and $B<4c$.
Both $a$ and $2\sqrt p-c/\sqrt p$ are then positive.
Squaring the desired strict inequality shows the exact equivalence
\[
 \delta\sqrt p<c
 \quad\Longleftrightarrow\quad
 B<4c-\frac{c^2}{p}
 \quad\Longleftrightarrow\quad
 p>\frac{c^2}{4c-B}.                                   \tag{NT-E10e}
\]
There are only finitely many possibilities $B=Db^2<4c$.
If this set is nonempty, let $B_{\max}$ be its largest element and
take
\[
 p>\max\left\{\frac c2,\frac{c^2}{4c-B_{\max}}\right\}.
                                                               \tag{NT-E10f}
\]
Then the final inequality in (NT-E10e) holds simultaneously for
all those possibilities. If the set is empty, no split prime
qualifies. This proves (NT-E10b), including exclusion of equality.

For a good prime unramified and nonsplit in $F/k$, the two CM
eigenspaces are interchanged, so $a_{\mathfrak p}=0$, as proved in
NT-E08. Then $\delta\sqrt p=2p$, which fails the endpoint
inequality for $p>c/2$. The remaining ramified primes form a
fixed finite set. If $4c\le D$, (NT-E10a) gives $B\ge D\ge4c$
at every split prime, establishing the stated finiteness and
explicit restriction on the possible unramified exceptions.

\paragraph{What the endpoint gate says.}
The remaining possible rational primes belong to the finite
collection of quadratic norm families
\[
 p=\frac{a^2+Db^2}{4},\qquad
 0<|b|<2\sqrt{c/D},\qquad
 a>0,\qquad a\equiv b\,\operatorname{disc}K\pmod2.        \tag{NT-E10g}
\]
An actual prime of $k$ must still have the indicated
\emph{curve-specific Frobenius}: arbitrary prime values of these
forms are not asserted to have the desired trace for this
particular $E/k$. Thus the reduction neither discards
congruence/sign conditions nor proves infinitude of quadratic
polynomial prime values. It locates exactly the boundary left
open by NT-E09.

Geometrically, the endomorphism field forces the Frobenius norm
equation into the integral lattice $\mathcal O_K$ of the CM
torus. At this endpoint, the transverse lattice coordinate $b$
is bounded, rather than growing with the norm. Work/run:
\texttt{GRH\_NIGHT/number\_theory},
\texttt{cm\_endpoint\_gate.tex}, NT-E10; complete proof above,
using the CM eigenspace and good-reduction framework of NT-E06--08.

\subsection*{NT-A06--A07. Gaussian prime angles at almost every fixed center}

Write \(\mathbb T_G=\mathbf R/(\pi/2)\mathbf Z\), with probability
measure \(d\mu_G(\theta)=2\,d\theta/\pi\). For a nonzero ideal
\(\mathfrak a=(\alpha)\) of \(\mathbf Z[i]\), its angle is
\(\theta_{\mathfrak a}=\arg\alpha\pmod{\pi/2}\). Put
\[
 \Xi_m((\alpha))=(\alpha/\overline\alpha)^{2m},
 \qquad m\in\mathbf Z .
\]
These are well-defined conductor-one unitary Hecke characters.
For \(m\ne0\) their finite-part functions are entire, with completion
\[
 \pi^{-(s+2|m|)}\Gamma(s+2|m|)L(s,\Xi_m)
 =
 \pi^{-(1-s+2|m|)}\Gamma(1-s+2|m|)L(1-s,\Xi_m).
\]
The precise analytic assumption in this subsection is RH for
\(L(s,\Xi_m)\) for every positive integer \(m\); negative \(m\)
give the same functions. Conditional input \(M\) supplies this family.
The mean prime count below uses only the unconditional prime ideal
theorem for \(\mathbf Q(i)\). Input \(A\) alone does not supply the
nonzero angular infinity types.

\paragraph{Classical variance input.}
Rudnick--Waxman, \emph{Angles of Gaussian primes},
\texttt{arXiv:1705.07498v2}, Theorem~3.2, author p.~12,
proved in \S4.4, p.~20, gives the following. Fix nonzero nonnegative
functions \(f\in C_c^\infty((-1/2,1/2))\), even, and
\(\Phi\in C_c^\infty((1/2,1))\). Define
\[
 F_K(\theta)=\sum_{j\in\mathbf Z}
 f\left(\frac{K}{\pi/2}(\theta-j\pi/2)\right),\qquad
 P_{K,X}(\theta)=
 \sum_{\mathfrak p}\Phi(N\mathfrak p/X)\log N\mathfrak p\,
                         F_K(\theta_{\mathfrak p}-\theta).
\]
In the range used here, \(3\le K\le X\le K^2\), that theorem gives
\[
 \operatorname{Var}_{\mu_G}(P_{K,X})
       \ll_{f,\Phi}\frac{X}{K}(\log K)^2.             \tag{A06.1}
\]
Lemma~3.1, author pp.~11--12, gives
\[
 \mathbb E_{\mu_G}P_{K,X}
       \sim c_{f,\Phi}\frac XK,\qquad
 c_{f,\Phi}=\left(\int_{\mathbf R}f\right)
            \left(\int_0^\infty\Phi\right)>0.        \tag{A06.2}
\]
The asymptotic is uniform in \(K\): its entire \(K\)-dependence is
the exact factor \(1/K\) in the integral of \(F_K\).
The nonzero Fourier modes in \S4 use precisely the functions
\(L(s,\Xi_m)\) above; no zero-correlation or simplicity hypothesis
is part of the variance upper bound. The full asymptotic variance
conjecture of that paper is not used.

\paragraph{NT-A06: an explicit exceptional-measure formulation.}
Let \(J_K(\theta)\) be the open centered arc of length \(\pi/(2K)\).
For \(X\) sufficiently large and \(3\le K\le X\le K^2\),
let \(E_{K,X}\) be the set of centers whose arc \(J_K(\theta)\)
contains no prime ideal angle of norm in \((X/2,X)\), and let
\(E^{(1)}_{K,X}\) impose the same absence for absolute-degree-one
prime ideals. Then
\[
 \mu_G(E_{K,X})\ll\frac{K(\log K)^2}{X},\qquad
 \mu_G(E^{(1)}_{K,X})\ll\frac{K(\log K)^2}{X}+\frac1K .
                                                               \tag{A06.3}
\]
The implied constants are fixed once the two auxiliary functions
are fixed; no numerical constant is claimed.

\begin{proof}
For all sufficiently large \(X\), (A06.2) gives
\(\mathbb E P_{K,X}\ge c_{f,\Phi}X/(2K)>0\), uniformly in \(K\).
If \(E_{K,X}\) occurs, every summand in \(P_{K,X}\)
vanishes by its support. Chebyshev's inequality therefore gives
\[
 \mu_G\{P_{K,X}=0\}
 \le \frac{\operatorname{Var}(P_{K,X})}
           {(\mathbb E P_{K,X})^2}
 \ll \frac{K(\log K)^2}{X}.
 \]
The only prime ideals of absolute degree greater than one in
\(\mathbf Q(i)\) are those above inert rational primes, all with
angle zero. The set of centers whose arc \(J_K(\theta)\) contains
zero has measure \(1/K\). Outside that set any prime ideal found
by the first assertion has degree one. This proves the second
assertion, including the endpoint conventions (endpoints form a
measure-zero set).
\end{proof}

\paragraph{NT-A07: almost every center, at every sufficiently fine scale.}
Fix \(\epsilon>0\). Under the same angular RH family, for
\(\mu_G\)-almost every fixed \(\theta\), there is
\(K_0(\theta,\epsilon)\) such that for every real \(K\ge K_0\),
\(J_K(\theta)\) contains the angle of a degree-one Gaussian prime
ideal satisfying
\[
 N\mathfrak p\le C_\epsilon
              K(\log K)^3(\log\log K)^{1+\epsilon}.       \tag{A07.1}
\]
The constant \(C_\epsilon\) is independent of the center.
For almost every center this produces infinitely many such prime
ideals. The threshold may depend on the center. This is a
complete deduction from the cited variance bound, not a claim
of priority over other almost-sure formulations.

\begin{proof}
For integers \(j\) sufficiently large put
\[
 K_j=2^j,\qquad
 X_j=K_j(\log K_j)^3(\log\log K_j)^{1+\epsilon}.
\]
Then \(K_j\le X_j\le K_j^2\). Let \(E_j\) be the event that
no degree-one prime ideal of norm in \((X_j/2,X_j)\) has angle
in \(J_{K_j}(\theta)\). Formula (A06.3) gives
\[
 \mu_G(E_j)\ll_\epsilon
      \frac1{j(\log j)^{1+\epsilon}}+2^{-j}.
\]
This series converges. The first Borel--Cantelli lemma, which
requires no independence, says that almost every center belongs
to only finitely many \(E_j\). Consequently every sufficiently
large dyadic scale has such a prime.

For arbitrary sufficiently large real \(K\), choose \(j\) with
\(K\le 2^j<2K\). The centered arc \(J_{K_j}(\theta)\) is contained
in \(J_K(\theta)\). For all sufficiently large \(K\),
\(\log K_j\le2\log K\) and
\(\log\log K_j\le2\log\log K\); hence
\[
 X_j\le 2^{5+\epsilon}
           K(\log K)^3(\log\log K)^{1+\epsilon}.
\]
This proves (A07.1), for example with \(C_\epsilon=2^{5+\epsilon}\)
after increasing the threshold. This constant comes only from
the dyadic comparison, not from an explicit numerical variance
estimate. The selected dyadic primes have norm at least \(X_j/2\),
which tends to infinity, so infinitely many distinct prime ideals
occur.
\end{proof}

\paragraph{Quantifier boundary and geometric placement.}
Almost every fixed center cannot be replaced by a bound uniform
over all centers at this norm scale. If a nonreal Gaussian
integer \(a+ib\) has \(a,b>0\) and norm at most \(X\), then
\(\arg(a+ib)\ge\arctan(1/\sqrt X)\). Hence the open arc
\((0,\arctan(1/\sqrt X))\) contains no degree-one prime angle
of norm at most \(X\). A Gaussian prime may be represented in
the closed first quadrant after multiplication by a unit;
the degree-one primes on the axes do not give angles in that
open arc. This elementary gap explains why the exceptional set
and its order of quantifiers matter.

In the Grothendieck/Frobenius dictionary the prime ideals are
closed points of \(\operatorname{Spec}\mathbf Z[i]\), and the
\(\Xi_m\) are the integer weights of its compact CM torus.
The variance averages over the center of a torus neighborhood;
the first Borel--Cantelli step turns that averaged bound into
a statement about almost every fixed neighborhood system.
No independent random primes or independent zeros are assumed.
Run location: \texttt{GRH\_NIGHT/number\_theory};
the exact source is \texttt{S-NT-RW}, and the full new deductions
are (A06.3) and NT-A07 above.


\subsection*{NT-A08--A11. Prescribed ray class and Hecke angle, with uniform parameters}

Let \(K\) have degree \(n\) and absolute discriminant \(D\).
For a primitive unitary Hecke character \(\lambda\), use the standard
finite-part completion, with the complex factors split by duplication:
\[
 \Lambda_K(s,\lambda)=
 q_\lambda^{s/2}\prod_{j=1}^{n}\Gamma_{\mathbf R}(s+\kappa_{\lambda,j})
 L_K(s,\lambda),\qquad
 q_\lambda=D\,N\mathfrak f_\lambda,\qquad
 \Gamma_{\mathbf R}(s)=\pi^{-s/2}\Gamma(s/2).
                                                        \tag{A08.1}
\]
An irrelevant constant of absolute value one in the functional equation
is allowed. At a real place the parameter is
\(\epsilon+i\tau\), \(\epsilon\in\{0,1\}\); at a complex place the
two parameters are \(|k|/2+i\tau\) and \(|k|/2+i\tau+1\).
Put
\[
 {\cal C}(\lambda)=q_\lambda
                    \prod_{j=1}^{n}(3+|\kappa_{\lambda,j}|).
\]
These are the local Hecke conventions of Ramakrishnan--Valenza
Theorem~7-19 and its local-factor discussion, also checked in
MT-M28; Rajan, \S1, pp.~281--283, records the conductor and
infinity-type behavior under powers.

\paragraph{The exact classical prime-sum input.}
Iwaniec--Kowalski, \emph{Analytic Number Theory}, Theorem~5.15,
printed p.~114, equation~(5.56), states, under RH and the
Ramanujan bound for a function of degree \(d\), that
\[
 \psi_f(x)=r_fx+
 O\bigl(\sqrt x\log(2x)
             [\log{\cal C}(f)+d\log(2x)]\bigr),\qquad x\ge2,
                                                        \tag{A08.2}
\]
with an absolute constant. Here \(\psi_f\) uses the coefficients
of \(-L_f'/L_f\), \(r_f\) is the pole order at \(1\), and
\({\cal C}(f)\) is the analytic conductor at zero, defined in
\S5.1, (5.6)--(5.7), p.~95.
Its preceding proof sentence uses the truncated explicit formula
(5.53), pp.~112--113, and the local zero count (5.26).
This is a cited classical theorem; its preceding exercise and
full analytic development are not reproved here.

The book's definition in \S5.1, pp.~93--95, requires its infinity
parameters to be real or occur in conjugate pairs. An arbitrary
individual unitary Hecke function need not meet that convention.
The following complete reduction removes that restriction without
enlarging the order of the error term.

\paragraph{NT-A08: a uniform complex Hecke prime-sum bound.}
Suppose \(\lambda\) is not a pure unitary norm character and
RH holds for \(L_K(s,\lambda)\). Then, with an absolute constant,
\[
 \left|\sum_{N\mathfrak a\le x}
       \Lambda_K(\mathfrak a)\lambda(\mathfrak a)\right|
 \ll \sqrt x\log(2x)
          [\log{\cal C}(\lambda)+n\log(2x)].             \tag{A08.3}
\]
The value at an ideal not coprime to the conductor is zero.
Consequently, for any integral ideal \(\mathfrak f\) divisible by
\(\mathfrak f_\lambda\),
\[
 \sum_{\substack{N\mathfrak p\le x\\\mathfrak p\nmid\mathfrak f}}
            \lambda(\mathfrak p)
 \ll \sqrt x[\log{\cal C}(\lambda)+n\log(2x)]
                       +\log N\mathfrak f.             \tag{A08.4}
\]
For \(\lambda=1\), the same bound applies after subtracting
\(\operatorname{Li}_2(x)\) from the left side, using RH for
\(\zeta_K\). This principal case is handled separately.

\begin{proof}
Write
\[
 -\frac{L'}L(s,\lambda)=\sum_{r\ge1}b_r r^{-s}.
 \]
Grouping the ideal Euler factors by rational prime gives
\(|b_r|\le n\Lambda(r)\). In fact at \(r=p^a\) the coefficient
is a sum of terms \(f_{\mathfrak p}\log p\) times unit complex
numbers, over residue degrees \(f_{\mathfrak p}\mid a\), and
\(\sum_{\mathfrak p\mid p}f_{\mathfrak p}\le n\).
The local roots of the resulting degree-\(n\) rational-prime
Euler product have absolute value at most one.

For real \(t\) form
\[
 F_t(s)=L_K(s+it,\lambda)L_K(s-it,\overline\lambda).
 \]
Because \(\lambda\) is not pure norm, both factors are entire;
there are no shifted poles. Its degree is \(2n\), its conductor
is \(q_\lambda^2\), and its infinity parameters are
\(\kappa_{\lambda,j}+it\) and their conjugates. Its functional
equation has root number one. These statements include the
bad-prime Euler factors. The Hecke continuation, finite order
and boundary nonvanishing supply the remaining axioms of
Iwaniec--Kowalski. RH for the two translated factors follows
from RH for \(\lambda\) and complex conjugation.
For \(|t|\le2\),
\[
 \log{\cal C}(F_t)\le
           2\log{\cal C}(\lambda)+2n\log2 .
 \]
Its logarithmic derivative has coefficient
\(2\operatorname{Re}(b_r r^{-it})\).
Thus (A08.2), with pole order zero, bounds
\[
 R_t(y):=\sum_{r\le y}\operatorname{Re}(b_r r^{-it}),
 \qquad
 |R_t(y)|\le C\sqrt y\log(2y)
           [\log{\cal C}(\lambda)+n\log(2y)]              \tag{A08.5}
\]
uniformly in \(|t|\le2\).

Here are both steps recovering the imaginary part. Fix \(x\ge4\),
put \(a=\sqrt x\) and \(t=\pi/(2\log x)\), and hold this \(t\)
fixed while \(u\) ranges over \([a,x]\).
Denote the right side of (A08.5), with \(y=x\), by \(B_x\).
Write \(b_r=u_r+iv_r\). First partial summation gives
\[
 \begin{split}
 \sum_{a<r\le y}u_r\cos(t\log r)
  ={}&[R_0(y)-R_0(a)]\cos(t\log y)\\
    &+\int_a^y[R_0(u)-R_0(a)]
                       \frac{t\sin(t\log u)}u\,du .
 \end{split}
\]
The total variation of the cosine on this interval is at most
one, so this sum is \(O(B_x)\), uniformly for \(a\le y\le x\).
Since
\(\operatorname{Re}(b_r r^{-it})
   =u_r\cos(t\log r)+v_r\sin(t\log r)\),
it follows that
\[
 V(y):=\sum_{a<r\le y}v_r\sin(t\log r)=O(B_x).
 \]
Throughout \([a,x]\) the sine is at least \(1/\sqrt2\).
A second partial summation gives
\[
 \sum_{a<r\le x}v_r
   =\frac{V(x)}{\sin(t\log x)}
    +\int_a^x V(u)
                \frac{t\cos(t\log u)}
                     {u\sin^2(t\log u)}\,du
   =O(B_x),
 \]
because \(1/\sin(t\log u)\) has total variation at most
\(\sqrt2-1\). The initial segment has absolute sum at most
\(n\psi(\sqrt x)\ll n\sqrt x\), by the elementary Chebyshev
bound. It is absorbed in \(B_x\), since
\(\log{\cal C}(\lambda)\ge n\log3\).
The real part was already bounded by \(R_0(x)\).
This proves (A08.3); \(2\le x<4\) follows directly from the
coefficient bound.

Removing higher prime powers costs \(O(n\sqrt x\log(2x))\),
which is absorbed in (A08.3). Indeed, their underlying rational
primes are at most \(\sqrt x\), and the total weight above
each is at most \(n\log x\).
Abel summation removes the prime logarithm: if the weighted
error is \(O(\sqrt u\log(2u)E(u))\), with
\(E(u)=\log{\cal C}(\lambda)+n\log(2u)\), the endpoint error
is \(O(\sqrt xE(x))\), and its integral is bounded by
\[
 E(x)\int_2^x\frac{\log(2u)}
                     {\sqrt u(\log u)^2}\,du
                       \ll\sqrt x E(x).
 \]
Omitting additionally the primes dividing \(\mathfrak f\)
costs at most \(\omega(\mathfrak f)\le
\log N\mathfrak f/\log2\).
This proves (A08.4). For the principal character use the
Dedekind prime-count estimate NT-C02, with the same deletion.
In particular the principal factor is never put into \(F_t\).
\end{proof}

\paragraph{NT-A09: joint class and arc discrepancy.}
Let \(\mathfrak m=\mathfrak m_0\mathfrak m_\infty\) be a ray
modulus, \(G=\operatorname{Cl}_{\mathfrak m}(K)\) and \(h=|G|\).
Let \(\eta\) be unitary and not of the form finite order times
a unitary norm power. Put
\[
 \begin{aligned}
 \mathfrak f&=\operatorname{lcm}(\mathfrak m_0,\mathfrak f_\eta),
 & B_\eta&=4+\max_j|\kappa_{\eta,j}|,\\
 Q&=D\,N\mathfrak f\,B_\eta^n,
 & {\cal L}(x)&=\log Q+2n\log(3x).
 \end{aligned}                                                     \tag{A09.1}
\]
For \(a\in G\) and an arc \(J\subset\mathbf R/\mathbf Z\), let
\(N_{a,J}(x)\) count primes \(\mathfrak p\nmid\mathfrak f\)
of norm at most \(x\) with ray class \(a\) and
\(\eta(\mathfrak p)\in e^{2\pi iJ}\).
Assume RH for the primitive Hecke functions attached to
\[
             \chi\eta^m,\qquad
     \chi\in\widehat G,\quad m=0,1,2,\ldots .             \tag{A09.2}
\]
Conditional input \(M\) suffices. Then
\[
 N_{a,J}(x)=\frac{|J|}{h}\operatorname{Li}_2(x)
           +O\bigl(\sqrt x\log(2x){\cal L}(x)\bigr),
 \qquad x\ge2,                                         \tag{A09.3}
\]
with an absolute constant, uniformly in \(K,\eta,\mathfrak m,a,J\).
The same bound holds if only absolute-degree-one primes are counted.
In particular there is no factor \(h\) in the error.

\begin{proof}
For \(\chi\in\widehat G\), its finite conductor divides
\(\mathfrak m_0\); its infinity type has only real signs.
The finite conductor of \(\chi\eta^m\) divides \(\mathfrak f\).
At complex places its angular and imaginary parameters grow
by \(m\); at real places its parity remains in \(\{0,1\}\).
The conventions (A08.1) therefore give directly
\[
 {\cal C}(\chi\eta^m)\le Q(m+1)^n,\qquad m\ge0.           \tag{A09.4}
\]
For \(m>0\), none of these characters is pure norm. Otherwise,
if \(\chi\eta^m=|\cdot|^{it}\), the character
\(\eta|\cdot|^{-it/m}\) would have finite order: its \(m\)-th
power is the finite-order character \(\chi^{-1}\).
This contradicts the hypothesis on \(\eta\).
For \(m=0\), only \(\chi=1\) is principal.

Apply (A08.4), treating that principal character separately.
The additional \(\log N\mathfrak f\) is absorbed by
\(\sqrt x\log Q\). For each \(m\ge0\) we obtain
\[
 \sum_{\substack{N\mathfrak p\le x\\\mathfrak p\nmid\mathfrak f}}
       \chi(\mathfrak p)\eta(\mathfrak p)^m
 =\mathbf1_{\chi=1,m=0}\operatorname{Li}_2(x)
 +O\left(\sqrt x\{\log Q+n\log[(m+1)(2x)]\}\right).
                                                               \tag{A09.5}
\]
The ray-class indicator is
\[
 \mathbf1_{[\mathfrak p]=a}
       =\frac1h\sum_{\chi\in\widehat G}
                        \overline{\chi(a)}\chi(\mathfrak p).
 \]
Its \(h\) terms cancel the prefactor \(1/h\) when the errors
in (A09.5) are summed in absolute value. Thus the number
\(N_a(x)\) in class \(a\) has main term \(\operatorname{Li}_2(x)/h\),
and every nonzero angular Fourier sum in that class has
the same error bound without \(h\).

Apply the Erd\H{o}s--Tur\'an inequality S-NT-ET to the finite
list of angles in this single class, with \(H=\lceil x\rceil\).
Its zero-frequency truncation error is
\(O(N_a(x)/(H+1))=O(n)\): at most \(n\) prime ideals lie over
each rational prime of size at most \(x\).
For \(1\le m\le H\),
\(\log Q+n\log[(m+1)(2x)]\ll{\cal L}(x)\).
Summing the \(1/m\) weights gives (A09.3), first for half-open
arcs. Singleton bounds follow by decreasing arcs to an endpoint;
changing either endpoint convention changes the count by at
most the same error. This also covers arcs crossing zero.
Finally there are at most \(n\pi(\sqrt x)=O(n\sqrt x)\)
prime ideals of absolute residue degree at least two and norm
at most \(x\). Their deletion is absorbed in the error.
\end{proof}

\paragraph{NT-A10: shrinking arcs with separate class-number dependence.}
Fix \(c>0\) and \(0<\alpha<1/2\), and let
\(u_0=\max\{2,c^{1/\alpha}\}\).
For a fixed center let the arc at norm \(u>u_0\) have length
\(c u^{-\alpha}\). Let \(N_{a,c,\alpha}(x)\) count precisely
the primes \(\mathfrak p\nmid\mathfrak f\) with
\(u_0<N\mathfrak p\le x\), ray class \(a\), and angle in
that arc at \(u=N\mathfrak p\).
For \(x>\max\{4,u_0\}\), this joint count satisfies
\[
 \begin{split}
 N_{a,c,\alpha}(x)
  ={}&\frac c h\int_{u_0}^x\frac{u^{-\alpha}}{\log u}\,du\\
   &+O_{\alpha,c}\left(
       \sqrt x\log(2x){\cal L}(x)
       +h^{-1/2}x^{3/4-\alpha/2}\sqrt{{\cal L}(x)}
                  \right).                            
 \end{split} \tag{A10.1}
\]
The constant is uniform in the center and in the parameters
of (A09.1). For varying fields and ray moduli, a sufficient
uniform asymptotic regime is
\[
 h\,{\cal L}(x)(\log x)^2=o(x^{1/2-\alpha}).              \tag{A10.2}
\]
For fixed data this holds automatically. The same assertions
hold after restricting to degree-one primes.

\begin{proof}
Set \(T(x)=\log(2x){\cal L}(x)\). Partition \((u_0,x]\)
geometrically with ratio \(1+\delta\), \(0<\delta\le1/2\).
Nested centered arcs bound the moving count on each interval
between its two fixed-endpoint arc counts. Apply (A09.3) to
each difference of prime-counting functions. The geometric
sum of endpoint errors is
\[
 O\bigl(\delta^{-1}\sqrt x\,T(x)\bigr),
 \]
since \(\sum \sqrt{u_j}\ll\delta^{-1}\sqrt x\), and
\(T(u_j)\le T(x)\).
The two main terms bracket the displayed integral, and their
difference from it is
\[
 O_\alpha\left(\frac c h\delta
           \int_{u_0}^x\frac{u^{-\alpha}}{\log u}\,du\right)
 \ll_{\alpha,c}\frac c h\delta
                   \frac{x^{1-\alpha}}{\log x}
 \]
for \(x\) above a threshold depending only on \(c,\alpha\).
This follows from the ratio \((1+\delta)^\alpha=1+O_\alpha(\delta)\);
the last incomplete interval obeys the same bounds.

Write \(A=\sqrt x\,T(x)\) and
\(B=(c/h)x^{1-\alpha}/\log x\).
Taking \(\delta=\min\{1/2,\sqrt{A/B}\}\) gives
\[
 A/\delta+B\delta\ll A+\sqrt{AB}.
 \]
Indeed when \(\sqrt{A/B}\le1/2\) both terms equal \(\sqrt{AB}\);
otherwise \(B<4A\). Since
\[
 \sqrt{AB}
  =\sqrt{c/h}\,x^{3/4-\alpha/2}
       \sqrt{\frac{\log(2x)}{\log x}{\cal L}(x)},
 \]
this proves (A10.1). Its main term is asymptotic to
\(c\,x^{1-\alpha}/[h(1-\alpha)\log x]\).
The ratios of the two error terms to this main term are,
up to fixed constants,
\[
 h{\cal L}(x)(\log x)^2x^{\alpha-1/2},
 \qquad
 \left[h{\cal L}(x)(\log x)^2x^{\alpha-1/2}\right]^{1/2}.
 \]
They tend to zero under (A10.2).
Using the degree-one version of (A09.3) in the identical
partition proves the last assertion.
\end{proof}

\paragraph{NT-A11: a least prime in a prescribed class and arc.}
Let \(0<\mu\le1\) be the length of a fixed arc. Put
\[
 d=\frac h\mu,\qquad
 V=\log(eQ)+n\log(e\,n\,d).
 \]
Under (A09.2), there is an absolute constant \(C>0\) such
that every ray class and every arc of length \(\mu\) contains
a degree-one prime \(\mathfrak p\nmid\mathfrak f\) with
\[
                       N\mathfrak p\le C d^2 V^6 .      \tag{A11.1}
\]
No numerical value of \(C\) is claimed.

\begin{proof}
Use (A09.3) with its degree-one version. Let \(A_0\) be an
absolute error constant there, and put \(x=C d^2 V^6\).
We have \(d\ge1\), \(V\ge n\), \(n\log d\le V\) and
\(n\log V\le2V\). The last inequality follows by writing
\(V=nv\), since \(v\ge\log(e\,n)\) and
\(\log V=\log n+\log v\le2v\).
Consequently, for \(C\ge4\),
\[
 \log x\le(\log C+14)V/n,\qquad
 \log(2x)\le(\log C+15)V/n,\qquad
 {\cal L}(x)\le(32+2\log C)V .
 \]
Also \(\operatorname{Li}_2(x)\ge x/(2\log x)\) for \(x\ge4\).
Thus a sufficient condition for the main term
\(\operatorname{Li}_2(x)/d\) to exceed the error is
\[
 \sqrt C>
 2A_0(\log C+14)(\log C+15)(32+2\log C).
 \]
An absolute \(C\) satisfying this exists, because \(\sqrt C\)
eventually dominates that cubic polynomial in \(\log C\).
The resulting count is positive, proving (A11.1).
\end{proof}

\paragraph{Geometric and run placement.}
The finite condition is a Frobenius fiber in the ray-class
cover of \(\operatorname{Spec}\mathcal O_K[1/\mathfrak f]\);
the angular condition lies in a compact rank-one Hecke torus.
Their joint character algebra is
\(\widehat G\times\mathbf Z\), and (A09.2) is exactly its
Hecke \(L\)-function family. For normalized algebraic Hecke
characters this is the tensor-character description of a CM
torus with a finite Artin factor. For a general unitary
infinity type no algebraic-motive realization is asserted.
The formulas keep the finite mass \(1/h\), the finite
conductor, the infinity parameters and the field degree
separate. They strengthen the separate class and angle
formulations NT-C06/A04--A05. Run:
\texttt{GRH\_NIGHT/number\_theory}; full new proofs are
NT-A08--A11 above, with classical input S-NT-IK and S-NT-ET.



\subsection*{NT-C09--C13. Smaller class-group generators and exact computational input}

Let \(K\) be a number field of degree \(n\ge2\), discriminant
\(D=|\operatorname{disc}K|\), and \(r_1\) real places. All class
groups below are ordinary ideal class groups. Put
\[
 G=\operatorname{Cl}(K),\qquad
 G(T)=\langle[\mathfrak p]:N\mathfrak p\le T\rangle .
\]
The analytic input used for the uniform bounds is
\[
 {\cal R}_K:\quad
 \text{RH for }\zeta_{\mathbf Q}
 \text{ and for every }L_K(s,\chi),\quad \chi\in\widehat G,
                                                               \tag{C09.1}
\]
including \(\chi=1\), whose function is \(\zeta_K\).
Equivalently, it suffices to assume RH for \(\zeta_{\mathbf Q}\)
and \(\zeta_{H_K}\), where \(H_K\) is the ordinary Hilbert class
field: the all-prime abelian factorization NT-H11 has no
factor poles in the open critical strip. Thus \(M\) supplies
(C09.1). The designated input \(A\) suffices if \(H_K/\mathbf Q\)
is abelian; abelianness of \(K/\mathbf Q\) alone does not give
that conclusion.

\paragraph{NT-C09 and NT-C10: the precise published bounds.}
Greni\'e--Molteni, \emph{Breaking the 4 barrier for the bound of
a generating set of the class group}, Math.\ Comp.\ 94 (2025),
3145--3163, Theorems~1--4, author pp.~2--3, give the following
conclusions under (C09.1):
\[
 \begin{array}{ll}
 G(B_n)=G,&
 B_n=\left(4-\dfrac1{2n}+\dfrac1{2n^2}\right)(\log D)^2,\\[5pt]
 G(B'_n)\supseteq G^2,&
 B'_n=\left(4-\dfrac2{3n}+\dfrac3{8n^2}\right)(\log D)^2 .
 \end{array}                                                    \tag{C09.2}
\]
Without the respective \(n^{-2}\) corrections, these conclusions
hold for every \(n\ge5\). For \(n=2,3,4\), the verified ranges
are \(\log D\le l_-\) or \(\log D\ge l_+\), with
\[
 \begin{array}{c|rrr}
 n&2&3&4\\ \hline
 l_-&63&353&833\\
 l_+&12184&18455&27911\\ \hline
 l'_-&65&379&993\\
 l'_+&7968&9832&13108
 \end{array} .                                                  \tag{C09.3}
\]
The primed rows apply to \(G^2\).
The excluded middle intervals are intervals where the sharper
bound is not asserted, rather than a list of actual counterexamples.
For odd \(|G|\), squaring is an automorphism and both primed
bounds generate all of \(G\).

The proofs are in \S4, author pp.~10--12, and \S5,
pp.~12--13. Their analytic inputs are visible in \S2,
pp.~3--5 (class-character and Dedekind zero sums), and
Proposition~1, pp.~8--9 (RH for \(\zeta_{\mathbf Q}\)).
The small-discriminant input is their 2018
\emph{Explicit bounds for generators of the class group},
Theorem~3.5, equation~(3.10), author p.~6, using the
same class-character and Dedekind inputs.
The threshold computations, Odlyzko bounds and certified zero
tables remain credited parts of the published proof; they were
not numerically rerun here.

The containment in (C09.2) is the precise conclusion proved
at the start of the 2025 paper's \S5. Its introductory definition
of \(T'_K\) uses the phrase of generating \(G^2\); arbitrary
small prime classes need not themselves belong to \(G^2\).
We use the unambiguous subgroup-containment formulation throughout.

\paragraph{NT-C11: a finite certificate for a quotient, with exclusions.}
Here a smaller RH input suffices. Fix a subgroup \(H\le G\).
Assume RH only for the nontrivial class characters \(\chi\)
that are trivial on \(H\). No assumption on \(\zeta_K\) or
\(\zeta_{\mathbf Q}\) is needed in this paragraph.

Let \(L>0\), and let \(F:\mathbf R\to\mathbf R_{\ge0}\) be even, continuous,
compactly supported in \([-L,L]\), with nonnegative Fourier
transform. Assume the regularity needed in the Weil--Poitou
formula, for instance that \(F\) is piecewise \(C^1\) and
\((F(0)-F(x))/x\) is of bounded variation on \((0,\infty)\).
Define
\[
 \begin{split}
 I(F)&=\int_0^\infty
                  \frac{F(0)-F(x)}{2\sinh(x/2)}\,dx,\\
 J(F)&=\int_0^\infty\frac{F(x)}{2\cosh(x/2)}\,dx,\\
 W_F(\mathfrak p)&=\log N\mathfrak p
       \sum_{m\ge1}\frac{F(m\log N\mathfrak p)}
                              {(N\mathfrak p)^{m/2}},\\
 R_K(F)&=F(0)\{\log D-(\gamma+\log8\pi)n\}
                                      +nI(F)-r_1J(F).
 \end{split}                                                    \tag{C11.1}
\]
Let \(E\) be any set of prime ideals to be omitted. If
\[
       2\sum_{\mathfrak p}
             (-1)^{\mathbf1_E(\mathfrak p)}W_F(\mathfrak p)
                          >R_K(F),                            \tag{C11.2}
\]
then the classes of primes \(\mathfrak p\notin E\) with
\(N\mathfrak p<e^L\), together with \(H\), generate \(G\).
Only finitely many terms in (C11.2) are nonzero. This is a
quotient formulation of Greni\'e--Molteni (2025),
Theorem~5, author pp.~13--14, with the support stated in
the logarithmic variable in which the test function is evaluated.

\begin{proof}
Suppose the stated classes and \(H\) generate a proper
subgroup \(U<G\). A nontrivial character of \(G/U\), inflated
to \(G\), gives a character \(\chi\ne1\) trivial on \(H\)
and on every selected prime. Write
\[
       \phi(s)=\int_{\mathbf R}F(x)e^{(s-1/2)x}\,dx.
\]
The Weil--Poitou identity, in the normalization of the
2025 source (3), author p.~3, gives
\[
 2\sum_{\mathfrak p}\log N\mathfrak p
       \sum_{m\ge1}
       \frac{\operatorname{Re}(\chi(\mathfrak p)^m)
                          F(m\log N\mathfrak p)}
                              {(N\mathfrak p)^{m/2}}
       =R_K(F)-\sum_{\rho_\chi}\phi(\rho_\chi).
                                                               \tag{C11.3}
\]
The supplied RH hypothesis and Fourier positivity make every
zero term nonnegative. A prime outside \(E\) contributing to
the sum has norm below \(e^L\), hence character value one.
Indeed continuity and compact support force \(F(L)=0\),
so the endpoint creates no term. At a prime in \(E\),
\(\operatorname{Re}(\chi(\mathfrak p)^m)\ge-1\).
Since \(F\ge0\), the left side of (C11.3) is at least
the left side of (C11.2), whereas its right side is at most
\(R_K(F)\). This contradicts the strict inequality and
proves the assertion. The zero sum can equivalently be
read as the symmetric-height limit in the classical formula;
its terms are nonnegative here, so the comparison is valid.
\end{proof}

There is a directly computable sufficient version. Let \(T>1\),
\(L=\log T\), \(F(x)=(1-|x|/L)_+\), and write
\[
 {\mathfrak G}=\sum_{j\ge0}\frac{(-1)^j}{(2j+1)^2}.
\]
Then it suffices that
\[
 \begin{split}
 2\sum_{(N\mathfrak p)^m<T}
   (-1)^{\mathbf1_E(\mathfrak p)}
   \frac{\log N\mathfrak p}{(N\mathfrak p)^{m/2}}
        \left(1-\frac{m\log N\mathfrak p}{\log T}\right)
 >{}&\log D-n(\gamma+\log8\pi)-\frac{\pi r_1}{2}\\
    &+\frac{\pi^2 n/2+4{\mathfrak G}r_1}{\log T}.
 \end{split}                                                    \tag{C11.4}
\]
For \(H=1\) and \(E=\varnothing\), this is exactly the
Belabas--Diaz y Diaz--Friedman criterion, Corollary~2.2,
author pp.~6--7; its rounded Theorem~1.1 is on p.~3.
Here \(H=1\) denotes the trivial subgroup.

\begin{proof}
The Fourier transform of the triangular function is
\(L(\sin(tL/2)/(tL/2))^2\), with the value at zero taken
by continuity. It is nonnegative and meets all the
regularity assumptions. Also
\[
 1-F(x)\le x/L\quad(x\ge0),\qquad
 \int_0^\infty\frac{dx}{2\cosh(x/2)}=\frac\pi2.
\]
Thus
\[
 I(F)\le\frac{\pi^2}{2L},\qquad
 J(F)\ge\frac\pi2-\frac{4{\mathfrak G}}L.
 \]
For completeness, the constants follow by expanding
\[
 \frac1{2\sinh(x/2)}
       =\sum_{j\ge0}e^{-(j+1/2)x},\qquad
 \frac1{2\cosh(x/2)}
       =\sum_{j\ge0}(-1)^j e^{-(j+1/2)x}
\]
and integrating \(x\) times the series. The first integration
is monotone; the second is justified by the integrable
absolute sum of the \(x\)-weighted integrals. Their values
are \(\pi^2/2\) and \(4{\mathfrak G}\), respectively.
These inequalities bound \(R_K(F)\) above by the right
side of (C11.4); now apply (C11.2).
\end{proof}

\paragraph{NT-C12: what can actually be enumerated in polynomial time.}
Suppose an integral basis of the maximal order \(\mathcal O_K\)
and its full integer multiplication table are supplied.
Under (C09.1), one can deterministically enumerate, in
polynomial time in this input's bit length, prime ideals
whose classes generate \(G\), all of norm at most \(B_n+2\).
With the additional promise that \(|G|\) is odd, \(B'_n+2\)
may be used instead. The harmless additive two permits
certified rational rounding of the displayed logarithmic
bound; no unproved separation bound from an integer is needed.
The output is a list of integral ideal lattices.
This does not compute a relation lattice, the class number,
units, or a presentation of the class group.

\begin{proof}
First, the discriminant is computable by exact linear algebra
from the trace form. If all multiplication constants have
absolute value at most \(2^b\), the matrix of multiplication
by a basis element has that entry bound. Therefore each
trace-form entry has absolute value at most \(n^2 2^{2b}\).
Hadamard's inequality gives
\[
       \log_2 D\le 2bn+\frac{5n}{2}\log_2 n.                  \tag{C12.1}
\]
Consequently any integer \(T=O(1+(\log D)^2)\) is polynomial
in the supplied input length.

An integer \(T\) with \(B_n\le T\le B_n+2\) is obtainable
in polynomial time. Compute a rational upper approximation
to \(B_n\) with error less than \(1/2\), and take its ceiling.
For example, write \(D=2^m u\), \(1\le u<2\), and evaluate
\[
 \log u=2\sum_{j\ge0}
       \frac{1}{2j+1}\left(\frac{u-1}{u+1}\right)^{2j+1},
\]
as well as the same series for \(\log2\).
The ratios have absolute value at most \(1/3\); geometric
tail bounds give certified accuracy \(O((1+\log D)^{-1})\)
for \(\log D\), dividing the error allowance for \(\log2\)
by \(m+1\) before forming \(m\log2+\log u\).
This uses polynomially many rational bit operations. Squaring
and multiplying by the fixed rational coefficient gives
the required approximation. The same procedure applies
to \(B'_n\).

We describe an unconditional procedure to list all prime
ideals with norm at most this integer \(T\).
List rational primes \(p\le T\), even by trial division.
For each such \(p\), form the \(n\)-dimensional algebra
\[
                   A=\mathcal O_K/p\mathcal O_K
\]
over \(\mathbf F_p\), using the given table. Its Frobenius
\(F:a\mapsto a^p\) is \(\mathbf F_p\)-linear and is
computed by repeated squaring. Choose \(e\) with \(p^e\ge n\).
Then
\[
                   \operatorname{Nil}(A)=\ker F^e.           \tag{C12.2}
\]
Indeed a nilpotent in an \(n\)-dimensional algebra has
nilpotence index at most \(n\), whereas Frobenius is
injective on every finite residue field. Form the reduced
quotient \(R=A/\operatorname{Nil}(A)\). It is a product of
finite fields, and
\[
          B=\{b\in R:b^p=b\}\simeq\mathbf F_p^{\,g},
                                                               \tag{C12.3}
\]
where \(g\le n\) is the number of factors.

A basis of \(B\) is obtained by solving linear equations.
Use its elements successively to split \(R\) into the
simultaneous eigenspaces of their multiplication maps.
Every eigenvalue belongs to \(\mathbf F_p\), so it is enough
to scan \(a=0,\ldots,p-1\) and compute kernels of
\(m_b-aI\) on each current component. These maps are
diagonalizable because \(b^p=b\) and \(X^p-X\) has distinct
roots. The common eigenspaces are exactly the field factors:
a basis of \(\mathbf F_p^g\) separates every pair of
coordinates. All kernels are ideals, and the nonzero
resulting components give the maximal ideals of \(R\)
as kernels of the corresponding projections. Pulling
these back to \(A\) lists all maximal ideals \(W\subset A\).

If \(f=\dim_{\mathbf F_p}(A/W)\), the associated prime
ideal of \(\mathcal O_K\) has norm \(p^f\); retain it
exactly when \(p^f\le T\). It is explicitly the inverse
image of the subspace \(W\). To output an integer basis,
put a basis of \(W\) in reduced row echelon form, lift its
entries to \(\{0,\ldots,p-1\}\), and add \(p\) times each
coordinate vector at a nonpivot position. These \(n\)
vectors generate precisely the inverse-image lattice:
subtracting the lifted rows using the integral pivot
coordinates leaves zero pivot coordinates and multiples
of \(p\) elsewhere. Their determinant has absolute value
\(p^f\), as also follows from the quotient size.

All algebra dimensions are at most \(n\). The number of
scanned scalars is polynomial in \(p,n\), and \(p\le T\)
is polynomial in the input length by (C12.1). All modular
entries use \(O(\log T)\) bits. There are at most \(n\)
prime ideals above each \(p\), so the number of output
ideals is at most \(nT\), with polynomial total encoding
length. This proves the claimed deterministic bit
complexity. The generation conclusion follows from
(C09.2) because \(T\ge B_n\), or \(T\ge B'_n\) with
odd class number.
\end{proof}

The supplied maximal-order basis is essential to this
formulation. Starting only with a defining polynomial
does not supply it for free. The proof makes no complexity
claim about finding that basis, identifying relations,
or proving an odd-class-number promise.

\paragraph{NT-C13: finite Frobenius detection of ordinary class-field quotients.}
Under (C09.1), let \(E/K\) be abelian, unramified at every
finite place and split at every real place. If all primes
of norm at most \(B_n\) split completely in \(E\), then
\(E=K\). If the smaller bound \(B'_n\) is used instead,
then \(\operatorname{Gal}(E/K)\) has exponent at most two.
The correction-free bounds in (C09.3) have the same
respective conclusions.

\begin{proof}
Ordinary global class field theory gives a surjection
\[
       G\twoheadrightarrow\operatorname{Gal}(E/K)
\]
sending each prime class to its Frobenius. Complete
splitting kills the indicated small-prime subgroup.
At \(B_n\) that subgroup is \(G\), so the quotient is
trivial. At \(B'_n\) it contains \(G^2\); consequently
the quotient is a quotient of \(G/G^2\) and has exponent
at most two. The sharper-range versions follow identically.
\end{proof}

In geometric language, \(G=\operatorname{Pic}
(\operatorname{Spec}\mathcal O_K)\), and these tests concern
abelian finite \'etale arithmetic covers with the ordinary
real-place splitting condition. The finite prime-sum
certificate proves generation of a prescribed quotient
by its closed-point Frobenius elements. The enumeration
proof exposes the precise integral-algebra input needed
to make these points concrete. Run:
\texttt{GRH\_NIGHT/number\_theory}; primary inputs
\texttt{S-NT-GM2025-CG}, \texttt{S-NT-BDDF},
and \texttt{S-NT-GM2018-CG}.



\subsection*{NT-A12--A14. A nonabelian Frobenius class and a Hecke angle}

Let \(L/K\) be a finite Galois extension with group \(G\), and let
\(C\subset G\) be one conjugacy class. Choose \(g\in C\) and put
\[
 \begin{aligned}
 H&=\langle g\rangle,\quad h=|H|,\quad F=L^H,\quad d=[F:K],\\
 \delta&=\frac{|C|}{|G|},\quad n_F=[F:\mathbf Q],\quad
                              n_L=[L:\mathbf Q]=h n_F .
 \end{aligned}
                                                               \tag{A12.1}
\]
Let \(\eta\) be a unitary Hecke character of \(K\) which is not a
finite-order character times a unitary norm power. Write
\(\nu=\eta\circ N_{F/K}\), and use the primitive character attached
to every product below. Let \(S\) be a finite set of prime ideals of
\(K\) containing the primes ramified in \(L/K\) and the support of
\(\mathfrak f_\eta\). Extra prescribed omissions may be included.
Set
\[
 \begin{aligned}
 \mathfrak s_F&=\prod_{\mathfrak q\mid\mathfrak p,\,
                                      \mathfrak p\in S}\mathfrak q,\\
 B_\nu&=4+\max_j|\kappa_{\nu,j}|,\\
 Q_C&=D_L
       \bigl(N_F\mathfrak f_\nu\,N_F\mathfrak s_F\bigr)^h
       B_\nu^{n_L},\\
 {\cal L}_C(x)&=\log Q_C+2n_L\log(3x).
 \end{aligned}                                                 \tag{A12.2}
\]
Here \(D_L\) is the absolute discriminant; the gamma parameters
\(\kappa_{\nu,j}\) have precisely convention (A08.1).
All conductors and norms in (A12.2) are over \(F\).

The exact analytic input is RH for
\[
 L_F(s,\chi\nu^m),\qquad
       \chi\in\widehat H,\quad m=0,1,2,\ldots,                 \tag{A12.3}
\]
where class field theory regards \(\chi\) as a finite-order Hecke
character for the cyclic extension \(L/F\). In particular the
\((\chi,m)=(1,0)\) member is \(\zeta_F\).
Conditional input \(M\) supplies this family. No holomorphy or RH
for an irreducible nonabelian Artin function is assumed.

\paragraph{NT-A12: uniform joint discrepancy.}
For an arc \(J\subset\mathbf R/\mathbf Z\), count prime ideals
\(\mathfrak p\notin S\) with \(N_K\mathfrak p\le x\),
Frobenius class \(C\), and \(\eta(\mathfrak p)\in e^{2\pi iJ}\).
Then
\[
 N_{C,J}(x)
   =\delta |J|\operatorname{Li}_2(x)
      +O\bigl(\delta\sqrt x\log(2x){\cal L}_C(x)\bigr),
       \qquad x\ge2.                                      \tag{A12.4}
\]
The constant is absolute, uniformly in all the displayed data,
the arc and its endpoint conventions. The same bound holds when
only absolute-degree-one primes of \(K\) are counted.

\begin{proof}
We first check the possible norm degeneracy. The diagonal idele
map \(j:C_K\longrightarrow C_F\) satisfies
\[
 N_{F/K}(j(c))=c^d,\qquad |j(c)|_F=|c|_K^d.
\]
If \(\nu=\theta|\cdot|_F^{it}\) with \(\theta\) of finite order,
composition with \(j\) gives
\((\eta|\cdot|_K^{-it})^d=\theta\circ j\).
If the latter has order dividing \(r\), the former character
raised to \(dr\) is trivial. This contradicts the hypothesis on
\(\eta\). Thus \(\nu\) also is not finite order times norm.
For \(m>0\), an equality
\(\chi\nu^m=|\cdot|_F^{it}\) would make
\((\nu|\cdot|_F^{-it/m})^m=\chi^{-1}\), giving the same
contradiction. Hence every \(m>0\) member of (A12.3) is entire,
with no pure-norm pole. For \(m=0\) only \(\chi=1\) is principal.
The principal member is treated separately by NT-A08.

Next we retain the sum of conductors instead of replacing each
one by its largest value. Locally, the conductor of a product
divides the product of conductors, and taking a power cannot
increase the conductor. Therefore
\[
 \mathfrak f_{\chi\nu^m}\mid\mathfrak f_\chi\mathfrak f_\nu,
 \qquad
 {\cal C}(\chi\nu^m)
   \le D_F N_F\mathfrak f_\chi N_F\mathfrak f_\nu
                       [B_\nu(m+1)]^{n_F}.               \tag{A12.5}
\]
The gamma-factor inequality follows directly from (A08.1):
a finite-order character changes only the real parities, while
the imaginary and complex angular parameters of \(\nu^m\)
grow linearly in \(m\). It also holds at \(m=0\).
The conductor--discriminant formula for the cyclic extension
\(L/F\) gives the exact identity
\[
       D_L=D_F^h\prod_{\chi\in\widehat H}
                                      N_F\mathfrak f_\chi .
                                                               \tag{A12.6}
\]
This includes all ramified local factors. The precise classical
locator is Neukirch, Chapter VII, (11.9), printed p.~534,
with the supplied conductor/induction proof on pp.~533--534
and the one-dimensional conductor identification (11.10),
p.~535; see S-NT-NEUKIRCH. Summing the logarithms in (A12.5)
thus yields
\[
 \sum_{\chi\in\widehat H}\log{\cal C}(\chi\nu^m)
 \le \log D_L+h\log N_F\mathfrak f_\nu
          +n_L\log B_\nu+n_L\log(m+1).                 \tag{A12.7}
\]

Every character in (A12.3) is unramified outside the primes
dividing \(\mathfrak s_F\). Indeed a local norm sends units
to units, and the characters of \(L/F\) ramify only over the
ramification locus already in \(S\). The prime version of
NT-A08, deleting the additional primes in \(\mathfrak s_F\),
therefore gives
\[
 \begin{split}
 P_{\chi,m}(x)
  &:=
    \sum_{\substack{N_F\mathfrak q\le x\\
                      \mathfrak q\nmid\mathfrak s_F}}
                 \chi(\mathfrak q)\nu(\mathfrak q)^m\\
  &=\mathbf1_{\chi=1,m=0}\operatorname{Li}_2(x)
    +O\left(\sqrt x\{\log{\cal C}(\chi\nu^m)
                       +n_F\log(2x)\}
                   +\log N_F\mathfrak s_F\right).
 \end{split}                                                   \tag{A12.8}
\]
Its constant is absolute. The last error is simply the
number of additional omitted prime ideals, at most
\(\log N_F\mathfrak s_F/\log2\); a conductor exponent is
not required to count these omissions. Formula (A12.8)
also follows directly from the proof of (A08.4).

We now establish the induction identity used at primes:
\[
       \mathbf1_C
       =\delta\sum_{\chi\in\widehat H}
                         \overline{\chi(g)}\,
                            \operatorname{Ind}_H^G\chi .
                                                               \tag{A12.9}
\]
Orthogonality on \(H\) says
\(\sum_\chi\overline{\chi(g)}\chi(y)=h\mathbf1_{\{g\}}(y)\).
By the definition of induced class functions,
\(\operatorname{Ind}_H^G\mathbf1_{\{g\}}\) has value
\(|Z_G(g)|/h\) on \(C\) and zero elsewhere. Since
\(\delta |Z_G(g)|=1\), (A12.9) follows.

For a prime \(\mathfrak p\notin S\), the cosets in \(G/H\)
fixed by a Frobenius element correspond exactly to prime
ideals \(\mathfrak q\) of \(F\) over \(\mathfrak p\) with
relative residue degree one. The trace of
\(\operatorname{Ind}_H^G\chi\) is the sum of the corresponding
\(\chi(\operatorname{Frob}_{\mathfrak q})\).
For these primes,
\(\nu(\mathfrak q)=\eta(\mathfrak p)\), because the ideal
norm is \(\mathfrak p\). Use the same arithmetic Frobenius
convention in this identity and in class field reciprocity.
Multiplying (A12.9) by \(\eta(\mathfrak p)^m\) and summing
therefore gives
\[
 T_m^C(x):=
 \sum_{\substack{N_K\mathfrak p\le x,\ \mathfrak p\notin S\\
                  \operatorname{Frob}_{\mathfrak p}\in C}}
                        \eta(\mathfrak p)^m
 =\delta\sum_{\chi\in\widehat H}
              \overline{\chi(g)}P_{\chi,m}(x)
           +O(\delta n_L\sqrt x).                       \tag{A12.10}
\]
To verify the last error, the extra primes in the right-hand
sum have relative residue degree at least two, hence lie
over primes of \(K\) with norm at most \(\sqrt x\).
There are at most \(d[K:\mathbf Q]\sqrt x=n_F\sqrt x\)
such primes in \(F\). Each of the \(h\) character sums has
terms of modulus one, so their total contribution is at
most \(\delta h n_F\sqrt x\). Primes in \(S\) have been
deleted on both sides; no Frobenius value is assigned to
a ramified prime.

Combining (A12.7)--(A12.10), and using
\(\log Q_C\ge n_L\log4\), proves
\[
 T_m^C(x)=\delta\mathbf1_{m=0}\operatorname{Li}_2(x)
     +O\left(\delta\sqrt x
        \{\log Q_C+n_L\log[(m+1)(2x)]\}\right).
                                                               \tag{A12.11}
\]
All constants remain absolute. Notice that the \(h\)
summands have already been absorbed in the exact
discriminant and in \(n_L\); they are not replaced by
an additional factor \(|G|\).

Apply Erd\H{o}s--Tur\'an, S-NT-ET Theorem III, to the angles
of the primes with Frobenius class \(C\), with cutoff
\(\lceil x\rceil\). The zero-frequency truncation error
is \(O([K:\mathbf Q])\), since there are at most
\([K:\mathbf Q]x\) prime ideals of norm at most \(x\).
It is absorbed because
\(\delta n_L=|C|[K:\mathbf Q]\ge[K:\mathbf Q]\).
For the other terms, the braces in (A12.11) are
\(O({\cal L}_C(x))\) throughout \(1\le m\le\lceil x\rceil\);
their harmonic weights sum to \(O(\log(2x))\).
The \(m=0\) instance of (A12.11) supplies the main term
in (A12.4). Decreasing arcs to an endpoint bounds its
singleton count by the same error; thus all endpoint
conventions are allowed.
Finally the number of absolute-degree-at-least-two primes
of \(K\) with norm at most \(x\) is at most
\([K:\mathbf Q]\sqrt x\le\delta n_L\sqrt x\).
Deleting them proves the degree-one version.
\end{proof}

\paragraph{NT-A13: a shrinking angle in the specified class.}
Fix \(c>0\), \(0<\alpha<1/2\) and
\(u_0=\max\{2,c^{1/\alpha}\}\). Fix an angular center.
Let \(N_{C,c,\alpha}(x)\) count exactly the primes
\(\mathfrak p\notin S\) with \(u_0<N_K\mathfrak p\le x\),
Frobenius class \(C\), and angle in the centered arc
of length \(c(N_K\mathfrak p)^{-\alpha}\). Under (A12.3),
\[
 \begin{split}
 N_{C,c,\alpha}(x)
   ={}&\delta c\int_{u_0}^x\frac{u^{-\alpha}}{\log u}\,du\\
     &+O_{c,\alpha}\left(
       \delta\sqrt x\log(2x){\cal L}_C(x)
       +\delta x^{3/4-\alpha/2}\sqrt{{\cal L}_C(x)}
                     \right),\qquad x>\max\{4,u_0\}.
 \end{split}                                                   \tag{A13.1}
\]
The constant is uniform in the center and all the field and
character parameters. A sufficient relative asymptotic regime is
\[
         {\cal L}_C(x)(\log x)^2=o(x^{1/2-\alpha}).
                                                               \tag{A13.2}
\]
For fixed arithmetic data this holds automatically, and the
main term is asymptotic to
\(\delta c x^{1-\alpha}/[(1-\alpha)\log x]\).
The same assertions hold for absolute-degree-one primes.

\begin{proof}
Partition \((u_0,x]\) geometrically with ratio
\(1+\varepsilon\), \(0<\varepsilon\le1/2\).
The centered arcs are nested, so their endpoint arcs bracket
the moving count on each subinterval. Apply (A12.4) to both
endpoint prime counts and sum. Since
\(\sum_j\sqrt{u_j}\ll\sqrt x/\varepsilon\), the total
discrepancy is at most a constant times \(A/\varepsilon\),
where
\[
 A=\delta\sqrt x\log(2x){\cal L}_C(x),\qquad
 B=\frac{\delta c x^{1-\alpha}}{\log x}.
\]
The two main terms bracket the integral in (A13.1).
Their difference is \(O_{c,\alpha}(B\varepsilon)\):
on one mesh interval the ratio of the arc lengths is at most
\((1+\varepsilon)^\alpha=1+O_\alpha(\varepsilon)\), and
\(\int_{u_0}^x u^{-\alpha}du/\log u
 \ll_{c,\alpha}x^{1-\alpha}/\log x\).
The finite initial range in this integral bound depends only
on \(c,\alpha\). Its possible contribution to the count is
absorbed uniformly using
\([K:\mathbf Q]\le\delta n_L\le\delta\log Q_C/\log4\).

Choose \(\varepsilon=\min\{1/2,\sqrt{A/B}\}\).
Then \(A/\varepsilon+B\varepsilon\ll A+\sqrt{AB}\), and
\[
 \sqrt{AB}
  =\delta\sqrt c\,x^{3/4-\alpha/2}
       \sqrt{\frac{\log(2x)}{\log x}{\cal L}_C(x)}.
\]
This proves (A13.1). After division by the main term,
the two errors are bounded, up to fixed \(c,\alpha\)
constants, by
\[
 R(x)={\cal L}_C(x)(\log x)^2x^{\alpha-1/2},
 \qquad \sqrt{R(x)}.
\]
They tend to zero under (A13.2). Applying the identical
partition to the degree-one version of (A12.4) proves
the final assertion.
\end{proof}

\paragraph{NT-A14: a least prime with both conditions.}
For an arc of length \(0<\mu\le1\), put
\[
       V_C=\log(eQ_C)+n_L\log(e n_L/\mu).
\]
There is an absolute constant \(C_0>0\) such that (A12.3)
implies the existence of an absolute-degree-one prime
\(\mathfrak p\notin S\), with Frobenius class \(C\) and
angle in this arc, satisfying
\[
                  N_K\mathfrak p\le C_0\mu^{-2}V_C^6 .
                                                               \tag{A14.1}
\]
No numerical value of \(C_0\) is claimed.

\begin{proof}
Let \(A_0\) be an absolute error constant in (A12.4) for
degree-one primes. Set \(x=C_0\mu^{-2}V_C^6\), with
\(C_0\ge4\). We have \(V_C\ge n_L\),
\(n_L\log(1/\mu)\le V_C\), and
\(n_L\log V_C\le2V_C\).
For the last inequality write \(V_C=n_L v\), note
\(v\ge\log(e n_L)\), and use
\(\log(n_L v)\le2v\). Consequently
\[
 \begin{split}
 \log x&\le(\log C_0+14)V_C/n_L,\\
 \log(2x)&\le(\log C_0+15)V_C/n_L,\\
 {\cal L}_C(x)&\le(32+2\log C_0)V_C .
 \end{split}
\]
Since \(\operatorname{Li}_2(x)\ge x/(2\log x)\),
the main term \(\delta\mu\operatorname{Li}_2(x)\)
exceeds the error in (A12.4) as soon as
\[
 \sqrt{C_0}>
  2A_0(\log C_0+14)(\log C_0+15)(32+2\log C_0).
\]
Here the common factor \(\delta\) cancels; dropping the
remaining factor \(n_L^{-2}\le1\) makes this condition
only stronger. An absolute \(C_0\) satisfying it exists.
The count is then positive, proving (A14.1).
\end{proof}

\paragraph{Geometric meaning and exact work placement.}
The finite condition is a conjugacy-class Frobenius condition
for the finite \'{e}tale \(G\)-cover over the arithmetic open
\(\operatorname{Spec}\mathcal O_K\setminus S\). Passing to
\(F=L^{\langle g\rangle}\) is the corresponding intermediate
cover; its cyclic character decomposition is the explicit
induction (A12.9). The second condition is a rank-one
compact Hecke torus coordinate, pulled back by the norm.
For normalized algebraic Hecke characters this is a finite
Artin factor coupled to a CM torus representation. For
arbitrary unitary infinity type no algebraic-motive
realization is asserted. Formula (A12.3) gives the exact
rank-one analytic family replacing any unproved general
Artin-holomorphy input.
Full proofs: \texttt{frobenius\_hecke\_angles.tex},
NT-A12--A14, also in \texttt{section.tex}.
Run: \texttt{GRH\_NIGHT/number\_theory}.
Classical premises: NT-A08 (S-NT-IK), S-NT-ET and
S-NT-NEUKIRCH; the simultaneous formulations are proved
here, with no literature-priority claim.


\subsection*{NT-E11. CM Hasse deficits inside a nonabelian Frobenius class}

Let \(E/K\) be an elliptic curve all of whose complex
multiplication is defined over \(K\), and let \(\eta\) be its
normalized unitary CM Hecke character. Thus at good primes
\[
 a_{\mathfrak p}
     =\sqrt{N_K\mathfrak p}
          \{\eta(\mathfrak p)+\overline{\eta(\mathfrak p)}\}.
                                                               \tag{E11.1}
\]
The CM character premise and its precise classical locator
are those of NT-E06 (S-NT-BFK, Proposition 2.5 and the CM
identification in the footnote to Proposition 4.7).
Its nonzero angular infinity type excludes finite order
times a unitary norm power.
Choose any finite Galois \(L/K\), any conjugacy class \(C\),
and the notation \(\delta,Q_C,{\cal L}_C\) of NT-A12.
Include in \(S\) all bad-reduction primes, as well as the
ramification and conductor primes required there.
Assume exactly the cyclic rank-one RH family (A12.3).
Conditional input \(M\) supplies it.

Fix \(c>0\), \(-1/2<\lambda<1/2\), and put
\[
 \begin{aligned}
 \alpha&=\frac14-\frac{\lambda}{2},\qquad
 u_0=\max\{2,c^{1/(1/2-\lambda)}\},\\
 w(u)&=\frac1\pi
       \arccos\left(1-\frac c2 u^{\lambda-1/2}\right)
                  \quad(u\ge u_0).
 \end{aligned}                                             \tag{E11.2}
\]
For either sign, count the primes \(\mathfrak p\notin S\)
with \(u_0<N_K\mathfrak p\le x\), Frobenius class \(C\), and
\[
       0\le 2\sqrt{N_K\mathfrak p}\mp a_{\mathfrak p}
                      <c(N_K\mathfrak p)^\lambda .
\]
Their number \(D_{C;c,\lambda}^{\pm}(x)\) satisfies, for
\(x>\max\{4,u_0\}\),
\[
 \begin{split}
 D_{C;c,\lambda}^{\pm}(x)
   ={}&\delta\int_{u_0}^x\frac{w(u)}{\log u}\,du\\
     &+O_{c,\lambda}\left(
        \delta\sqrt x\log(2x){\cal L}_C(x)
        +\delta x^{5/8+\lambda/4}\sqrt{{\cal L}_C(x)}
                        \right).                         
 \end{split} \tag{E11.3}
\]
All other dependence is explicit in (A12.2). The same formula
holds when only absolute-degree-one primes are counted.
A sufficient uniform relative asymptotic regime is
\[
       {\cal L}_C(x)(\log x)^2=o(x^{1/4+\lambda/2}).
                                                               \tag{E11.4}
\]
For fixed arithmetic data it gives
\[
 D_{C;c,\lambda}^{\pm}(x)
    \sim \frac{\delta\sqrt c}{\pi(3/4+\lambda/2)}
                    \frac{x^{3/4+\lambda/2}}{\log x}.
                                                               \tag{E11.5}
\]

In particular let \(M_C^\pm(x)\) count the primes outside \(S\)
of norm at most \(x\), Frobenius class \(C\), and
\(a_{\mathfrak p}=\pm\lfloor2\sqrt{N_K\mathfrak p}\rfloor\).
For fixed \(E,K,L,C,S\),
\[
 \begin{split}
 M_C^\pm(x)
     &=\frac{\delta}{\pi}
          \int_2^x\frac{u^{-1/4}}{\log u}\,du
                +O_{E,K,L,C,S}\bigl(x^{5/8}\sqrt{\log(2x)}\bigr),\\
 M_C^\pm(x)&\sim
             \frac{4\delta}{3\pi}\frac{x^{3/4}}{\log x}.
 \end{split}                                                   \tag{E11.6}
\]
The integral, rather than only its leading asymptotic term, is
retained in the error-bearing formula. No general non-CM
elliptic \(L\)-function RH is used.

\begin{proof}
Write \(\eta(\mathfrak p)=e^{2\pi i\theta_{\mathfrak p}}\).
For the positive sign the displayed deficit condition says
\[
 \cos(2\pi\theta_{\mathfrak p})
                      >1-\frac c2 q^{\lambda-1/2},
                \qquad q=N_K\mathfrak p .
 \]
It is exactly the arc about zero of normalized length \(w(q)\),
with its outer endpoints omitted. The negative sign gives
the arc of the same length centered at \(1/2\). Equality
at the Hasse endpoint is allowed and is already in these
arcs. All endpoint conventions are covered by NT-A12.

Here is the required uniform control of the moving length.
Set \(v=cu^{-2\alpha}\), so \(0<v\le1\). Direct integration
and differentiation give
\[
 w(u)=\frac1\pi\int_0^v\frac{dz}{\sqrt{4z-z^2}},
 \qquad
 -u w'(u)=\frac{\alpha\sqrt v}{\pi\sqrt{1-v/4}},
 \qquad w(u)\ge\frac{\sqrt v}{\pi}.
 \]
Consequently
\[
 w(u)\asymp_{c,\lambda}u^{-\alpha},\qquad
    0\le-\frac{u w'(u)}{w(u)}
                    \le\frac{2\alpha}{\sqrt3}.            \tag{E11.7}
\]
On a multiplicative interval of ratio at most
\(1+\varepsilon\), with \(0<\varepsilon\le1/2\),
the endpoint lengths therefore have ratio
\(1+O_\lambda(\varepsilon)\).

Partition \((u_0,x]\) by such intervals. The arcs centered
at zero, or at \(1/2\), are nested. Their two endpoint
counts, using (A12.4), bracket the moving count. As in the
complete proof of NT-A13, the total error is
\[
 \begin{gathered}
 O_{c,\lambda}(A/\varepsilon+B\varepsilon),\\
 A=\delta\sqrt x\log(2x){\cal L}_C(x),\qquad
 B=\frac{\delta x^{1-\alpha}}{\log x}.
 \end{gathered}                                           \tag{E11.8}
\]
For completeness the first term follows from
\(\sum_j\sqrt{u_j}\ll\varepsilon^{-1}\sqrt x\).
For the second, (E11.7) bounds the difference of the
upper and lower integral sums by
\[
 O_\lambda\left(\delta\varepsilon
                    \int_{u_0}^x\frac{w(u)}{\log u}\,du\right)
             \ll_{c,\lambda} B\varepsilon .
 \]
Their common limiting integral is precisely the first
term in (E11.3), so no small-angle expansion has been
substituted into this comparison. The initial bounded
range is uniform after absorption by
\([K:\mathbf Q]\le\delta n_L\) and
\(\log Q_C\ge n_L\log4\).

Take \(\varepsilon=\min\{1/2,\sqrt{A/B}\}\).
Then (E11.8) is \(O_{c,\lambda}(A+\sqrt{AB})\), and
\[
 \sqrt{AB}\ll
          \delta x^{3/4-\alpha/2}\sqrt{{\cal L}_C(x)}
       =\delta x^{5/8+\lambda/4}\sqrt{{\cal L}_C(x)}.
 \]
This proves (E11.3). Using the degree-one form of (A12.4)
throughout proves its degree-one version.
Since \(w(u)\sim(\sqrt c/\pi)u^{-\alpha}\), the integral
has asymptotic
\[
 \frac{\delta\sqrt c}{\pi(1-\alpha)}
                          \frac{x^{1-\alpha}}{\log x}.
 \]
The relative errors are bounded by a constant times
\(R(x)+\sqrt{R(x)}\), where
\(R(x)={\cal L}_C(x)(\log x)^2x^{\alpha-1/2}\).
This proves (E11.4)--(E11.5).

For (E11.6), take \(c=1\), \(\lambda=0\). For any integer
\(a\le y\), \(a=\lfloor y\rfloor\) is equivalent to
\(a>y-1\); this includes integral \(y\).
Apply this with \(a=a_{\mathfrak p}\), or
\(-a_{\mathfrak p}\), and \(y=2\sqrt q\).
The floor event is exactly the strict deficit condition
in (E11.3), even when \(q\) is a square.
The primes of norm exactly two contribute at most
\([K:\mathbf Q]\le\delta n_L\), absorbed by its error.
The elementary expansion on \(u\ge2\) is
\[
 w(u)=\frac1\pi u^{-1/4}+O(u^{-3/4}).
 \]
Its integral remainder is \(O(x^{1/4})\) and is absorbed
by the first error in (E11.3).
For fixed arithmetic data
\({\cal L}_C(x)=O_{E,K,L,C,S}(\log(2x))\).
Thus the second error is \(O(x^{5/8}\sqrt{\log(2x)})\);
the first, \(O(\sqrt x\log^2(2x))\), is smaller for
sufficiently large \(x\), and is absorbed after adjusting
the fixed-data constant. Finally
\(\int_2^x u^{-1/4}du/\log u
 \sim(4/3)x^{3/4}/\log x\), which proves (E11.6).
\end{proof}

\paragraph{Scope and geometric placement.}
The finite étale cover may have an arbitrary nonabelian
Galois group; for example it may be a division-field
cover attached to another, non-CM elliptic curve.
The trace in this theorem is still the trace of the
specified CM curve \(E\). The proof couples the finite
Frobenius condition to its rank-one CM torus by the cyclic
intermediate covers of NT-A12. It does not infer GRH for
the standard \(L\)-function of that second curve, or a
Koblitz prime-order statement. The assumption that all
CM is defined over \(K\) is retained; no unproved descent
factor is inserted into this simultaneous formulation.
At \(\lambda=-1/2\), the positive exponent in (E11.4)
vanishes and the argument gives no asymptotic; the
finite norm-form gate in NT-E10 remains applicable.
For general \(\lambda\), the exact arcsine integral in
(E11.3) is essential to the stated error, as explained
in NT-E09. This is a newly proved formulation, with no
claim of literature priority.
Work/run: \texttt{GRH\_NIGHT/number\_theory},
\texttt{cm\_frobenius\_deficits.tex}, NT-E11, also
\texttt{section.tex}; full analytic input and its
classical locators are (A12.3), NT-A08 and NT-E06.


\subsection*{NT-A15--A17. Several Hecke coordinates and the exact independence gate}

Retain \(L/K,G,C,g,H,F,h,\delta,n_L,S,\mathfrak s_F\)
from NT-A12. Fix \(r\ge1\) unitary Hecke characters
\(\eta_1,\ldots,\eta_r\) of \(K\). Require that the character
\[
             \prod_{i=1}^r\eta_i^{m_i},
                    \qquad 0\ne\mathbf m\in\mathbf Z^r,  \tag{A15.1}
\]
is never finite order times a unitary norm power.
Enlarge \(S\) to contain all their conductor primes.
Write \(\nu_i=\eta_i\circ N_{F/K}\), and put
\[
 \begin{aligned}
 \mathfrak f_*&=\operatorname{lcm}_i\mathfrak f_{\nu_i},\\
 B_*&=4+\sum_{i=1}^r\max_j|\kappa_{\nu_i,j}|,\\
 Q_*&=D_L(N_F\mathfrak f_*\,N_F\mathfrak s_F)^h B_*^{n_L},\\
 {\cal L}_*(x)&=\log Q_*+2n_L\log(3x).
 \end{aligned}                                             \tag{A15.2}
\]
The exact RH family is
\[
 L_F\left(s,\chi\prod_{i=1}^r\nu_i^{m_i}\right),
       \qquad \chi\in\widehat H,\quad\mathbf m\in\mathbf Z^r .
                                                               \tag{A15.3}
\]
It is a subfamily of \(M\). The mixed products in this family
are part of the hypothesis; separate assertions only for
the individual coordinates are not substituted for them.

\paragraph{Classical discrepancy input, with a precise locator.}
For \(N\) points \(\boldsymbol\theta_j\in[0,1)^r\), the
Erd\H{o}s--Tur\'an--Koksma inequality bounds their normalized
rectangular discrepancy by
\[
 \begin{gathered}
 R(\mathbf m)=\prod_i\max(1,|m_i|),\\
 D_N\le(3/2)^r\left\{\frac2{H+1}
     +\sum_{0<\|\mathbf m\|_\infty\le H}
       \frac{1}{R(\mathbf m)}
       \left|\frac1N\sum_{j=1}^N
                  e^{2\pi i\mathbf m\cdot\boldsymbol\theta_j}
       \right|\right\}.
 \end{gathered}                                             \tag{A15.4}
\]
The exact statement is Dick--Pillichshammer,
\emph{Digital Nets and Sequences}, author preprint
Theorem 3.26, printed/PDF p.~84 (S-NT-DP).
That page refers its proof to Drmota--Tichy,
\emph{Sequences, Discrepancies and Applications},
\S1.2.2. The theorem page was image checked; the
referred proof is a cited premise, not a newly audited
argument in this run.

\paragraph{NT-A15: joint boxes and a least prime.}
For a product \(J=J_1\times\cdots\times J_r\) of circular
arcs, let \(\mu=\prod_i|J_i|\). Under (A15.1)--(A15.3),
the count of primes outside \(S\), Frobenius class \(C\),
norm at most \(x\), and coordinate angles in \(J\) is
\[
 N_{C,J}^{(r)}(x)=\delta\mu\operatorname{Li}_2(x)
       +O_r\bigl(\delta\sqrt x
                    (\log(2x))^r{\cal L}_*(x)\bigr).
                                                               \tag{A15.5}
\]
The constant depends only on \(r\); all other dependence
is displayed. Endpoint conventions are arbitrary, and
the same bound holds for absolute-degree-one primes.
For \(\mu>0\), this box contains such a degree-one prime
of norm at most
\[
 C_r\mu^{-2}V_*^{\,2r+4},\qquad
       V_*=\log(eQ_*)+n_L\log(e n_L/\mu),                 \tag{A15.6}
\]
where \(C_r\) depends only on \(r\).

\begin{proof}
Apply the diagonal-idele argument of NT-A12 to each
nonzero monomial in (A15.1). Its norm pullback is still
not finite order times norm. Multiplication by a
finite-order \(\chi\) therefore cannot make it pure norm.
Only \(\chi=1,\mathbf m=0\) is principal.
For negative powers the finite conductor is unchanged;
for any product its conductor divides the least common
multiple of the individual conductors, since all factors
are trivial on each sufficiently deep local unit group.
Thus
\[
 {\cal C}\left(\chi\prod_i\nu_i^{m_i}\right)
  \le D_FN_F\mathfrak f_\chi N_F\mathfrak f_*
       [B_*(1+\|\mathbf m\|_\infty)]^{n_F}.              \tag{A15.7}
\]
At a complex place the absolute value of the sum of
angular parameters is at most the sum of their absolute
values; imaginary parameters add. At real places the
parity remains zero or one. These facts prove the
gamma-factor part of (A15.7), including negative \(m_i\).
Summing logarithms over \(\chi\) uses exactly
\(D_L=D_F^h\prod_\chi N_F\mathfrak f_\chi\).

The whole cyclic induction and omitted-prime argument
(A12.8)--(A12.10) now applies to the monomial.
In particular its relative-degree correction remains
\(\delta n_L\sqrt x\), independently of \(\mathbf m\).
It gives the complete Fourier estimate
\[
 \begin{split}
 T_{\mathbf m}^C(x)
  &:=
    \sum_{\substack{N_K\mathfrak p\le x,\ \mathfrak p\notin S\\
                     \operatorname{Frob}_{\mathfrak p}\in C}}
                          \prod_i\eta_i(\mathfrak p)^{m_i}\\
  &=\delta\mathbf1_{\mathbf m=0}\operatorname{Li}_2(x)
      +O\left(\delta\sqrt x\{
         \log Q_*+n_L\log[(1+\|\mathbf m\|_\infty)(2x)]
                                  \}\right).
 \end{split}                                                   \tag{A15.8}
\]
The constant here is absolute: any dependence on \(r\)
has been put in \(B_*\), not in the Hecke prime estimate.
The principal factor is again handled separately.

Use (A15.4) in unnormalized form with \(H=\lceil x\rceil\).
The sum of the frequency weights is exactly
\[
 \sum_{0<\|\mathbf m\|_\infty\le H}\frac1{R(\mathbf m)}
       =\left(1+2\sum_{j=1}^H\frac1j\right)^r-1
                     \ll_r(\log(2H))^r .
 \]
The zero-frequency truncation contributes
\(O_r([K:\mathbf Q])\), absorbed as in NT-A12 by
\(\delta n_L\le\delta\log Q_*/\log4\).
All the braces in (A15.8) are \(O({\cal L}_*(x))\)
on this frequency cube. These facts prove (A15.5)
for half-open boxes not crossing the circle origins.
A circular box is a disjoint union of at most \(2^r\)
such boxes. Shrinking a coordinate interval onto a
specified endpoint bounds each coordinate-face count
by the same error. Summing the finitely many faces
permits all endpoint conventions.
The number of omitted absolute-degree-at-least-two
primes is at most \([K:\mathbf Q]\sqrt x\), again
absorbed. If the prime list is empty, the unnormalized
discrepancy assertion follows directly from the
zero-frequency estimate, so no division by zero is used.

For (A15.6), set \(V=V_*\), \(n=n_L\), and
\(x=C_r\mu^{-2}V^{2r+4}\). The inequalities
\(V\ge n\), \(n\log(1/\mu)\le V\), and
\(n\log V\le2V\), proved in NT-A14, imply
\[
 \begin{aligned}
 \log x&\le(\log C_r+4r+10)V/n,\\
 \log(2x)&\le(\log C_r+4r+11)V/n,\\
 {\cal L}_*(x)&\le(8r+24+2\log C_r)V.
 \end{aligned}
\]
Let \(A_r\) be a constant in the degree-one error
(A15.5). Since \(\operatorname{Li}_2(x)\ge x/(2\log x)\),
positivity follows, after cancellation of \(\delta\),
from
\[
 \sqrt{C_r}>
  2A_r(\log C_r+4r+10)(\log C_r+4r+11)^r
                    (8r+24+2\log C_r).
 \]
An appropriate \(C_r\ge4\) exists for every fixed \(r\).
The discarded factor \(n^{-(r+1)}\) is at most one.
This proves the least-prime assertion.
\end{proof}

\paragraph{NT-A16: simultaneously shrinking boxes.}
Fix \(c_i>0\), \(\alpha_i>0\), and put
\[
 A=\sum_i\alpha_i<\tfrac12,\qquad
 c_*=\prod_i c_i,\qquad
 u_0=\max\{2,c_1^{1/\alpha_1},\ldots,c_r^{1/\alpha_r}\}.
 \]
Fix every coordinate center and count \(u_0<N_K\mathfrak p\le x\)
in Frobenius class \(C\) whose \(i\)-th angular coordinate
lies in the arc of length \(c_i(N_K\mathfrak p)^{-\alpha_i}\).
Then
\[
 \begin{split}
 N_{C,\mathbf c,\boldsymbol\alpha}(x)
   ={}&\delta c_*\int_{u_0}^x\frac{u^{-A}}{\log u}\,du\\
     &+O_{\mathbf c,\boldsymbol\alpha,r}\left(
        \delta\sqrt x(\log(2x))^r{\cal L}_*(x)\right.\\
     &\hspace{38mm}\left.
        +\delta x^{3/4-A/2}
             (\log(2x))^{(r-1)/2}\sqrt{{\cal L}_*(x)}
                                    \right).
 \end{split}                                                   \tag{A16.1}
\]
The error is uniform in the centers and all arithmetic
parameters already displayed. The degree-one version
holds too. A sufficient relative asymptotic regime is
\[
          {\cal L}_*(x)(\log x)^{r+1}
                                  =o(x^{1/2-A}).           \tag{A16.2}
\]

\begin{proof}
The boxes are nested because every side length decreases.
Partition \((u_0,x]\) with multiplicative ratio
\(1+\varepsilon\), \(0<\varepsilon\le1/2\).
Applying (A15.5) to both endpoints gives a total error
at most a constant times \(A_0/\varepsilon\), where
\[
 A_0=\delta\sqrt x(\log(2x))^r{\cal L}_*(x).
 \]
Indeed the endpoint square roots form the same geometric
series as in NT-A13. The volume is exactly \(c_*u^{-A}\);
its ratio on a mesh interval is at most
\((1+\varepsilon)^A=1+O_A(\varepsilon)\).
Consequently the integral bracketing error is
\(O_{\mathbf c,\boldsymbol\alpha}(B_0\varepsilon)\), with
\[
                  B_0=\frac{\delta c_* x^{1-A}}{\log x}.
 \]
The canonical \(u_0\) bounds the finite initial range
using only the fixed \(\mathbf c,\boldsymbol\alpha\);
as before its field degree is absorbed by
\([K:\mathbf Q]\le\delta n_L\).
Choose \(\varepsilon=\min\{1/2,\sqrt{A_0/B_0}\}\).
The resulting error is \(O(A_0+\sqrt{A_0B_0})\), and
\[
 \sqrt{A_0B_0}
  =\delta\sqrt{c_*}\,x^{3/4-A/2}
          \sqrt{\frac{(\log(2x))^r}{\log x}{\cal L}_*(x)}.
 \]
This is the second term of (A16.1). Its main term is
asymptotic to
\(\delta c_*x^{1-A}/[(1-A)\log x]\).
The two relative errors are bounded by fixed constants
times \(R+\sqrt R\), where
\(R={\cal L}_*(x)(\log x)^{r+1}x^{A-1/2}\).
This proves (A16.2), and the degree-one argument is
identical. No assertion is made at \(A=1/2\).
\end{proof}

\paragraph{NT-A17: necessity of the monomial independence condition.}
For fixed \(L/K,C,S\), under \(M\) the joint prime angles
in this conjugacy class are Haar equidistributed on
\((\mathbf R/\mathbf Z)^r\) if and only if (A15.1) holds.
More precisely, if
\(\prod_i\eta_i^{m_i}=\theta|\cdot|^{it}\) for nonzero
\(\mathbf m\), with \(\theta\) of finite order dividing \(b\),
rename the real norm parameter so that the ideal value of
the \(b\)-th power is \(q^{i\tau}\). Then its normalized
prime Fourier moment is
\[
 \frac{1}{N_C(x)}
 \sum_{\substack{N_K\mathfrak p\le x,\ \mathfrak p\notin S\\
                  \operatorname{Frob}_{\mathfrak p}\in C}}
          \prod_i\eta_i(\mathfrak p)^{b m_i}
       =\frac{x^{i\tau}}{1+i\tau}+o(1).                  \tag{A17.1}
\]
Here \(N_C(x)\) is the number of primes in the sum.
For \(\tau=0\) the left side is exactly one once \(N_C(x)>0\).
Thus neither separate infinite order nor separate
coordinate equidistribution replaces (A15.1).

\begin{proof}
Sufficiency follows from (A15.5) for every fixed box:
its error is \(o(\delta\operatorname{Li}_2(x))\) for fixed
data, and the zero-frequency count has that denominator
asymptotic. Boxes characterize Haar equidistribution.

For necessity, raise the proposed relation to the \(b\)-th
power. Outside the conductor set, every summand in (A17.1)
is exactly \((N_K\mathfrak p)^{i\tau}\).
The fixed-field Chebotarev estimate under RH for \(\zeta_L\),
which \(M\) supplies, gives
\[
 N_C(u)=\delta\operatorname{Li}_2(u)
                       +O_{L,K,C,S}(\sqrt u\log(2u)).
 \]
Abel summation, with fixed \(\tau\), therefore yields
\[
 \sum_{\mathfrak p\text{ in the sum}}
                  (N_K\mathfrak p)^{i\tau}
    =\delta\int_2^x\frac{u^{i\tau}}{\log u}\,du
                +O_{L,K,C,S,\tau}(\sqrt x\log(2x)).
 \]
The lower endpoint can add only a fixed constant.
Differentiating
\(u^{1+i\tau}/[(1+i\tau)\log u]\) shows that this integral is
\[
       \frac{x^{1+i\tau}}{(1+i\tau)\log x}
                         +O_\tau(x/(\log x)^2).
 \]
Division by \(N_C(x)\sim\delta x/\log x\) proves (A17.1).
Its modulus tends to \(1/\sqrt{1+\tau^2}>0\).
The corresponding nonzero Fourier character has Haar
integral zero, so Haar equidistribution fails.
\end{proof}

\paragraph{Geometric meaning and work placement.}
The character lattice \(\mathbf Z^r\) describes the
compact torus coordinates; (A15.1) removes every
finite character or norm drift in that lattice.
For normalized algebraic Hecke characters this is the
torus part of the tensor category of the corresponding
CM objects, coupled to the finite \'{e}tale \(G\)-cover.
The mixed Hecke family (A15.3) records every tensor
character needed by the proof. This is not statistical
independence of zeros or a new assertion of functoriality
for general automorphic representations.
Full proofs: \texttt{joint\_hecke\_tori.tex}, NT-A15--A17,
also \texttt{section.tex}; run
\texttt{GRH\_NIGHT/number\_theory}.
Classical premises are NT-A08/NT-A12, S-NT-NEUKIRCH,
S-NT-DP and the explicit Chebotarev result NT-C02.


\subsection*{NT-E12. Distinct CM fields obstruct simultaneous integral extremes}

Let \(E_1,E_2\) be elliptic curves over a number field \(K\),
with all their complex multiplication defined over \(K\).
Suppose their imaginary quadratic CM fields \(M_1,M_2\)
are distinct. At a prime of common good reduction, put
\[
 q=N_K\mathfrak p,\qquad
 a_i=q+1-\#E_i(\mathcal O_K/\mathfrak p).
 \]
Then the following implication is unconditional:
\[
            a_1^2=a_2^2
       \quad\Longrightarrow\quad
       q\text{ is a square and }a_1^2=a_2^2=4q.
                                                               \tag{E12.1}
\]
In particular, at an absolute-degree-one prime the two
absolute traces are different. Simultaneous integral
extremal events
\[
            |a_1|=|a_2|=\lfloor2\sqrt q\rfloor
                                                               \tag{E12.2}
\]
can occur only at square norms and exact Hasse endpoints.
For all primes of norm at most \(x\) satisfying (E12.2),
the elementary bound
\[
        \#\{\mathfrak p:\text{(E12.2)},\ N_K\mathfrak p\le x\}
                              \le [K:\mathbf Q]\sqrt x
                                                               \tag{E12.3}
\]
holds. No RH assumption is used.

\begin{proof}
We give the centralizer argument, including the rational
Frobenius exception. Let \(E/K\) be either curve, with
CM field \(M\), and let \(\overline E\) be its good
reduction over \(\mathbf F_q\). Good-reduction
specialization injects \(\operatorname{End}_K(E)\) into
\(\operatorname{End}_{\mathbf F_q}(\overline E)\).
Thus \(M\) embeds in
\(A=\operatorname{End}_{\overline{\mathbf F}_q}^0(\overline E)\).
Every element in that embedded field commutes with the
\(q\)-power Frobenius \(\pi\), since its endomorphisms
are defined over \(\mathbf F_q\).

Choose any auxiliary prime \(\ell\ne\operatorname{char}\mathbf F_q\).
The faithful Tate-module action injects
\(A\otimes\mathbf Q_\ell\) into
\(\operatorname{End}_{\mathbf Q_\ell}(V_\ell\overline E)\),
whose underlying space has dimension two.
The action of the two-dimensional \'{e}tale algebra
\(M\otimes\mathbf Q_\ell\) is faithful and has rank one
over that algebra. If the algebra is a field, this
follows from the two dimensions; if it is
\(\mathbf Q_\ell\times\mathbf Q_\ell\), faithfulness
makes each of its two idempotent summands nonzero,
so each has dimension one.
Its centralizer in the two-by-two endomorphism algebra
is therefore exactly \(M\otimes\mathbf Q_\ell\).

Commutation with a basis of \(M\) is a finite system of
rational linear equations in \(A\). The centralizer of
\(M\) in \(A\) consequently has rational dimension at
most two, and already contains \(M\). It equals \(M\).
Since \(\pi\) commutes with \(M\), we conclude \(\pi\in M\).
These specialization and faithful-action facts are
the same classical good-reduction interface used in
the full proof of NT-E10; no assertion about a
Tate-conjecture surjectivity is needed.

The characteristic polynomial on \(V_\ell\) is
\[
                       X^2-aX+q .
 \]
If \(\pi\notin\mathbf Q\), then \(\mathbf Q(\pi)=M\);
its minimal polynomial has degree two and must be the
displayed characteristic polynomial. In particular
\[
                 M=\mathbf Q(\sqrt{a^2-4q}),
                         \qquad a^2-4q<0.               \tag{E12.4}
 \]
If instead \(\pi=t\in\mathbf Q\), its action is scalar.
Its trace and determinant give \(a=2t\) and \(q=t^2\).
As \(q\) is an integer, a rational square root \(t\)
is an integer. Thus \(q\) is a square and \(a^2=4q\).
Conversely a repeated root of the displayed polynomial
forces \(\pi=a/2\): inside the field \(M\), an element
whose square \((\pi-a/2)^2\) is zero is zero.

Now suppose \(a_1^2=a_2^2\). If this common value is
not \(4q\), both Frobenius elements are nonrational,
and (E12.4) identifies both \(M_1\) and \(M_2\) with
the same imaginary quadratic field. This contradicts
their distinctness. Hence the common value is \(4q\),
which proves (E12.1). If \(q=p\) is a rational prime
it is not a square, giving the degree-one conclusion.
More generally every odd absolute residue degree is
excluded.

Event (E12.2) implies equality of squared traces, so
it can occur only as stated. At square norm \(q=p^{2f}\le x\),
the underlying rational prime satisfies \(p\le\sqrt x\).
There are at most \([K:\mathbf Q]\) prime ideals of \(K\)
over each rational prime. Summing this bound over
the rational primes at most \(\sqrt x\) proves
(E12.3), even with the weaker count by all integers.
\end{proof}


\paragraph{Earlier rational special case.}
Bucur--Fit\'e--Kedlaya, Remark~4.11, published p.~1743
(author v2 p.~22), already gives the obstruction
$a_p(E_1)=a_p(E_2)=\lfloor2\sqrt p\rfloor$ over
$\mathbf Q$ for distinct CM fields and $p>3$,
crediting Andrew Sutherland. The formulation here
retains arbitrary base fields, equal absolute traces,
the precise square-norm exception, and the quantitative
count of higher-degree primes. The rational positive-extreme
special case is not claimed as a new literature result.

\paragraph{Meaning of the shrinking-boundary gate.}
Each separate extremal CM trace asks for an angle
of width comparable to \(q^{-1/4}\). Two such
requirements have total shrinking exponent \(1/2\),
exactly the boundary excluded from NT-A16.
Equations (E12.1)--(E12.3) exhibit a genuine integral
Frobenius obstruction there: continuous torus Haar
mass cannot be used to infer degree-one prime
infinitude at this scale. The obstruction concerns
the two integral traces and their quadratic
endomorphism algebras; it is compatible with
equidistribution in every fixed positive-volume box.
It persists after imposing any finite Frobenius
condition, since that only removes primes.
Geometrically the scalar commutant of each CM action
determines the nonrational Frobenius field, so two
distinct CM algebras cannot give the same noncentral
characteristic polynomial.
Full proof: \texttt{cm\_joint\_endpoint\_obstruction.tex},
NT-E12, also \texttt{section.tex}; run
\texttt{GRH\_NIGHT/number\_theory}.


\subsection*{NT-E13. The endpoint obstruction descends to potential CM curves}

Let \(E_1,E_2\) be elliptic curves over a number field \(k\)
with distinct imaginary quadratic potential-CM fields
\(M_1,M_2\). All endomorphisms need not be defined over \(k\).
Put \(F_i=kM_i\), \(F=F_1F_2\), and let \(S\) contain
the bad-reduction primes of both curves and the finite
primes ramified in \(F/k\). At any prime
\(\mathfrak p\notin S\), write \(q=N_k\mathfrak p\)
and let \(a_i\) be the two Frobenius traces.
Then
\[
       a_1^2=a_2^2\ne0
       \quad\Longrightarrow\quad
       \begin{cases}
       \mathfrak p\text{ splits completely in }F/k,\\
       q\text{ is a square and }a_1^2=4q .
       \end{cases}
                                                               \tag{E13.1}
\]
Consequently the simultaneous integral extreme
\[
              |a_1|=|a_2|=\lfloor2\sqrt q\rfloor
                                                               \tag{E13.2}
\]
is impossible at every absolute-degree-one prime outside
\(S\). Among all prime ideals outside \(S\) of norm at most
\(x\), its count is at most \([k:\mathbf Q]\sqrt x\).
These assertions are unconditional.
The nonzero condition in (E13.1) is essential: both curves
may have trace zero in nonsplit CM components.

\begin{proof}
For a characteristic-zero CM elliptic curve the field
over which all endomorphisms are defined is \(kM_i\).
One way to see the exact quadratic action is to let an
endomorphism act on the one-dimensional tangent space.
In characteristic zero this is a faithful map from the
CM field into \(\overline k\), compatible with Galois
action. The action on endomorphisms therefore factors
through the two automorphisms of \(M_i\), with kernel
precisely the group fixing \(kM_i\).
This is the same endomorphism-field premise used in
NT-E08 and NT-E10.

Fix an auxiliary prime \(\ell\) different from the
residue characteristic. Over an algebraic closure of
\(\mathbf Q_\ell\), the faithful CM action decomposes
the two-dimensional Tate module into its two distinct
one-dimensional eigenspaces. For an arithmetic
Frobenius at \(\mathfrak p\), conjugation of
endomorphisms agrees with its action on \(M_i\).
If the restriction to \(F_i/k\) is nontrivial, Frobenius
interchanges those two eigenspaces. Its matrix in that
decomposition has zero diagonal, hence trace zero.
The Frobenius trace is \(a_i\), since the curve has
good reduction and \(\mathfrak p\) is unramified in
the indicated quadratic extension.

Thus \(a_i\ne0\) forces trivial Frobenius in \(F_i/k\).
If \(a_1^2=a_2^2\ne0\), the Frobenius is trivial on
both \(F_i\), hence on their compositum \(F\).
This proves the splitting assertion in (E13.1).
Choose a prime \(\mathfrak P\) of \(F\) over \(\mathfrak p\).
It has the same residue field, norm \(q\), and both
traces \(a_i\). The two base-changed curves now have
all CM defined over \(F\), with the same distinct
fields \(M_1,M_2\). NT-E12 applies and forces square
norm and exact Hasse equality, proving (E13.1).

Since \(q\ge2\), its extremal integer
\(\lfloor2\sqrt q\rfloor\) is positive. Event (E13.2)
therefore satisfies the nonzero premise of (E13.1).
It cannot occur at norm a rational prime.
For a square norm \(q=p^{2f}\le x\), the underlying
rational prime is at most \(\sqrt x\), and at most
\([k:\mathbf Q]\) prime ideals of \(k\) lie above each
one. This gives the stated count bound.
\end{proof}

\paragraph{An explicit pair over \(\mathbf Q\).}
For
\[
        E_1:y^2=x^3-x,\qquad E_2:y^2=x^3+1,
                                                               \tag{E13.3}
\]
there is no rational prime \(p>3\) at which both
\(|a_p(E_i)|=\lfloor2\sqrt p\rfloor\).
Indeed the automorphism
\((x,y)\mapsto(-x,iy)\) gives CM by \(\mathbf Z[i]\)
on \(E_1\), while
\((x,y)\mapsto(\zeta_3x,y)\) gives CM by
\(\mathbf Z[\zeta_3]\) on \(E_2\).
These are the maximal orders in the distinct
quadratic fields \(\mathbf Q(i)\) and
\(\mathbf Q(\sqrt{-3})\); their displayed generators
are defined over those fields.
The two integral equations have discriminants
\(64\) and \(-432\), respectively, so both curves
have good reduction at every \(p>3\).
The compositum is \(\mathbf Q(\zeta_{12})\), ramified
only at \(2,3\). Every \(p>3\) is therefore within
the hypotheses just proved.


\paragraph{Earlier rational special case.}
Bucur--Fit\'e--Kedlaya, Remark~4.11, published p.~1743
(author v2 p.~22), already gives the obstruction
$a_p(E_1)=a_p(E_2)=\lfloor2\sqrt p\rfloor$ over
$\mathbf Q$ for distinct CM fields and $p>3$,
crediting Andrew Sutherland. The formulation here
retains arbitrary base fields, equal absolute traces,
the precise square-norm exception, and the quantitative
count of higher-degree primes. The rational positive-extreme
special case is not claimed as a new literature result.

\paragraph{Input, scope and geometric placement.}
The proof combines the exact CM component action with
the rational-Frobenius exception of NT-E12. It does
not presume that a nonsplit component has a continuous
circle distribution; its trace is zero instead.
It gives a concrete obstruction over \(\mathbf Q\)
at the total shrinking exponent \(1/2\), and remains
true after any further finite Frobenius restriction.
The relevant geometry is the component group acting
on the two CM eigenlines, followed on the split
component by the two distinct quadratic endomorphism
algebras. No GRH or Artin-holomorphy input occurs.
Full proof: \texttt{potential\_cm\_joint\_endpoint.tex},
NT-E13; run \texttt{GRH\_NIGHT/number\_theory}.

\subsection*{NT-A18--A19. Reciprocal-prime weighting removes the norm-drift obstruction}

Keep the finite Galois extension \(L/K\), conjugacy class
\(C\), cyclic subgroup \(H=\langle g\rangle\), field
\(F=L^H\), and \(\delta=|C|/|G|\) of NT-A12.
Fix unitary Hecke characters \(\eta_1,\ldots,\eta_r\)
of \(K\), and let \(S\) contain all ramification and
conductor primes. Write
\(\nu_i=\eta_i\circ N_{F/K}\).
For this result the independence condition is only
\[
 \prod_i\eta_i^{m_i}\text{ has infinite order}
              \qquad(0\ne\mathbf m\in\mathbf Z^r).         \tag{A18.1}
\]
Nontrivial unitary norm powers are allowed.
The sufficient analytic input is RH for \(\zeta_L\)
and for every primitive, non-pure-norm member of
\[
 L_F\left(s,\chi\prod_i\nu_i^{m_i}\right),
      \quad \chi\in\widehat H,\quad
                         0\ne\mathbf m\in\mathbf Z^r .
                                                               \tag{A18.2}
\]
Pure norm members require only RH for \(\zeta_F\),
already supplied by RH for \(\zeta_L\). Thus \(M\)
supplies exactly the stated rank-one input, without
general Artin holomorphy. All data in this subsection
are fixed; no uniformity as frequencies approach a
norm relation is asserted.

\paragraph{NT-A18: finite harmonic constants and the exact weighted torus gate.}
Let \(\mathcal P_C\) be the primes outside \(S\) with
Frobenius class \(C\). For \(q_{\mathfrak p}=N_K\mathfrak p\),
put
\[
 W_C(x)=\sum_{\substack{\mathfrak p\in\mathcal P_C\\
                                      q_{\mathfrak p}\le x}}
                       q_{\mathfrak p}^{-1},\qquad
 W_{\mathbf m}(x)=
 \sum_{\substack{\mathfrak p\in\mathcal P_C\\
                                      q_{\mathfrak p}\le x}}
       \frac{\prod_i\eta_i(\mathfrak p)^{m_i}}{q_{\mathfrak p}}.
 \]
There are finite constants \(B_C\in\mathbf R\) and
\(B_{\mathbf m}\in\mathbf C\), for every nonzero
\(\mathbf m\), such that
\[
 \begin{aligned}
 W_C(x)&=\delta\log\log x+B_C
                         +O_{L,K,S}\!\left(\frac{\log(2x)}{\sqrt x}\right),\\
 W_{\mathbf m}(x)&=B_{\mathbf m}
          +O_{\mathbf m,\boldsymbol\eta,L,K,S}\!
              \left(\frac1{\log x}+\frac{\log(2x)}{\sqrt x}\right).
 \end{aligned}                                             \tag{A18.3}
\]
If none of the factors for this \(\mathbf m\) in
(A18.2) is pure norm, its \(1/\log x\) term can be
omitted. The series defining \(B_{\mathbf m}\) is
ordered by prime norm; absolute convergence of its
first powers is not claimed.

The probability measures
\[
 \frac1{W_C(x)}
   \sum_{\substack{\mathfrak p\in\mathcal P_C\\q_{\mathfrak p}\le x}}
        \frac{\delta_{(\eta_1(\mathfrak p),\ldots,\eta_r(\mathfrak p))}}
             {q_{\mathfrak p}}
                                                               \tag{A18.4}
\]
therefore tend weakly to Haar measure on \((S^1)^r\).
Every fixed product of arcs has limiting mass equal
to its product length, with arbitrary endpoints.
For each fixed nonzero Fourier frequency the error
is \(O_{\mathbf m}(1/\log\log x)\).
The same conclusions, with possibly changed finite
constants, hold if only absolute-degree-one primes
are counted.

Conversely a finite-order nonzero monomial prevents
full Haar convergence for these weights. Thus, under
\(M\), (A18.1) is the exact condition for full
reciprocal-prime-weighted Haar distribution. It is
strictly weaker than the ordinary-prime condition
in NT-A17.

There is also a convergent Euler-product formulation.
For every nonzero \(\mathbf m\), the products
\[
 \prod_{\substack{\mathfrak p\in\mathcal P_C\\q_{\mathfrak p}\le x}}
 \left(1-\frac{z\prod_i\eta_i(\mathfrak p)^{m_i}}
                    {q_{\mathfrak p}}\right)^{-1}
                                                               \tag{A18.5}
\]
converge locally uniformly for \(|z|<2\) to a
nonvanishing holomorphic function. In particular
their values at \(z=1\) have finite nonzero limits.

\begin{proof}
First consider a single primitive unitary Hecke
character \(\lambda\) of \(F\), with all conductor
primes removed. Additional finitely many primes
may also be removed. If \(\lambda\) is not pure norm,
NT-A08 under its RH gives, for fixed data,
\[
 A_\lambda(u):=\sum_{N_F\mathfrak q\le u}\lambda(\mathfrak q)
                        =O_\lambda(\sqrt u\log(2u)).
 \]
Abel summation, including any primes of norm two,
gives
\[
 \sum_{N_F\mathfrak q\le x}
            \frac{\lambda(\mathfrak q)}{N_F\mathfrak q}
   =\frac{A_\lambda(x)}x+
                         \int_2^x\frac{A_\lambda(u)}{u^2}\,du .
 \]
The integral over \([2,\infty)\) is absolutely
convergent. Its tail and the endpoint term are
\(O_\lambda(\log(2x)/\sqrt x)\). This proves the
claimed finite limit in the non-pure case.

Now let the ideal values of a pure norm character
be \((N_F\mathfrak q)^{i\tau}\), with \(\tau\ne0\).
Under RH for \(\zeta_F\), the prime count after the
same finite deletions is
\(\operatorname{Li}_2(u)+E(u)\), where
\(E(u)=O_{F,S}(\sqrt u\log(2u))\).
Stieltjes integration by parts with
\(u^{i\tau-1}\) gives
\[
 \begin{split}
 \sum_{N_F\mathfrak q\le x}(N_F\mathfrak q)^{i\tau-1}
  ={}&\int_2^x\frac{u^{i\tau-1}}{\log u}\,du
       +x^{i\tau-1}E(x)\\
    &-(i\tau-1)\int_2^x E(u)u^{i\tau-2}\,du .
 \end{split}                                                   \tag{A18.6}
\]
The \(E\)-integral converges absolutely, and its
tail together with the endpoint is
\[
        O_{F,S}\left((1+|\tau|)
                               \frac{\log(2x)}{\sqrt x}\right).
 \]
The main integral becomes
\(\int_{\log2}^{\log x}e^{i\tau y}\,dy/y\).
It converges at infinity by integration by parts,
with the explicit tail bound
\[
       \left|\int_{\log x}^\infty
                      \frac{e^{i\tau y}}y\,dy\right|
                         \le\frac2{|\tau|\log x}.         \tag{A18.7}
\]
These computations prove actual convergence to a
finite constant in the pure norm case too.
There is no application of the shifted-pole
realification argument in this case.

Condition (A18.1) survives norm pullback:
if a pulled-back monomial had finite order \(b\),
composition with the diagonal idele map would
make its original monomial have order dividing
\(b[F:K]\). Thus, for nonzero \(\mathbf m\),
\(\chi\prod_i\nu_i^{m_i}\) is never principal.
If it is pure norm, its parameter \(\tau\) is
therefore nonzero. RH for \(\zeta_L\) supplies
RH for \(\zeta_F\) by the cyclic factorization
\(\zeta_L=\prod_{\chi\in\widehat H}L_F(s,\chi)\):
the nontrivial finite characters are entire, so
a nontrivial zero of any factor cannot be
cancelled in the open critical strip.

Apply the cyclic class-function identity (A12.9)
to the weighted prime sum. The part over primes
of \(F\) with relative residue degree one is
exactly \(W_{\mathbf m}(x)\), with coefficient
\(\delta\overline{\chi(g)}\) in the \(\chi\)-sum.
Each full prime sum over \(F\) has a finite limit
by (A18.6)--(A18.7) or the preceding non-pure
argument. The additional primes of relative
residue degree at least two give an absolutely
convergent correction, uniformly in their unit
character values.

Here is a quantitative check of that correction.
Every such prime has absolute residue degree
at least two. For \(n_F=[F:\mathbf Q]\),
\[
 \begin{split}
 \sum_{\substack{\mathfrak q:\ f_{\mathfrak q}\ge2\\
                                   N_F\mathfrak q>x}}
                \frac1{N_F\mathfrak q}
 &\le n_F\sum_{j=2}^{n_F}\sum_{p>x^{1/j}}p^{-j}\\
 &\le n_F\sum_{j=2}^{n_F}
       \left(x^{-1}+\frac{x^{-1+1/j}}{j-1}\right)
                       \ll_F x^{-1/2}.
 \end{split}                                                   \tag{A18.8}
\]
The inner bound follows by comparison with the
integral of \(t^{-j}\), and the rational primes
may be enlarged to all integers. The empty
sum for \(n_F=1\) causes no exception.
Combining these finitely many \(\chi\)-sums and
corrections proves the second line of (A18.3).

For the zeroth frequency, NT-C02 and finite
deletions give
\[
 \pi_C(u)=\delta\operatorname{Li}_2(u)+R_C(u),
          \qquad R_C(u)=O_{L,K,S}(\sqrt u\log(2u)).
 \]
The complete reciprocal-sum calculation of
NT-C08 applies verbatim: define
\[
 B_C=-\delta\log\log2+
                             \int_2^\infty\frac{R_C(u)}{u^2}\,du .
 \]
Then the error is exactly
\(R_C(x)/x-\int_x^\infty R_C(u)du/u^2\).
This proves the first line of (A18.3), and
\(W_C(x)\to\infty\).
Deleting absolute-degree-at-least-two primes
of \(K\) changes each Fourier sum by a convergent
series with the tail (A18.8), now over \(K\).
This proves every assertion about the degree-one
version.

Divide each nonzero Fourier sum by \(W_C(x)\).
Its limit is zero, with the stated fixed-frequency
rate. For completeness, trigonometric polynomials
are uniformly dense in continuous functions on
the torus, for instance by product Fej\'er means.
For a polynomial the limiting integral is Haar
integration by the just-proved Fourier limits.
For a continuous function approximate uniformly;
both the measures in (A18.4) and Haar measure
have mass one, so the approximation error is
at most twice the uniform error. This proves
weak convergence. Continuous upper and lower
approximations to a fixed box, with arbitrarily
small Haar integral gap at its faces, prove the
box assertion and every endpoint convention.

If instead a nonzero monomial is finite order
dividing \(b\), its \(b\)-th power is identically
one outside the conductor set. The corresponding
nonzero Fourier coefficient of (A18.4) is exactly
one whenever the denominator is positive.
Its Haar integral is zero, which proves necessity.

Finally set \(u_{\mathfrak p}=\prod_i\eta_i(\mathfrak p)^{m_i}\).
For \(|z|<2\), use the power-series logarithm of
each factor in (A18.5). Its first-degree sum is
\(zW_{\mathbf m}(x)\), converging locally uniformly.
For each fixed \(R<2\), all higher-degree sums
are uniformly absolutely convergent on \(|z|\le R\),
since
\[
 \sum_{\mathfrak p\in\mathcal P_C}\sum_{j\ge2}
       \frac{R^j}{j q_{\mathfrak p}^{\,j}}
 \le \frac{R^2}{2(1-R/2)}
                 \sum_{\mathfrak p\in\mathcal P_C}q_{\mathfrak p}^{-2}
                                                    <\infty .
 \]
The resulting logarithm is holomorphic. Its
exponential is the asserted locally uniform
Euler-product limit and never vanishes.
\end{proof}

\paragraph{NT-A19: a Dirichlet-GRH-only progression example.}
Fix primitive Dirichlet characters \(\chi_i\) of
conductors \(f_i\), and real numbers \(t_1,\ldots,t_r\)
which are linearly independent over \(\mathbf Q\).
For primes \(p\nmid\prod_i f_i\), put
\[
                    z_i(p)=\chi_i(p)p^{it_i}.
 \]
Fix a reduced arithmetic progression \(p\equiv a\pmod q\),
and put \(Q_0=\operatorname{lcm}(q,f_1,\ldots,f_r)\).
RH for the primitive Dirichlet characters with
conductor dividing \(Q_0\), including \(\zeta\) for
conductor one, suffices for the following assertions.
This is a finite subfamily of the user-adopted input
\(A_D\), and uses no additional \(M\) input.

The \(1/p\)-weighted points \((z_1(p),\ldots,z_r(p))\)
in that progression tend to full Haar measure on
\((S^1)^r\). Every nonzero mixed prime harmonic sum
converges to a finite constant. In contrast, the
ordinary equal-weight prime points fail to be
Haar equidistributed.

\begin{proof}
Partition the progression into its finitely many
reduced residue classes \(b\bmod Q_0\) satisfying
\(b\equiv a\bmod q\), omitting the finite conductor
primes. On each such class, a Fourier frequency
\(\mathbf m\ne0\) has value
\[
       \prod_i z_i(p)^{m_i}
        =\left(\prod_i\chi_i(b)^{m_i}\right)p^{i\tau_{\mathbf m}},
       \qquad \tau_{\mathbf m}=\sum_i m_i t_i\ne0.
 \]
The progression estimate NT-P12, with the
moduli-one-and-two convention there, gives
\[
 \pi(u;Q_0,b)
   =\frac{\operatorname{Li}_2(u)}{\varphi(Q_0)}
                              +O_{Q_0}(\sqrt u\log(2u)).
 \]
Apply exactly the Stieltjes calculation (A18.6),
with the main integral multiplied by
\(1/\varphi(Q_0)\). Equation (A18.7) proves
convergence of each Fourier harmonic sum on
each class, hence on their finite union.
The zeroth sum is
\(\varphi(q)^{-1}\log\log x+O_{q,\boldsymbol f}(1)\).
The same Fourier approximation argument therefore
proves full weighted Haar distribution.

For the ordinary weighting, choose any coordinate
\(i\) and an integer \(b_i\ge1\) with
\(\chi_i^{b_i}=1\). Its nonzero Fourier frequency
\(b_i\mathbf e_i\) has value \(p^{ib_i t_i}\).
The ordinary Abel calculation in NT-A17, now
using NT-P12 in the fixed progression, gives
\[
 \frac{\sum_{p\le x,\ p\equiv a\bmod q}p^{ib_i t_i}}
               {\pi(x;q,a)}
       =\frac{x^{ib_i t_i}}{1+ib_i t_i}+o(1).
 \]
Omitting conductor primes does not affect the
limit. Its modulus tends to the positive number
\((1+b_i^2t_i^2)^{-1/2}\), whereas the Haar
Fourier integral is zero. This proves the
ordinary-weight obstruction.
\end{proof}

\paragraph{Geometric and analytic meaning.}
The measures are reciprocal-norm measures on
closed points in a specified finite \'{e}tale
Frobenius class, with values in a compact torus
of rank-one Hecke coordinates. Their radial
norm drift is averaged by the logarithmic
measure \(dy/y\); a finite character relation
cannot be averaged away.
For general unitary infinity types this is an
idele-class formulation, not a claimed algebraic
motive realization. No zero-independence or
joint-zero-statistics hypothesis is used.
The constants in (A18.7) can grow when a nonzero
frequency approaches a norm relation, so no
uniform shrinking-box result is inferred here.
Full proofs: \texttt{reciprocal\_prime\_tori.tex},
NT-A18--A19; run \texttt{GRH\_NIGHT/number\_theory}.
Exact classical inputs are NT-A08 (S-NT-IK),
NT-C02 (S-NT-GM), the cyclic induction already
proved in NT-A12, and NT-P12 for the \(A_D\)
example. These are newly proved formulations
in this run, with no literature-priority claim.


\subsection*{NT-A20--A21. Euclidean ideal classes and one-point localization}

\paragraph{NT-A20: the exact Picard-group criterion.}
Let $K$ be a number field, let $S$ be a finite set of places
containing every archimedean place, and put $R=\mathcal O_{K,S}$.
Assume $|S|\ge2$. For every positive squarefree integer $m$, set
\[
 E_m=K\bigl(\zeta_m,(R^\times)^{1/m}\bigr).
 \tag{A20.1}
\]
These are finite number fields: the finitely generated group
$R^\times$ has only finitely many generators whose $m$th roots
need to be adjoined. Assume precisely RH for all $\zeta_{E_m}$.
Input $M$ supplies this family through its trivial Hecke characters.

A fractional ideal class $C\in\operatorname{Pic}(R)$ is
Euclidean in Lenstra's sense if and only if
\[
                 \operatorname{Pic}(R)=\langle C\rangle.
 \tag{A20.2}
\]
Thus, if $\operatorname{Pic}(R)$ is cyclic of order $h_S$,
it has exactly $\varphi(h_S)$ Euclidean ideal classes,
where $\varphi(1)=1$; otherwise it has none.
At most one of these classes can be Euclidean for the ideal norm.
In particular, when $h_S>2$ is cyclic, at least
$\varphi(h_S)-1$ classes admit a Euclidean reduction but
do not admit the norm reduction.

Here is the precise definition and the full elementary part
of the argument. Write
\[
 \mathcal E(R)=\{I:\ I\text{ is a nonzero fractional ideal of }R,\
                                      R\subseteq I\}.
\]
For a representative $\mathfrak c$ of $C$, a Euclidean
reduction is a map $\psi:\mathcal E(R)\to W$ to a well-ordered
set such that for every $I\in\mathcal E(R)$ and
$x\in I\mathfrak c\setminus\mathfrak c$, there is
$y\in\mathfrak c$ satisfying
\[
 J=(x-y)^{-1}I\mathfrak c\in\mathcal E(R),
                    \qquad \psi(J)<\psi(I).
 \tag{A20.3}
\]
The element $x-y$ is nonzero. The displayed membership
holds because $x-y\in I\mathfrak c$.
Replacing $\mathfrak c$ by a principal multiple leaves
the definition unchanged.

\begin{proof}
For necessity, suppose there is an ideal class outside
$\langle C\rangle$. Every fractional ideal class has a
representative in $\mathcal E(R)$: for a nonzero
$a$ in a fractional ideal $\mathfrak a$, the ideal
$a^{-1}\mathfrak a$ contains $R$.
Choose $I\in\mathcal E(R)$ outside that subgroup with
$\psi(I)$ smallest. Then $I\ne R$, and invertibility
of $\mathfrak c$ gives $I\mathfrak c\ne\mathfrak c$.
Choose $x$ in their difference and apply (A20.3).
Since $[J]=[I]C$, the class $[J]$ is still outside
$\langle C\rangle$, contradicting the minimal choice.
This proves necessity without RH.

For sufficiency, apply Lenstra's Proposition~(2.2),
printed p.~127 of \emph{Euclidean ideal classes},
Ast\'erisque~61 (1979), 121--131. Its hypotheses are
exactly $|S|\ge2$ and RH for (A20.1), and its conclusion
is that every class generating $\operatorname{Pic}(R)$
has such a reduction. The printed inequality is
$|S|\ge2$; the searchable OCR incorrectly renders it
as a strict inequality. This is a cited classical
analytic input. The paper refers its sufficiency
argument to the methods of Lenstra~1977, reference~[9];
it does not print a full new analytic proof at p.~127.
We do not represent the statement check as a fresh
verification of that omitted analytic argument.

The number of generators of a finite cyclic group of
order $h_S$ is $\varphi(h_S)$, giving the count.
For the norm assertion, use the single fixed map
\[
                       \psi_N(I)=N_R(I)^{-1}\in\mathbf N.
 \tag{A20.4}
\]
It is integral because $I^{-1}$ is integral and
$N_R(I)^{-1}=\#(R/I^{-1})$.
If two classes were Euclidean for this same map,
choose an $I\ne R$ of smallest $\psi_N(I)$.
For either class, (A20.3) produces a $J$ with smaller
value. Minimality among nontrivial overideals forces
$J=R$, so both classes must equal $[I]^{-1}$.
The set from which $I$ is chosen is nonempty:
there is a prime ideal outside the finite set $S$.
Finally,
\[
 \psi_N\bigl((x-y)^{-1}I\mathfrak c\bigr)<\psi_N(I)
                  \quad\Longleftrightarrow\quad
                     N_R(x-y)<N_R(\mathfrak c),
\]
which is exactly norm-Euclidean reduction. This proves
the claimed uniqueness and the lower count of classes
that are not norm-Euclidean.
\end{proof}

The necessity argument and fixed-function uniqueness
are elementary reformulations of Lenstra's
Lemma~(1.5), Theorem~(1.6), and the paragraph after
Exercise~(1.7), printed pp.~124--125; those complete
elementary proofs were read. The mathematical content
of the sufficiency is classical. The explicit family,
class count, and comparison with $M$ here make its
input and consequences reviewable. No finite-order
Dirichlet input alone supplies the general unit
Kummer family. No complexity bound for evaluating
the well-ordered reduction is claimed.

\paragraph{NT-A21: a bounded single closed point makes a Euclidean ring.}
Let $K$ be any number field whose ordinary class group
is cyclic, of order $h$. Write $n=[K:\mathbf Q]$ and
$D=|\operatorname{disc}K|$. Under $M$, for every
generator $C$ of $\operatorname{Cl}(K)$ there is a
prime ideal $\mathfrak p$ of absolute residue degree
one such that
\[
 [\mathfrak p]=C,\qquad
 N\mathfrak p\le
                \bigl[h(4\log D+2.5n)+5\bigr]^2,
 \tag{A21.1}
\]
and the ring
\[
             R=\mathcal O_{K,S_\infty\cup\{\mathfrak p\}}
 \tag{A21.2}
\]
is Euclidean. In particular, this holds for fields
whose original unit group is finite. The norm bound
controls the one closed point removed; it is not a
bound for the running time of the Euclidean reduction.

The exact analytic inputs, weaker than universal $M$,
are RH for $\zeta_{\mathbf Q}$ and for the ordinary
Hilbert class-field zeta function $\zeta_{H_K}$,
followed by RH for the fields (A20.1) formed from
the particular ring (A21.2). The latter family
depends on the selected prime. No claim that the
first two zeta functions imply that family is made.

\begin{proof}
NT-C07, the checked specialization of Bach--Sorenson
Theorem~5.1, gives (A21.1). Its discriminant
substitution is $D_{H_K}=D^h$ and $[H_K:\mathbf Q]=hn$.
For $K=\mathbf Q$ take $\mathfrak p=(2)$ directly;
it satisfies (A21.1).

Extension of ideals to $R$ gives the exact quotient
\[
 \operatorname{Pic}(R)=
       \operatorname{Cl}(K)/\langle[\mathfrak p]\rangle.
 \tag{A21.3}
\]
Indeed every ideal of $R$ is represented by a finite
product of primes different from $\mathfrak p$,
so extension is onto. If an ideal becomes principal,
its divisor differs from a principal divisor only at
$\mathfrak p$, showing that the kernel is the displayed
cyclic subgroup. Conversely every power of
$\mathfrak p$ becomes the unit ideal.
Since $[\mathfrak p]$ generates $\operatorname{Cl}(K)$,
the group in (A21.3) is zero.

There are at least two places in
$S_\infty\cup\{\mathfrak p\}$. Apply NT-A20 to its
principal ideal class; equivalently apply the
classical PID statement NT-A03. This makes $R$
Euclidean. For clarity, the equivalence between a
principal Euclidean ideal class and a Euclidean domain
can also be seen directly. For nonzero $b\in R$ put
$\phi(b)=\psi(b^{-1}R)$, with $\psi$ from (A20.3) and
$\mathfrak c=R$. If $a/b\notin R$, apply (A20.3) to
$I=b^{-1}R$ and $x=a/b$ to find $y\in R$. Then
$r=a-by\ne0$ and
\[
       (a/b-y)^{-1}b^{-1}R=r^{-1}R,\qquad
                                  \phi(r)<\phi(b).
 \tag{A21.4}
\]
If $a/b\in R$, the remainder is zero, so the usual
division property holds in all cases.

Finally $\mathfrak p^h=(a)$ for some $a\in\mathcal O_K$,
because $[\mathfrak p]$ has order $h$. Its divisor
is supported only at $\mathfrak p$, and hence
$R=\mathcal O_K[1/a]$. Thus (A21.2) is the actual
affine open obtained by deleting this single closed
point, even when $\mathfrak p$ is not principal.
\end{proof}

Geometrically, NT-A20 concerns reduction in an
invertible sheaf generating
$\operatorname{Pic}(\operatorname{Spec}R)$.
NT-A21 removes one quantitatively bounded closed point
whose divisor class kills that Picard group, and then
uses the unit-Kummer RH family to supply Euclidean
reduction. Cyclicity remains a hypothesis; the
argument neither asserts that all class groups are
cyclic nor that the ideal norm is a Euclidean function.


\subsection*{NT-E14--E16. Finite-field detection of characteristic-zero isogeny}

The analytic input here is Dedekind RH for certain finite
number fields. It is not RH for the standard $L$-function
of a non-CM elliptic curve.
For an even squarefree integer $m\ge2$, let
$\mathcal D_m$ be the collection of finite Galois
extensions $T/\mathbf Q$ satisfying
\[
 [T:\mathbf Q]\le384,\qquad
 T/\mathbf Q\text{ is unramified at every prime }p\nmid m.
 \tag{E14.1}
\]
The hypothesis RH for $\zeta_{\mathbf Q}$ and for
$\zeta_T$, $T\in\mathcal D_m$, is sufficient for the
applications below. This is a finite family: the
discriminant bound (E16.4) below and Hermite's theorem
give finiteness. A precise reference is Milne,
\emph{Algebraic Number Theory}, v3.08 (2020),
Theorems~8.42--8.43, printed pp.~151--152.
Input $M$ supplies this family.
The more restricted fields actually used in the
argument are specified explicitly below.

\paragraph{NT-E14: the published non-CM isogeny test.}
Let $E_1,E_2/\mathbf Q$ be elliptic curves without
complex multiplication, with conductors $N_1,N_2$,
and put $m=\operatorname{rad}(2N_1N_2)$.
If the two curves are not $\mathbf Q$-isogenous,
then there is a prime $p\nmid2N_1N_2$ with
\[
 a_p(E_1)\ne a_p(E_2),\qquad
 p\le \bigl(482\log m+2880\bigr)^2.
 \tag{E14.2}
\]
Consequently equality of the point counts at all
such primes up to this bound is equivalent to
$\mathbf Q$-isogeny.
This is Mayle--Wang, \emph{On the effective version
of Serre's open image theorem},
Bull.\ London Math.\ Soc.~56 (2024), 1399--1416,
Theorem~2, printed p.~1401; complete proof
pp.~1408--1410. Its proof omits $2$ as well as all
bad primes, although the theorem statement only
requires common good reduction.

The precise source dependency is useful. Proposition~12,
pp.~1407--1408, constructs a finite Galois quotient field
\[
 F\subset\mathbf Q(E_1[2^\infty],E_2[2^\infty]),
                      \qquad [F:\mathbf Q]\le192.
 \tag{E14.3}
\]
It detects a trace difference. Corollary~6,
pp.~1404--1405, then applies effective Chebotarev to
$T=F(\sqrt m)$ to exclude every prime dividing $m$.
Thus only RH for this actual $\zeta_T$ and
$\zeta_{\mathbf Q}$ is used analytically.
The field $F$ is a quotient of the full $2$-adic image;
it is not asserted to be the first $2$-division field.
The convenient fixed family (E14.1) contains every
possible such $T$. The improved numerical constants
in (E14.2) use the published Bach--Sorenson Table~3
bounds through the source's computation
(13)--(16). Those numerical inputs are credited,
not represented as a fresh rerun of that table.

\paragraph{NT-E15: the published effective open-image bound.}
For a non-CM elliptic curve $E/\mathbf Q$ with
conductor $N_E$, input $M$ implies that
$\overline\rho_{E,\ell}$ is surjective for every
prime $\ell>B_E$, where
\[
                 B_E=964\log\operatorname{rad}(2N_E)+5760.
 \tag{E15.1}
\]
This is Mayle--Wang Theorem~1, printed p.~1401,
proved in \S5.2, pp.~1412--1413.
It retains both the base field $\mathbf Q$ and
the non-CM hypothesis. It does not prove the
uniform proposed cutoff $37$ for all curves.

Here a sufficient analytic family is (E14.1) for
$m=\operatorname{rad}(2N_E)$, together with
RH for $\zeta_{\mathbf Q}$. Indeed the proof
associates to a nonsurjective prime $\ell>37$ a
nontrivial quadratic twist $E_D$ with
$\operatorname{rad}(N_{E_D})\mid m$.
It applies (E14.2) to $E,E_D$; the two curves
are not $\mathbf Q$-isogenous by its Lemma~14.
The witnessing prime has a nonzero trace divisible
by $\ell$, and the Hasse bound gives
$\ell\le2\sqrt p$, producing (E15.1).
All analytic fields are therefore the fields
$F_D(\sqrt m)$ from (E14.3) for these twists,
of degree at most $384$ and unramified outside $m$.
The finitely many relevant squarefree $D$ divide
$N_E$, with their signs allowed.
The source's modular-curve exclusions, local
inertia lemma, Faltings theorem, and twist facts
are unconditional classical inputs; this reduction
does not claim to reprove them.

\paragraph{NT-E16: a complete decision formulation from integral equations.}
Let $E_1,E_2$ be given by nonsingular integral
Weierstrass equations over $\mathbf Q$, of
discriminants $\Delta_1,\Delta_2\ne0$.
No minimality, factorization, conductor, or
non-CM promise is required in this paragraph.
Put
\[
 R=2|\Delta_1\Delta_2|,\qquad
 b=\operatorname{bitlength}(R),\qquad
                      B=(1532b+10181)^2.
 \tag{E16.1}
\]
Under $M$, the curves are $\mathbf Q$-isogenous
if and only if
\[
 \#E_1(\mathbf F_p)=\#E_2(\mathbf F_p)
                  \quad\hbox{for every prime }p\le B,\ p\nmid R.
 \tag{E16.2}
\]
Testing (E16.2) by direct point counting is a
deterministic algorithm polynomial in the total
bit length of the integral coefficients.
This decides isogeny over $\mathbf Q$; it does
not construct an isogeny, decide geometric
isogeny over arbitrary extensions, or compute an
endomorphism ring.

More precisely, the sufficient RH family is
RH for $\zeta_{\mathbf Q}$ and
all $\zeta_T$ in (E14.1) with
$m=\operatorname{rad}(R)$.
The algorithm never needs to compute $m$.
For any nonisogenous pair only one actual
quotient field $F$ and the field $F(\sqrt m)$
are needed in the proof.

\begin{proof}
We give the finite-quotient argument explicitly;
it is the $\ell=2$, rank-two case of the
Faltings--Serre construction in Mayle--Wang
Proposition~12.
Faltings's isogeny theorem says that two abelian
varieties are isogenous over their number field
if and only if their local Euler factors agree
at all but finitely many places: Invent.\ Math.~73
(1983), Corollary~2, equivalence (i)--(iii),
printed p.~361. For elliptic curves over
$\mathbf Q$, the good local factor is determined
by $1-a_pT+pT^2$. If $E_1,E_2$ are not
$\mathbf Q$-isogenous, we can therefore choose
a common good odd prime where their traces differ.
A Frobenius element at that prime has differing
rational $2$-adic Tate traces. This argument also
applies to CM curves and needs no separate
character-criterion citation. Mayle--Wang recall
the elliptic criterion in (11), printed p.~1407.

Write $\rho=(\rho_1,\rho_2)$ for the product
of the integral rank-two Tate representations.
Let $A$ be the $\mathbf Z_2$-linear span of
$\rho(G_{\mathbf Q})$ in
$M_2(\mathbf Z_2)\times M_2(\mathbf Z_2)$.
Products of image elements remain image elements,
so $A$ is a ring containing the identity.
It is a finite free $\mathbf Z_2$-module of rank
at most eight. Let $e$ be the least finite
$2$-adic valuation of a nonzero difference
$\operatorname{tr}\rho_1(\sigma)-
 \operatorname{tr}\rho_2(\sigma)$.
The linear functional
\[
 \lambda:A\longrightarrow\mathbf Z_2,\qquad
 (u,v)\longmapsto
             2^{-e}(\operatorname{tr}u-\operatorname{tr}v)
 \tag{E16.3}
\]
is integral on $A$ and induces a nonzero functional
on at least one image element in $A/2A$.
Let $G$ be the image of $G_{\mathbf Q}$ in
$(A/2A)^\times$, and let $C$ be the nonempty
subset on which this functional is nonzero.
Trace invariance makes $C$ conjugacy-stable.
Every lift of an element of $C$ has different
two Tate traces.

Here $|G|\le2^8-1=255$: zero is not a unit.
The homomorphism to $G$ is continuous.
Indeed $A$ is a closed finite lattice in the
ambient matrix module, and $2A$ is open in $A$.
The fixed field $F$ of its kernel is consequently
a finite Galois extension of $\mathbf Q$,
contained in the joint $2$-adic division tower.
It lies in some finite level
$\mathbf Q(E_1[2^j],E_2[2^j])$, with $j\ge1$, so its degree
divides
$|\mathrm{GL}_2(\mathbf Z/2^j\mathbf Z)|^2
 =36\cdot2^{8j-8}$.
Among divisors of such an integer not exceeding
$255$, the largest is $192$.
Thus $[F:\mathbf Q]\le192$.
This argument uses $A/2A$ itself; it does not
identify it with the reduction of the ambient
matrix algebra, an identification that need
not be injective for an unsaturated lattice.

N\'eron--Ogg--Shafarevich gives that $F/\mathbf Q$
is unramified outside the prime divisors of $R$.
Here the exact cited form is Mayle--Wang
Theorem~8 and Corollary~9, printed pp.~1406--1407,
with their Remark~11 for a product of two curves;
Theorem~8 cites Serre--Tate, \emph{Good reduction
of abelian varieties}, Ann.\ Math.~88 (1968),
Theorem~1.
Set $m=\operatorname{rad}(R)$ and
$T=F(\sqrt m)$. This is Galois of degree
$d\le384$ and unramified outside $m$.
As $m$ is even and squarefree, every prime
dividing $m$ ramifies in $\mathbf Q(\sqrt m)$,
and hence in $T/\mathbf Q$.
The inverse image of $C$ in
$\operatorname{Gal}(T/\mathbf Q)$ is nonempty
and conjugacy-stable.
NT-C03 applied to any one of its conjugacy
classes provides, under RH for $\zeta_T$
and $\zeta_{\mathbf Q}$, an unramified prime
$p$ with differing Tate traces and
\[
               p\le(4\log D_T+2.5d+5)^2.
\]
In particular $p\nmid R$, so both input equations
have good reduction and the differing traces
are their actual $a_p(E_i)$.

The standard different bound, in the exact
form of Mayle--Wang Lemma~7, printed p.~1405,
gives
\[
 \begin{split}
 \log D_T
   &\le(d-1)\log\operatorname{rad}(D_T)+d\log d\\
   &\le383\log m+384\log384.
 \end{split}
 \tag{E16.4}
\]
For completeness the local calculation is
$v_{\mathfrak q}(\mathfrak D_T)
 \le e_{\mathfrak q}-1+
             e_{\mathfrak q}v_p(e_{\mathfrak q})$.
For a Galois field, $e_{\mathfrak q}\mid d$;
summing residue degrees over primes above $p$
bounds the exponent of $p$ in $D_T$ by
$d-1+d\,v_p(d)$.
Sum after multiplying by $\log p$ over
ramified $p$. The first terms contribute
$(d-1)\log\operatorname{rad}(D_T)$ and the
remaining terms are at most $d\log d$.
This proves (E16.4), retaining the classical
local different inequality as its input.

Since $\log384<6$, the Chebotarev expression is
bounded by
\[
 \begin{split}
 4\log D_T+2.5d+5
    &\le1532\log m+1536\log384+965\\
    &<1532\log R+10181
     \le1532b+10181.
 \end{split}
 \tag{E16.5}
\]
The last inequality uses $R<2^b$ and $\log2<1$.
This proves the existence of a differing point
count within (E16.1) whenever the two curves
are not $\mathbf Q$-isogenous.
Conversely a $\mathbf Q$-isogeny identifies
the rational Tate representations, so the
point counts agree at every common good
prime. This proves the equivalence.

Finally let $L$ be the total input bit length.
Each discriminant is a fixed polynomial in
the five integral Weierstrass coefficients;
it is computable in polynomial time and
has bit length $O(L)$. Thus $b=O(L)$ and
$B=O(L^2)$. Trial division or a sieve
enumerates all primes up to the integer $B$
in polynomial time. Testing $p\nmid R$ only
needs a remainder calculation.
At each retained prime enumerate all
$(x,y)\in\mathbf F_p^2$ satisfying the reduced
Weierstrass equation and add the unique point
at infinity. This uses $O(p^2)$ elementary
modular operations; even the sum over all
integers $p\le B$ is $O(B^3)$.
The operands have polynomial bit length.
Consequently the entire finite test is
polynomial in $L$, with no factoring of $R$,
conductor calculation, or point-counting
oracle. A negative output includes a prime
and two different counts as an unconditional
certificate; correctness of every positive
output uses the stated RH input.
\end{proof}

NT-E14--E15 are precise published theorem
transfers with their analytic field families
identified from the full source proofs.
NT-E16 is a fully proved formulation of the
finite-quotient method with a weaker elementary
constant and a direct input model; no
literature-priority claim is made.
Geometrically, this detects isomorphism of
the semisimple rank-two Tate local systems
using closed points of the good-reduction
open of $\operatorname{Spec}\mathbf Z$,
then invokes Faltings to recover
$\mathbf Q$-isogeny. The finite quotients
have nonabelian images in general, but their
Dedekind functions are trivial-character
instances of $M$. No non-CM $\mathrm{GL}_2$
GRH or general Artin holomorphy has entered.
Work/run: \texttt{effective\_elliptic\_detection.tex},
\texttt{GRH\_NIGHT/number\_theory}.

\subsection*{NT-E17--E18. A quadratic-Dirichlet input for CM recognition}

Write $A_{\rm quad}$ for RH for $\zeta_{\mathbf Q}$ and
for every primitive quadratic Dirichlet $L$-function.
This is a subfamily of the user-designated input $A_D$.
No Hecke function over a number field of degree greater
than one needs to be assumed separately below.
For a nonzero integer polynomial $H$, put
\[
 h=\deg H,\qquad
 \mathcal H=\max\{1,\log\max_i|H_i|\}.
\]
The height is padded by one to include polynomials
whose nonzero coefficients are units.

\paragraph{Exact classical inputs and scope.}
We use Cremona--Sutherland,
\emph{Computing the endomorphism ring of an elliptic
curve over a number field}, Contemporary Mathematics~796
(2024), 75--102, in author version
arXiv:2301.11169v4 (15 December 2023).
All page numbers in this subsection refer to that
author version.

Its Proposition~19, pp.~9--10, proves under GRH that
every imaginary quadratic order discriminant $D<0$
of class number $h$ satisfies
\[
 |D|\le
 \left(\frac{12e^{2\gamma}}{\pi}\right)^2
 h^2\bigl(\log\log(h+1)+4\bigr)^4.
 \tag{E17.1}
\]
The analytic input in this proof is
Languasco--Trudgian, \emph{Uniform effective estimates
for $|L(1,\chi)|$}, Corollary~1, author p.~2,
applied to the primitive quadratic character of the
fundamental discriminant.
The remaining reduction from maximal to arbitrary
orders uses the class-number formula (2.1) and an
unconditional Euler-totient estimate.
Thus $A_{\rm quad}$ supplies the analytic hypothesis
of (E17.1). The cited small-discriminant
tabulations and numerical range checks in that
published proposition are retained as classical
inputs; they were not rerun here.

Let $\Phi_\ell(X,Y)\in\mathbf Z[X,Y]$ be the classical
modular polynomial for cyclic $\ell$-isogenies.
For monic irreducible $H\in\mathbf Z[X]$ and a root $j$,
define
\[
 m_\ell(j)=
 \operatorname{ord}_H
    \operatorname{Res}_Y\bigl(H(Y),\Phi_\ell(X,Y)\bigr).
 \tag{E17.2}
\]
Cremona--Sutherland equation (2.2), p.~8, gives
this equality with the number of conjugates joined
to $j$ by a cyclic $\ell$-isogeny, with multiplicity.
Their Corollary~14, p.~8, gives $m_\ell(j)=2$
for a CM $j$ of order discriminant $D$ whenever
$\ell\nmid D$ splits in $\mathbf Q(\sqrt D)$.
Their Proposition~16, p.~9, gives
$m_\ell(j)\le1$ for every prime $\ell$ if $j$ is not CM.
These are classical algebraic statements, independent
of GRH. Their supplied derivations are read here;
the upstream theory of isogeny volcanoes and
$K$-curves remains cited.

The published Algorithm~1 has conditional correctness
in Theorem~24, p.~14, and deterministic complexity
in Theorem~27 and its full proof, p.~14.
For that particular algorithm the published time
and space bounds are respectively
$(h^2\mathcal H)^{1+o(1)}$ and
$(h\mathcal H)^{1+o(1)}$.
For our modified test we assert the weaker
polynomial bound, with its exact-arithmetic
verification supplied below.
We use the deterministic polynomial-time computation
of $\Phi_\ell$ explicitly cited in the proof
of Theorem~27.
A direct primary locator is
Blake--Csirik--Rubinstein--Seroussi,
\emph{On the Computation of Modular Polynomials
for Elliptic Curves}, \S2, pp.~4--7:
the exact power-series and CRT algorithm is
given there, with the bit-complexity discussion
and coefficient bound (9) on p.~7.
The modified test below examines all the bounded
degrees, so it does not require the source's
least-ordinary-prime selection or its finite list
of abelian discriminants.

\paragraph{NT-E17: a bounded algebraic CM certificate under $A_D$.}
Let $H\in\mathbf Z[X]$ be monic and irreducible,
of degree $h\ge1$, and set
\[
 b_h=\operatorname{bitlength}(h+1),\qquad
                     P_h=(72b_h+264)^2.
 \tag{E17.3}
\]
Under $A_{\rm quad}$, $H$ is a Hilbert class
polynomial for an imaginary quadratic order if
and only if there is a prime $5\le\ell\le P_h$ with
\[
 H(X)^2\ \bigm|\ 
       \operatorname{Res}_Y\bigl(H(Y),\Phi_\ell(X,Y)\bigr)
                           \quad\hbox{in }\mathbf Z[X].
 \tag{E17.4}
\]
Scanning these primes and performing the exact
divisibility tests is deterministic and takes
polynomial time in $h\mathcal H$.
A positive answer, accompanied by $\ell$ and
the exact divisibility identity, is valid
unconditionally. The universal correctness of
a negative answer uses $A_{\rm quad}$.
Irreducibility is part of the stated input promise.

\begin{proof}
If (E17.4) holds, (E17.2) gives $m_\ell(j)\ge2$.
The cited non-CM bound therefore forces $j$ to be CM.
Its monic irreducible polynomial is then exactly
the Hilbert class polynomial of its endomorphism
order. This implication is unconditional.

For the converse suppose $H=H_D$.
First we replace (E17.1) by an elementary bound
whose constants can be used without computing
$\gamma$, $\pi$, or logarithms in the algorithm.
For $h\ge1$,
\[
 \log\log(h+1)+4
 \le\log(h+1)+4
 \le h+4\le3(h+1).
\]
Using $\gamma<1$, $e<3$, and $\pi>3$ in (E17.1),
we obtain
\[
 |D|<1296\cdot81\,h^2(h+1)^4
       <2^{17}(h+1)^6.
 \tag{E17.5}
\]
This bound includes $D=-3,-4$ and nonmaximal orders
because (E17.1) does.

Set $m=\operatorname{rad}(6|D|)$ and
\[
 K=\mathbf Q(\sqrt D),\qquad
 T=K(\sqrt m).
\]
The integer $m$ is positive, even, and squarefree.
The field $\mathbf Q(\sqrt m)$ is real quadratic
and $K$ is imaginary quadratic, so $T/\mathbf Q$
is biquadratic of degree four.
It is unramified outside the primes dividing $m$.
Every prime dividing $m$ ramifies already in
$\mathbf Q(\sqrt m)$, since its discriminant is $4m$.
Consequently every unramified prime of $T$ is
prime to $6D$.

The full Euler factorization of $\zeta_T$ is
$\zeta_{\mathbf Q}$ times the three quadratic
Dirichlet functions of its quadratic subfields.
Thus $A_{\rm quad}$ supplies RH for $\zeta_T$
and $\zeta_{\mathbf Q}$.
Apply NT-C03 to the identity class of
$\operatorname{Gal}(T/\mathbf Q)$.
There is a prime $\ell$ splitting completely in $T$,
with $\ell\nmid6D$ and
\[
                 \ell\le(4\log|\operatorname{disc}T|+15)^2.
\]
The Galois different bound proved in NT-E16,
with absolute degree four, gives
\[
 \log|\operatorname{disc}T|
       \le3\log m+4\log4
       \le3\log(6|D|)+4\log4.
 \tag{E17.6}
\]
It follows that
\[
 \begin{aligned}
 4\log|\operatorname{disc}T|+15
   &\le12\log|D|+12\log6+16\log4+15\\
   &<12\log|D|+60\\
   &<72\log(h+1)+204\log2+60\\
   &<72b_h+264.
 \end{aligned}
 \tag{E17.7}
\]
For the second line one can use
$\log6<9/5$ and $\log4<7/5$, which give the
stronger numerical sum $59$.
Both inequalities follow by summing the exponential
series through degree six:
the resulting positive margins over $6$ and $4$
are respectively $42629/1250000$ and
$592609/11250000$.
The last line uses $h+1<2^{b_h}$ and $\log2<1$.
We have proved $5\le\ell\le P_h$, $\ell\nmid D$,
and splitting in $K$.
The cited CM formula gives $m_\ell(j)=2$,
and hence (E17.4).

Finally $P_h=O(\log^2(h+2))$, so both the number
and the sizes of the tested primes are polynomial
in the input length.
The cited deterministic construction of
$\Phi_\ell$ takes polynomial time in $\ell$,
and its output has polynomial bit size.
Here is an elementary polynomial bound for
the rest of the exact computation.
The Sylvester matrix defining the resultant
has size $N=h+\ell+1$, and its entries are
integer polynomials in $X$ of degree at most
$\ell+1$ with polynomially bounded coefficient
bit lengths.
Its determinant has degree at most
$d=N(\ell+1)$.
Evaluate at $X=0,\ldots,d$.
Each resulting integer matrix has entries of
polynomial bit length, and its determinant can
be computed by exact rational elimination
in polynomial time.
The determinant-size bound
$N!(\max|a_{ij}|)^N$ bounds all these outputs.
Rational interpolation at the $d+1$ distinct
integers recovers the resultant in polynomial
time; the factorial denominators of the
Lagrange basis have polynomial bit length.
Finally division by the monic polynomial $H^2$
uses polynomially many exact integer operations.
At each division step the coefficient bit bound
increases by at most the coefficient bound for
$H^2$ plus a logarithm of its degree, so it
remains polynomial throughout.
The remainder vanishes exactly when (E17.4)
holds. This proves the stated complexity
and both directions of the test.
\end{proof}

The only analytic characters in this proof are the
quadratic character used in (E17.1), the three
quadratic characters of the displayed biquadratic
field, and the trivial character.
In particular $M$, ring-class Hecke RH, general
Artin holomorphy, and elliptic $\mathrm{GL}_2$ RH
are not additional inputs.
This is a new formulation and input reduction,
with no claim of literature priority.

\paragraph{NT-E18: translation to an elliptic curve over a varying number field.}
Suppose a number field
$F=\mathbf Q[t]/(f(t))$ is given by a monic
irreducible $f\in\mathbf Q[t]$, and a nonsingular
Weierstrass equation for $E/F$ is given with
its coefficients represented in the power basis.
The defining polynomial is supplied as a dense
coefficient list, and field elements as dense
power-basis vectors. All rational coefficients
are encoded in binary, so $n=\deg f$ is at most
the total input length;
the assertion that $f$ defines a field is the
input promise.
Under $A_{\rm quad}$, geometric CM of $E$
can be decided deterministically in polynomial
time in this total input length.
No maximal order, conductor, integer factorization,
or endomorphism basis is required.

\begin{proof}
Compute $\Delta\ne0$ and $j=c_4^3/\Delta$ using
exact arithmetic in $F$.
Multiplication is reduction modulo $f$;
inversion is rational linear algebra.
These operations have polynomial bit complexity.
Let $n=\deg f$. Compute the coordinate vectors
of $1,j,\ldots,j^n$ and find the first dependence.
It gives the monic minimal polynomial
$H_j\in\mathbf Q[X]$, of degree $h\le n$.

For clarity the coefficient-size assertion does
not assume a precomputed integral basis.
The multiplication matrix of $j$ has rational
entries of polynomial bit length in the input.
After a common denominator is cleared, its first
$n$ powers have polynomial bit length: multiplying
integer matrices increases the logarithmic
entry bound by at most the previous bound plus
the fixed matrix-entry bound and $\log n$.
Cramer's determinant bounds for the first
linear dependence therefore give coefficients
of $H_j$ of polynomial bit length.
Exact elimination consequently takes polynomial
time as well.

If $H_j$ has a nonintegral coefficient, then $j$
is not an algebraic integer and cannot be a
CM $j$-invariant; return non-CM.
Otherwise $H_j$ is a monic irreducible integer
polynomial. Apply NT-E17 to it.
The degree and padded logarithmic height of
$H_j$ are polynomially bounded in the original
input length, so the complete computation is
polynomial. The characterization of CM
$j$-invariants as the roots of the Hilbert class
polynomials proves correctness.
\end{proof}

This recognizes whether
$\operatorname{End}_{\overline F}(E)$ is $\mathbf Z$
or an imaginary quadratic order.
It does not output the order discriminant or
explicit endomorphisms.
Cremona--Sutherland Algorithm~2 does output
the discriminant, and its correctness and
termination are Theorem~30, p.~15.
Its near-quasilinear expected running time is
Theorem~35, p.~17, whose hypotheses include
the separate prime-distribution heuristic (H)
and the computational heuristics in its
reference [45, \S7.1], as well as GRH.
Neither $A_D$ nor $M$ by itself removes those
extra hypotheses.
Likewise Enge's floating-point class-polynomial
complexity, Corollary~1.3 of arXiv:cs/0601104v2,
pp.~2--3, explicitly retains a rounding-error
heuristic; it is not used in either proof above.
These distinctions concern the cited algorithms,
not an assertion that unconditional CM recognition
is impossible.

Geometrically (E17.4) detects the quadratic
endomorphism algebra of $H^1(E)$ through the
cyclic-isogeny correspondences on the moduli
line of elliptic curves. The conditional
existence of a bounded correspondence degree
comes from a finite biquadratic Frobenius cover
of an arithmetic open, in the finite-etale
dictionary of SGA~1, Expose~V, \S7.
The result distinguishes CM torus monodromy
from the non-CM case without proving a
rank-two automorphic GRH.
Work/run: \texttt{cm\_recognition\_dirichlet.tex},
\texttt{GRH\_NIGHT/number\_theory}.
